% twn_fuse Linux driver 開發講義（v1.1，繁體中文）。建置：在本資料夾執行 latexmk，輸出 build/handout.pdf。
\documentclass[11pt,a4paper]{article}
\usepackage[top=20mm,bottom=20mm,left=20mm,right=20mm,headheight=14pt]{geometry}
\usepackage{fontspec}
\usepackage{xeCJK}
\setmainfont{Segoe UI}
\setsansfont{Segoe UI}
\setmonofont{Consolas}[Scale=0.92]
\setCJKmainfont{Microsoft JhengHei}[BoldFont={Microsoft JhengHei Bold}]
\setCJKsansfont{Microsoft JhengHei}[BoldFont={Microsoft JhengHei Bold}]
\setCJKmonofont{Microsoft JhengHei}
\xeCJKsetup{CJKecglue={\hskip 0.15em plus 0.05em}}
\linespread{1.3}
\setlength{\parindent}{0pt}
\setlength{\parskip}{0.55em}

\usepackage{xcolor}
\usepackage{graphicx}
\usepackage{booktabs}
\usepackage{tabularx}
\usepackage{xltabular}
\usepackage{array}
\usepackage{enumitem}
\setlist{itemsep=0.15em,topsep=0.25em,leftmargin=1.4em}
\usepackage[most]{tcolorbox}
\usepackage{listings}
\usepackage{tikz}
\usetikzlibrary{positioning,arrows.meta,fit,calc,backgrounds}
\usepackage{fancyhdr}
\usepackage[hidelinks]{hyperref}
\usepackage{titlesec}
\usepackage{needspace}
\usepackage{etoolbox}

\definecolor{ink}{HTML}{1F3A5F}
\definecolor{oldc}{HTML}{8A5A00}   \definecolor{oldbg}{HTML}{FFF6E5}
\definecolor{techc}{HTML}{1F5FA8}  \definecolor{techbg}{HTML}{EAF2FC}
\definecolor{gainc}{HTML}{2E7D32}  \definecolor{gainbg}{HTML}{EAF6EA}
\definecolor{qac}{HTML}{5B3F8C}    \definecolor{qabg}{HTML}{F3EEFA}
\definecolor{pitc}{HTML}{A23B2A}   \definecolor{pitbg}{HTML}{FCEDEA}
\definecolor{codebg}{HTML}{F7F7F5}
\definecolor{kern}{HTML}{DCE8F5}   \definecolor{user}{HTML}{F2EBDD}  \definecolor{hw}{HTML}{E3F1E3}

\titleformat{\section}{\Large\bfseries\color{ink}}{\thesection}{0.8em}{}
\titleformat{\subsection}{\large\bfseries\color{ink}}{\thesubsection}{0.6em}{}
\titlespacing*{\section}{0pt}{0.6em}{0.6em}
\titlespacing*{\subsection}{0pt}{0.9em}{0.35em}
% 小節標題下面至少要放得下 8 行，標題才不會單獨留在頁底
\pretocmd{\subsection}{\needspace{8\baselineskip}}{}{}
% 版面是 raggedbottom：頁底留白不算代價，所以任何負的 penalty 都會被當成最佳斷點，整頁後面的
% 內容只要放不下一個不可拆的框，就退回到那裡換頁、留下大片空白。標準的 \needspace 會先放一個
% \penalty-100，章節標題本身也有 \@secpenalty（-300）；兩者都拿掉，只在真的放不下時才強制換頁。
% 量剩餘空間之前先放一個 \penalty0：page builder 只在遇到斷點時才換頁，前一個不可拆的框（紀錄框）
% 若已經超出這一頁，\pagetotal 還是舊頁加上它的高度；先讓它把舊頁送出，量到的才是新的一頁。
\makeatletter
\@secpenalty=0
\renewcommand{\needspace}[1]{\par\penalty\z@\begingroup\setlength{\dimen@}{#1}%
  \dimen@ii\pagegoal\advance\dimen@ii-\pagetotal
  \ifdim\dimen@>\dimen@ii\ifdim\dimen@ii>\z@\vfil\fi\penalty-\@M\fi\endgroup}
\makeatother

\pagestyle{fancy}
\fancyhf{}
\fancyhead[L]{\small\color{gray}twn\_fuse Linux driver 開發講義}
\fancyhead[R]{\small\color{gray}\leftmark}
\fancyfoot[C]{\small\thepage}
\renewcommand{\headrulewidth}{0.3pt}
\renewcommand{\sectionmark}[1]{\markboth{\thesection\ #1}{}}

% ---- 每一章固定的框 ----
\newtcolorbox{oneline}{colback=ink!6,colframe=ink,boxrule=0.6pt,arc=1mm,left=2mm,right=2mm,
  fontupper=\bfseries,before skip=0.4em,after skip=0.6em}
% 「原本」可能是 v1.0.0 的 UIO 寫法，也可能是 driver v1.0；框的標題寫清楚是哪一個
\newtcolorbox{before}[1][原本怎麼做（v1.0.0 的 UIO 版本）]{colback=oldbg,colframe=oldc,title={#1},
  fonttitle=\bfseries,boxrule=0.6pt,arc=1mm,breakable}
\newtcolorbox{tech}{colback=techbg,colframe=techc,title={用了什麼技術},fonttitle=\bfseries,
  boxrule=0.6pt,arc=1mm,breakable}
\newtcolorbox{gain}{colback=gainbg,colframe=gainc,title={改善了什麼},fonttitle=\bfseries,
  boxrule=0.6pt,arc=1mm,breakable}
\newtcolorbox{qa}[1]{colback=qabg,colframe=qac,title={面試會問（完整版，約 1 分鐘）：#1},
  fonttitle=\bfseries,boxrule=0.6pt,arc=1mm,breakable}
\newtcolorbox{qashort}[1]{colback=qabg,colframe=qac,title={#1},fonttitle=\bfseries,
  boxrule=0.6pt,arc=1mm,breakable}
\newtcolorbox{pit}{colback=pitbg,colframe=pitc,title={踩過的坑},fonttitle=\bfseries,
  boxrule=0.6pt,arc=1mm,breakable}

% ---- 程式碼：一律由 extract.py 從 repo 的實際檔案抽出，保留原始行號 ----
\lstset{language=C,basicstyle=\ttfamily\footnotesize,numbers=left,numberstyle=\tiny\color{gray},
  numbersep=6pt,xleftmargin=1.6em,
  backgroundcolor=\color{codebg},frame=single,rulecolor=\color{gray!40},breaklines=true,
  breakatwhitespace=true,columns=fullflexible,keepspaces=true,showstringspaces=false,tabsize=4,
  keywordstyle=\color{techc}\bfseries,commentstyle=\color{gray!90!black}\itshape,
  stringstyle=\color{oldc},captionpos=t,aboveskip=0.5em,belowskip=0.5em,
  morekeywords={u32,u64,size_t,bool,dma_addr_t,__u32,__u64,uint32_t,uint64_t}}
% \snip[語言]{標籤}：印出 build/snippets/<標籤>.c，標題是「檔案 : 起始行–結束行」（舊版本的程式
% 另外標出 commit）。開始一個程式框之前，估計這一頁還放得下幾行程式（一行約 0.75 個內文行距，
% 從建出的 PDF 量的，略為保守；標題與框另算 3 行；行數用 extract.py 算好的 \snipvis，太長而折行的
% 算兩行）。框放不下而會被拆開時，若留在這一頁的少於 6 行就先換頁；若跑到下一頁的少於 5 行，
% 短的框先換頁，長的框（換頁會空下 25 行以上）則自己在倒數第 5 行之前拆開，下一頁至少有 5 行。
% 其他情形照常拆開。在框裡面（\ifinner）不檢查：框自己決定怎麼分頁。
\newcount\sniplines
\newcount\snipfit
\newcount\snipcut   % > 0：這一頁只印前 \snipcut 行，其餘換頁再印
\newcommand{\snip}[2][C]{%
  \input{build/snippets/#2.tex}%
  \sniplines=\snipvis\relax \snipcut=0
  \par\ifinner\else
    \penalty0 % 同 \needspace：先讓已經超出的舊頁送出，再量
    \dimen0=\pagegoal \advance\dimen0 by -\pagetotal
    \ifdim\dimen0>3\baselineskip
      \snipfit=\numexpr\dimexpr\dimen0-3\baselineskip\relax/\dimexpr\baselineskip*75/100\relax\relax
      \ifnum\snipfit<\sniplines
        \ifnum\numexpr\sniplines-\snipfit\relax<5
          \ifnum\snipfit<25 \vfil\penalty-10000
          \else \snipcut=\numexpr\sniplast-\snipfirst+1-5\relax \fi
        \else\ifnum\snipfit<6 \vfil\penalty-10000\fi\fi
      \fi
    \else\ifdim\dimen0>0pt \vfil\penalty-10000\fi\fi
  \fi
  \ifnum\snipcut>0
    \lstinputlisting[language=#1,firstnumber=\snipfirst,lastline=\the\snipcut,
      title={\scriptsize\ttfamily\snippath\ :\ \snipfirst--\sniplast}]{build/snippets/#2.c}%
    \par\vfil\penalty-10000
    \lstinputlisting[language=#1,firstline=\the\numexpr\snipcut+1\relax,
      firstnumber=\the\numexpr\snipfirst+\snipcut\relax,aboveskip=0pt]{build/snippets/#2.c}%
  \else
    \lstinputlisting[language=#1,firstnumber=\snipfirst,
      title={\scriptsize\ttfamily\snippath\ :\ \snipfirst--\sniplast}]{build/snippets/#2.c}%
  \fi}
% 段落的最後一兩行不單獨留到下一頁（紀錄框前面的引言常常只剩一行跟著框走）
\widowpenalties 3 10000 10000 150

% ---- 板上與模擬的實際輸出 ----
% \boardfile[說明]{檔名}{起-迄行}：直接從 results/raw/ 讀出紀錄原文，不手抄。檔名的第一個字母就是
% 導讀表格裡的開機組：A、B、C 與 dtb_diff 在 linux_driver_20261026/（driver v1.0 那一輪），其餘
% （D–I、S、first_bitstream/）在 linux_driver_v11_20261027/。標題自動帶上開機組。框左邊的灰色數字
% 是紀錄檔本身的行號，正文用「第 N 行」引用的就是它。listings 的 linerange 遇到跳行時會接著
% 往下編號，所以行號範圍裡的每一段（以逗號分開）各印成一個緊接著的框，從那一段的起始行號編起。
\lstdefinestyle{board}{language={},numbers=left,numberstyle=\tiny\color{gray},numbersep=4pt,
  xleftmargin=0pt,belowskip=0pt,
  basicstyle=\ttfamily\scriptsize\color{white},backgroundcolor=\color{black!82},
  rulecolor=\color{black!82}}
\ExplSyntaxOn
\seq_new:N \l__board_seq
\tl_new:N \l__board_first_tl
\tl_new:N \l__board_last_tl
\bool_new:N \l__board_first_bool
\cs_new_protected:Npn \__board_titled:nnn #1#2#3
  { \lstinputlisting[style=board,firstline=#1,lastline=#2,firstnumber=#1,title={#3}]{\boardpath} }
\cs_new_protected:Npn \__board_cont:nn #1#2
  { \lstinputlisting[style=board,firstline=#1,lastline=#2,firstnumber=#1,aboveskip=0pt]{\boardpath} }
\cs_generate_variant:Nn \__board_titled:nnn { VVn }
\cs_generate_variant:Nn \__board_cont:nn { VV }
% \boardranges{起-迄,起-迄,...}{標題}：每一段一個框，只有第一個帶標題
\NewDocumentCommand \boardranges { m m }
  {
    \bool_set_true:N \l__board_first_bool
    \clist_map_inline:nn {#1}
      {
        \seq_set_split:Nnn \l__board_seq { - } {##1}
        \tl_set:Ne \l__board_first_tl { \seq_item:Nn \l__board_seq { 1 } }
        \tl_set:Ne \l__board_last_tl { \seq_item:Nn \l__board_seq { -1 } }
        \bool_if:NTF \l__board_first_bool
          {
            \__board_titled:VVn \l__board_first_tl \l__board_last_tl {#2}
            \bool_set_false:N \l__board_first_bool
          }
          { \__board_cont:VV \l__board_first_tl \l__board_last_tl }
      }
  }
\ExplSyntaxOff
\newcommand{\rawdir}{../../../results/raw}
\newcommand{\logsub}[1]{\logsubi#1\relax}
\def\logsubi#1#2\relax{\if#1Alinux_driver_20261026\else\if#1Blinux_driver_20261026\else
  \if#1Clinux_driver_20261026\else\if#1dlinux_driver_20261026\else
  linux_driver_v11_20261027\fi\fi\fi\fi}
\newcommand{\whereof}[1]{\whereofi#1\relax}
\def\whereofi#1#2\relax{\if#1S（PC）\else\if#1d（PC）\else\if#1f（第一版 v1.1 bitstream）\else
  （開機組 #1）\fi\fi\fi}
% 每個框都不超過 25 行，放進 minipage 讓它不會跨頁斷開；框前面不准換頁，放不下時前一段的最後
% 幾行會跟著框一起移到下一頁，引言才不會和框分開。紀錄檔不存在時印出紅字，交付前不可以有
\newcommand{\boardfile}[3][]{%
  \edef\boardsub{\logsub{#2}}%
  \edef\boardpath{\rawdir/\boardsub/#2}%
  \IfFileExists{\boardpath}{%
  \par\nopagebreak\noindent\begin{minipage}{\linewidth}%
  \boardranges{#3}{\scriptsize\ttfamily results/raw/\expandafter\detokenize\expandafter{\boardsub}/\detokenize{#2}\ :\ #3\normalfont\whereof{#2}#1}%
  \end{minipage}\par}{\par\noindent\textcolor{red}{（紀錄待補：\detokenize{#2} #3）}\par}}
\lstnewenvironment{cmdline}{\lstset{language=bash,numbers=none,xleftmargin=0pt,
  basicstyle=\ttfamily\footnotesize}}{}

% \code{...}：行內程式碼，保留空白；在 _ / . 之後允許換行，長名稱才不會超出版面
\ExplSyntaxOn
\NewDocumentCommand \code { m }
  {
    \tl_set:Ne \l_tmpa_tl { \tl_to_str:n {#1} }
    \regex_replace_all:nnN { ([_/\.]) } { \1 \c{allowbreak} } \l_tmpa_tl
    \texttt { \tl_use:N \l_tmpa_tl }
  }
\ExplSyntaxOff
\emergencystretch=2em
\newcommand{\todo}[1]{\textcolor{red}{（待寫：#1）}}
\newcommand{\sref}[1]{\S\ref{#1}}
% \lineof{標籤}：該段程式碼在原始檔的起始行號（給速查表用，跟著程式碼自動更新）
\newcommand{\lineof}[1]{\begingroup\input{build/snippets/#1.tex}\snipfirst\endgroup}

\title{\bfseries twn\_fuse Linux driver 開發講義\\[0.3em]
  \large 從 UIO 到 kernel driver：probe、DMA、ioctl、中斷、記憶體隔離與發佈（v1.1）}
\author{Wei-Jhe Huang}
\date{2026 年 10 月 27 日}
\renewcommand{\contentsname}{目錄}
\renewcommand{\figurename}{圖}
% 講義裡會重複出現的量測數字只寫在這裡一次，來源都是 results/raw/linux_driver_20261026/。
% 改數字只改這個檔案。
\newcommand{\errChecks}{14}          % twn_errtest timeout 的檢查項數（A2/B2 第 2 步，14 行 PASS）
\newcommand{\errRefusals}{12}        % 其中該被拒絕的項數（另 2 項該成功）
\newcommand{\timeoutSec}{3.01}       % A2/B2 第 25 行
% 吞吐量：A2_board_check_uio_dtb.log 第 40-45 行（原本的 dtb，這是 v1.0.0 發佈的設定）
\newcommand{\fpsUioRuns}{230.13、230.22、229.78}
\newcommand{\fpsDrvRuns}{232.05、232.13、231.44}
\newcommand{\fpsUioMean}{230.04}
\newcommand{\fpsDrvMean}{231.87}
% 同一套量測在改過 compatible 的 dtb 下：B2_board_check_patched_dtb.log 第 40-45 行
\newcommand{\fpsUioRunsB}{229.59、229.23、229.13}
\newcommand{\fpsDrvRunsB}{230.83、232.17、231.74}
\newcommand{\fpsUioMeanB}{229.32}
\newcommand{\fpsDrvMeanB}{231.58}
% 一幀之後的 sysfs：A3_driver_state.txt 第 20 行
\newcommand{\cyclesTotal}{990,917}


\begin{document}
\maketitle
\thispagestyle{empty}
% 每一章的小節都是同樣的八個部分，目錄只列章
\setcounter{tocdepth}{1}
{\setlength{\parskip}{0pt}\tableofcontents}
\clearpage

\section*{導讀}
\addcontentsline{toc}{section}{導讀}

\subsection*{這份講義怎麼讀}
\todo{每一節的固定結構：一句話 → 原本怎麼做 → 用了什麼技術 → 改善了什麼 → 實際例子 → 原理 → 面試會問 → 關鍵函式速查}

\subsection*{整體架構：UIO 版本與 driver 版本}
\todo{一張圖：app、kernel、CSR、DDR arena、PL 在兩個版本裡的關係}

\subsection*{範例：原本的 frame 執行函式}
下面是 v1.0.0 跑一個 backbone frame 的函式，由 \texttt{extract.py} 從 repo 直接抽出（建置流程測試用，開工後會換成正式內容）。
\snip{pl_hal_fuse_c}{fuse_run_frame}

\clearpage
\section{Device tree、probe 與 binding：kernel 怎麼找到 driver}
\label{sec:probe}

\begin{oneline}
kernel 拿 device tree 節點的 \code{compatible} 字串去比對每個 driver 的比對表，對上了就呼叫那個 driver 的 \code{probe()} 接管硬體；binding 則把「這個節點該長什麼樣」寫成可以檢查的規格。
\end{oneline}

\begin{before}
PetaLinux 從 Vivado 匯出的硬體描述自動產生節點。v1.0.0 在 \code{system-user.dtsi} 把它的 \code{compatible} 改寫成 \code{"generic-uio"}：

\snip[{}]{v100_dt_uio_node}

開機參數 \code{uio_pdrv_genirq.of_id=generic-uio} 再讓 UIO driver 認這個字串。UIO driver 不知道這是什麼硬體，只把節點的位址範圍開成 \code{/dev/uio4}；app 自己去 \code{/sys/class/uio/} 找名字對的編號，再 \code{mmap} 進來：

\snip{v100_uio_hal_init}

問題：\code{"generic-uio"} 描述的是「要交給哪個 driver」，不是「這是什麼硬體」，而 device tree 只該描述硬體；要靠開機參數配合；kernel 對這顆硬體一無所知，不檢查、不配置資源，也不知道誰在用它。節點該有哪些屬性，也沒有寫在任何地方。
\end{before}

\begin{tech}
\begin{itemize}
\item \code{struct platform_driver}：SoC（把 CPU 與周邊做在同一顆晶片上）周邊的 driver 型別。這類周邊不在 PCI、USB 這種能自動列舉的匯流排上，要靠 device tree 告訴 kernel 它們在哪裡。
\item \code{struct of_device_id} 比對表（\code{of_match_table}）與 \code{MODULE_DEVICE_TABLE(of, ...)}：列出認得的 \code{compatible} 字串，並寫進 module 的資訊，讓系統知道哪個 module 該被自動載入（第 14 章）。
\item \code{probe()}／\code{remove()} 與 \code{devm_*}（device-managed）系列：資源跟著裝置走，driver 離開或 probe 失敗時自動釋放。
\item \textbf{binding}：用 dt-schema 格式（JSON Schema 寫成的 YAML）描述節點，\code{wjhuang,twn-fuse.yaml}。
\end{itemize}
\end{tech}

\begin{gain}
\begin{itemize}
\item 硬體由自己的名字 \code{"wjhuang,twn-fuse-1.1"} 認領；字串尾端的版本號讓有中斷的新硬體和沒有中斷的舊硬體各自對到正確的處理方式。
\item \code{probe()} 時就檢查硬體、配置記憶體、申請中斷、建立 \code{/dev/twn_fuse}，結果寫進 \code{dmesg}，不必等 app 跑起來才發現問題。
\item 節點的規格寫成檔案，並拿 PetaLinux 實際產生的 \code{system.dtb} 來比對；比對時發現 PetaLinux 多產生了 \code{interrupt-names}，事先沒有人知道。
\end{itemize}
\end{gain}

\subsection{實際例子}

\textbf{比對表和 driver 宣告。}比對表有兩筆：\code{-1.1} 是有中斷的硬體（第 6 章），\code{-1.0} 是 v1.0.0 那顆沒有中斷的硬體，最後一筆空的表示結束。\code{struct platform_driver} 把 probe、remove、比對表與 sysfs 屬性（第 8 章）綁在一起；\code{suppress_bind_attrs} 的理由在第 7 章：

\snip{drv_of_match}
\snip{drv_driver}

最後一行 \code{module_platform_driver(twn_driver);} 展開後就是 module 的 \code{init} 與 \code{exit}：載入時 \code{platform_driver_register()}，卸載時 \code{platform_driver_unregister()}。

\textbf{probe：接管硬體的五個步驟。}\code{probe()} 收到一個 \code{struct platform_device}，裡面已經有由 \code{reg} 轉成的位址範圍，以及指向 device tree 節點的指標（中斷要等第 4 步才從節點解析出來）。每一步失敗就回傳錯誤碼，kernel 放棄綁定並收回 \code{devm_} 配過的東西：

\snip{drv_probe}

第 1 步把暫存器映射進 kernel、第 2 步讀 NBLK 確認是這顆硬體（第 2 章），第 3 步配 DMA buffer（第 3 章），第 4 步申請中斷（第 6 章），第 5 步註冊 \code{/dev/twn_fuse}（第 5 章），最後建立 debugfs（第 8 章）。\code{devm_kzalloc} 配的 \code{struct twn_dev} 用 \code{platform_set_drvdata} 掛在裝置上，之後 \code{remove()} 與 sysfs 用 \code{platform_get_drvdata}／\code{dev_get_drvdata} 拿回來。

\textbf{remove：移除 debugfs 與裝置檔、關掉中斷致能，其餘交給 \code{devm_}。}

\snip{drv_remove}

\textbf{板上：module 帶著比對字串與參數。}\code{modinfo} 印出的 \code{alias} 就是 \code{MODULE_DEVICE_TABLE} 產生的，兩個版本各兩筆；\code{parm} 是三個 module parameter：

\boardfile{H10_driver_state.log}{1-11}

\textbf{板上：開機後已經綁好。}發佈映像的 \code{system-user.dtsi} 把節點改成 driver 認得的字串（\code{iommus} 是第 11 章的 SMMU）：

\snip[{}]{dtsi_fuse0}

開機第 4.8 秒 probe 成功，裝置的 \code{driver} 連結指向 \code{twn_fuse}。開機參數裡還留著 \code{uio_pdrv_genirq.of_id=generic-uio}，但已經沒有節點寫 \code{generic-uio}，UIO 不會接管任何東西：

\boardfile{H3_boot_state.log}{1-13}

\textbf{binding：節點的規格。}一定要有 \code{compatible} 和 \code{reg}；用 \code{if}／\code{then}／\code{else} 規定 \code{-1.1} 一定要有中斷、\code{-1.0} 不准有；\code{additionalProperties: false} 表示沒列出的屬性一律不准：

\snip[{}]{yaml_rules}

\code{test/binding_check.py} 把 PetaLinux 產生的 \code{system.dtb} 反組譯，找出加速器的節點，逐一檢查屬性都有被描述，也檢查 binding 的 \code{compatible} 和 driver 的比對表一致：

\boardfile[：節點原文與每一項檢查]{S3_binding_check.log}{1-3,23-45}

\subsection{原理}

\textbf{從開機到 probe。}U-Boot（開機載入程式）把 \code{system.dtb} 放進記憶體交給 kernel；kernel 把它展開成一棵節點樹，\code{amba_pl} 底下每個節點建成一個 \code{platform_device}，名字是「位址.節點名」（\code{a0000000.twn_fuse}），\code{reg} 變成記憶體資源。\code{interrupts} 這時還留在節點上：driver 呼叫 \code{platform_get_irq_optional()} 時，kernel 才用 \code{of_irq_get()} 把它翻成 Linux 的中斷號碼；中斷控制器還沒準備好，就回 \code{-EPROBE_DEFER}（第 6 章）。driver 註冊時，platform bus 拿它和每個還沒有 driver 的裝置比對，比對函式是 \code{drivers/base/platform.c} 的 \code{platform_match()}，順序是：裝置設了 \code{driver_override} 就只比 driver 名字；否則比 device tree 的 \code{compatible} 與 \code{of_match_table}；接著是 ACPI、\code{id_table}，最後比裝置名稱。比對成功就呼叫 \code{probe()}，回傳 0 才算綁定成功。

\textbf{device tree 的檔案從哪裡來。}PetaLinux 每次建置都從 Vivado 匯出的 XSA 重新產生最上層的 \code{system-top.dts}：它 include 晶片本身的 \code{zynqmp.dtsi}、時脈的 \code{zynqmp-clk-ccf.dtsi}、block design 裡 PL IP 的 \code{pl.dtsi}、PS 設定的 \code{pcw.dtsi}，最後才 include 我們維護的 \code{system-user.dtsi}（在 \code{linux/project-spec/} 底下），所以後者可以改寫前面的內容。\code{.dtsi} 是被 include 的片段，格式和 \code{.dts} 相同，就像 C 的 \code{.h}。\code{dtc}（device tree compiler）把它們編成二進位的 \code{system.dtb}，kernel 只讀這個檔。PetaLinux 為加速器產生的節點寫的是 \code{compatible = "xlnx,twn-fuse-1.1"}，沒有 driver 認得；\code{system-user.dtsi} 用 \code{&fuse_0}（修改既有節點）把它改成 driver 的字串，並加上 \code{iommus}。發佈時 PetaLinux 產生的這幾個檔案，以及反編譯回文字的 \code{system.dtb}，存在 \code{linux/device-tree/}，可以和 \code{system-user.dtsi} 對照著看。

\textbf{\code{compatible} 的格式與版本號。}格式是 \code{"製造商,型號"}，逗號前是 vendor prefix（代表廠商或作者）。Linux 在 \code{vendor-prefixes.yaml} 登記合法的前綴；\code{wjhuang} 不在表上，自己用沒問題，送進 mainline（Linus Torvalds 維護的官方 kernel 原始碼樹）才需要登記。結尾的版本號在這個專案真的用上了：加了中斷的 bitstream 是 \code{-1.1}；舊的 driver 只認得 \code{-1.0}，就不會去等一個它不知道的中斷。

\textbf{\code{devm_} 為什麼好用。}每一筆 \code{devm_*} 配置都登記在裝置上，driver 離開或 \code{probe()} 中途失敗時，kernel 按相反順序全部釋放。所以 \code{probe()} 每一步失敗只要 \code{return} 錯誤碼，不用寫一串 \code{goto} 收拾。

\textbf{binding 檢查的是什麼。}device tree 是資料，不是程式：寫錯一個屬性名稱，\code{dtc} 照樣編過，開機後 driver 才在 \code{probe()} 裡找不到東西。\code{dtc} 不看 binding，要另外跑 \code{make dtbs_check}（或本專案的 \code{binding_check.py}）才能在開機前抓到。\code{status}、\code{phandle}、\code{interrupt-parent} 是每個節點都能有的通用屬性，由 dt-schema 的核心規格描述，個別 binding 不必列。kernel 官方的 \code{make dt_binding_check} 會逐條驗證屬性的格式，這次沒有跑，原因與補法在第 13 章。

\textbf{一個沒處理的限制：PL 的時脈。}節點有 \code{clocks = <&zynqmp_clk 71>}（PL 的 240 MHz 時脈），但 \code{probe()} 沒有取用它。現在能運作，推測是因為開機參數有 \code{clk_ignore_unused}，kernel 不會關掉「沒有 driver 宣告在用」的時脈（沒有試過拿掉它；時脈 driver 從 PMU firmware 拿到的旗標也可能讓這個時脈本來就不被關）。若 PL 時脈被關掉，之後讀暫存器會讓 AXI 匯流排卡住。正確做法是在 \code{probe()} 用 \code{devm_clk_get_enabled()} 宣告使用；這次沒有做。

\subsection{面試會問}

\begin{qa}{compatible 字串怎麼讓 kernel 找到你的 driver？你怎麼確定節點寫對了？}
開機時 kernel 把 device tree 展開，PL 底下每個節點變成一個 platform device，\code{reg} 變成它的記憶體資源，中斷等 probe 時才從節點解析。我的 driver 有一張 \code{of_match_table}，module 載入時向 platform bus 註冊，bus 用 \code{platform_match()} 比對：先看 \code{driver_override}，再比 compatible。對上了就呼叫我的 \code{probe()}：ioremap 暫存器、配 DMA buffer、申請中斷、註冊 \code{/dev/twn_fuse}，全部用 \code{devm_}，失敗時自動收回。compatible 帶版本號，加了中斷的硬體是 \code{-1.1}。節點的規格我寫成 dt-schema 的 binding，規定 \code{-1.1} 一定要有中斷，並拿 PetaLinux 實際產生的 dtb 去比對，因此發現它多產生了 \code{interrupt-names}。
\end{qa}

\subsection{關鍵函式速查}

\begin{center}
\small
\begin{tabularx}{\textwidth}{l l X}
\toprule
名稱 & 位置 & 做什麼 \\
\midrule
\code{twn_of_match} & \code{twn_fuse.c}:\lineof{drv_of_match} & 比對表：\code{-1.1} 與 \code{-1.0} \\
\code{twn_driver} & \code{twn_fuse.c}:\lineof{drv_driver} & 把 probe、remove、比對表、sysfs 屬性綁在一起 \\
\code{twn_probe()} & \code{twn_fuse.c}:\lineof{drv_probe} & 接管硬體的五個步驟 \\
\code{twn_remove()} & \code{twn_fuse.c}:\lineof{drv_remove} & 移除 debugfs 與裝置檔、關中斷致能，其餘交給 \code{devm_} \\
\code{&fuse_0} & \code{system-user.dtsi}:\lineof{dtsi_fuse0} & 發佈映像的節點：compatible 與 \code{iommus} \\
\code{required}、\code{allOf} & \code{wjhuang,twn-fuse.yaml}:\lineof{yaml_rules} & binding 的必要屬性與依版本的規則 \\
\code{platform_match()} & kernel \code{drivers/base/platform.c} & 比對順序：override、compatible、ACPI、id table、名稱 \\
\bottomrule
\end{tabularx}
\end{center}

\clearpage
\section{MMIO：ioremap 與 readl / writel}

\begin{oneline}
\todo{一句話講完這一節}
\end{oneline}

\begin{before}
\todo{v1.0.0 怎麼做，引用 sw/app/pl\_hal\_fuse.c 的實際程式碼}
\end{before}

\begin{tech}
\todo{用到的 kernel API 與機制}
\end{tech}

\begin{gain}
\todo{具體改善了什麼，能驗證的寫法}
\end{gain}

\subsection{實際例子：程式碼與板上輸出}
\todo{\textbackslash snip 抽出的函式，逐段說明；板上實際 log}

\subsection{原理}
\todo{為什麼這樣做是對的；必要時一張圖}

\subsection{面試會問}
\begin{qa}{MMIO mapping：為什麼要先 ioremap，不能直接用指標存取 0xA000\_0000？}
\todo{30 秒版本的回答}
\end{qa}

\subsection{關鍵函式速查}
\todo{函式、在哪個檔案第幾行、做什麼}

\needspace{20\baselineskip}
\section{DMA buffer：配置、cache 與記憶體屬性}
\label{sec:dma}

\begin{oneline}
PL 讀寫的 DDR buffer 由 driver 透過 DMA API 配置，kernel 同時給出 CPU 用的位址與 PL 用的 DMA 位址，並讓 CPU 以 Normal non-cacheable 屬性存取它；在同一組開機檔上，這個屬性就是 driver 版本比 UIO 版本快約 1.1\% 的原因。
\end{oneline}

\begin{before}
device tree 保留一塊固定的 256 MB，標成 \code{no-map}（kernel 完全不碰）：

\snip[{}]{v100_dt_reserved}

app 再經 \code{/dev/mem}（第 148 行），用寫死的實體位址 \code{0x6000_0000}（第 140 行的 \code{FUSE_ARENA_PA_DEFAULT}）把整塊 \code{mmap} 進來：

\snip{v100_uio_arena_init}

問題：位址寫死在程式與 device tree 兩個地方；256 MB 只用到約 6 MB，其餘 kernel 也不能用；\code{/dev/mem} 能讀寫\textbf{任何}實體位址（第 10 章）；映射出來的屬性是 Device-nGnRnE，寫入不能合併、不能非對齊存取，連 \code{memset} 都不能用。
\end{before}

\begin{tech}
\begin{itemize}
\item \code{dma_set_mask_and_coherent(dev, DMA_BIT_MASK(40))}：告訴 kernel 這個裝置能定址 40 bits。
\item \code{dmam_alloc_coherent()}：配一塊 CPU 與裝置「一致」（coherent）的記憶體，回傳 CPU 位址與 \code{dma_addr_t}；\code{dmam_} 表示 driver 離開時自動歸還。
\item CMA（Contiguous Memory Allocator，連續記憶體配置器）：開機時劃出的一塊區域（這裡 256 MB），平時借給可以搬走的一般用途，driver 要大塊實體連續的記憶體時再清出來。
\item \code{dma_addr_t}：裝置看到的位址。沒有 IOMMU 時等於實體位址；有 SMMU 時是 I/O 虛擬位址（IOVA，第 11 章），程式不該假設兩者相同。
\end{itemize}
\end{tech}

\begin{gain}
\begin{itemize}
\item 位址不再寫死：每次載入由 kernel 決定，driver 再寫進硬體。driver 只配 8 MB（module 參數 \code{arena_mb}），app 需要 6 MB；256 MB 保留區拿掉了（第 10 章）。
\item 卸載時自動歸還：沒有 SMMU 時連續載入 10 次、每次 64 MB（共 640 MB，CMA 只有 256 MB），10 次都成功。
\item CPU 端屬性從 Device 變成 Normal non-cacheable，寫入可以合併：CPU 寫 buffer 快 3 到 5 倍，pipeline 的「寫輸入」從 \splitDevIn{} ms 縮到 \splitNcIn{} ms（每次執行印出的最後一幀，各 3 次），fps 高 1.2\%（開機組 D，只差這個屬性的兩組）。可以 \code{memset} 這一點是依 ARM 規格推得，app 沒有改用，也沒有在板上測。
\end{itemize}
\end{gain}

\subsection{實際例子}

\textbf{probe 的第 3 步。}先設 DMA mask，再配 buffer，最後檢查對齊：

\snip{drv_probe_dma}

\code{DMA_BIT_MASK(40)} 對應 CSR 的 40-bit 位址暫存器；kernel 配記憶體時保證位址在這個範圍內。\code{dmam_alloc_coherent} 的回傳值 \code{td->arena} 是 CPU 位址（kernel 用），\code{&td->arena_dma} 收到 DMA 位址（寫進 PL 的暫存器）。64 KB 對齊是因為 app 的配置以 64 KB 為單位（head 的兩個輸入 slot 相隔 \code{0x10000}）。

\textbf{8 MB 裡放什麼。}buffer 的配置由 app 決定，driver 只管「給一塊、檢查 offset 不越界」。\code{FUSE_ARENA_NEED} 是 6 MB，後面還有一個 \code{_Static_assert}，配置長大超過 6 MB 就編譯失敗：

\snip{app_arena_map}

依序是：backbone 權重、兩個 activation slot（17 個 block 輪流使用）、輸入影像、backbone 輸出；\code{0x400000} 起是 head 的權重、兩個輸入 slot、三層中間結果與兩個輸出 slot。

\textbf{板上：buffer 從哪裡來，取決於有沒有 SMMU。}關掉 SMMU 時（開機組 I），DMA 位址 \code{0x6a10_0000} 落在 CMA（從 \code{0x6a00_0000} 起 256 MB）裡，\code{CmaFree} 少了 8 MB 多；連續 10 次各要 64 MB，超過 CMA 的總量也都成功，表示每次卸載都真的歸還了：

\boardfile{I1_driver_state_nosmmu.log}{1,19-24,56-60}

發佈映像打開了 SMMU（開機組 H）：DMA 位址變成 IOVA \mbox{\code{0xff_ff80_0000}}，\code{CmaFree} 幾乎沒少，因為 SMMU 可以把分散的實體頁排成連續的 IOVA，DMA API 改從一般記憶體配頁（本章「原理」）：

\boardfile{H10_driver_state.log}{19-24}

\begin{pit}
一開始用 \code{/proc/meminfo} 的 \code{CmaFree} 檢查有沒有漏記憶體，卸載後數字卻沒有回升，看起來像漏掉了。其實不是：CMA 沒人用時會借給 page cache（kernel 快取檔案內容的記憶體）這類可搬移的用途，\code{CmaFree} 會隨讀檔變化，不能當證據。後來改用「載入總量超過 CMA」的方法，才得到能下結論的證據。
\end{pit}

\textbf{量測：記憶體屬性值多少。}driver v1.0 剛完成時比 UIO 版本快約 0.8\%（開機組 A，第 13 章），但當時一起改了三件事：等硬體的地方、buffer 位置、CPU 端屬性。在 v1.0.0 的開機檔上（開機組 D）做四組實驗，每組 pipeline 跑 6 次、分兩輪且順序相反：

\begin{center}
\small
\begin{tabularx}{\textwidth}{l l l l X}
\toprule
組別 & 等硬體的地方 & buffer 位置 & CPU 端屬性 & 跟上一組的差別 \\
\midrule
\code{uio} & app & \code{0x6000_0000} & Device-nGnRnE & 起點（v1.0.0） \\
\code{uio_68} & app & \code{0x6800_0000} & Device-nGnRnE & buffer 位置 \\
\code{drv_dev} & driver & CMA & Device-nGnRnE & 等硬體的地方與 buffer 來源（一起換） \\
\code{drv_nc} & driver & CMA & Normal non-cacheable & 只有 CPU 端屬性 \\
\bottomrule
\end{tabularx}
\end{center}

\code{drv_dev} 靠 driver 的 module 參數 \code{map_device=1}：\code{twn_mmap()} 改用 \code{/dev/mem} 那種 Device 屬性映射（第 4 章）。

\boardfile[，第 27--31 行：每組 6 次的平均]{D1_ab_measure.log}{1-2,27-31}

前三組都在 228.8 到 229.1 fps，只有 \code{drv_nc} 跳到 \abDrvNc{} fps。直接量 CPU 存取 buffer 的速度（\code{twn_bench} 連續寫、讀 1 MB），寫入差很多，讀取一樣：

\boardfile[：第 11 行是 Normal non-cacheable，第 20 行是 \code{map_device=1}]{D2_bench.log}{11-12,20}

pipeline 開 \code{PIPE_TIMING=1}，每次執行（120 幀）印出最後一幀每一步的時間：只有 \code{write-in}（CPU 把輸入影像寫進 buffer）變了，從 \splitDevIn{} ms 變成 \splitNcIn{} ms，和這三次執行平均的每幀週期差（約 0.04 ms）相當：

\boardfile[：前三次是 Normal non-cacheable，後三次是 \code{map_device=1}]{D6_stage_times.log}{3-21}

\subsection{原理}

\textbf{為什麼會有 cache 問題。}CPU 讀寫 DDR 會經過 cache；PL 經 HP0--HP3 port（PL 存取 DDR 的四個高速 AXI 埠）直接到 DDR 控制器，\textbf{不會}去查 CPU 的 cache。buffer 若是 cacheable：CPU 寫的輸入可能還停在 cache，PL 讀到舊資料；PL 寫進 DDR 的輸出，CPU 可能命中 cache 裡上一幀的內容。處理方法有三種：

\begin{center}
\small
\begin{tabularx}{\textwidth}{l X X}
\toprule
做法 & 怎麼做 & 代價 \\
\midrule
coherent buffer（這次用的） & \code{dma_alloc_coherent}：CPU 端映射成 non-cacheable，讀寫直接到 DDR & CPU 存取比 cacheable 慢；app 只寫輸入、讀輸出，量不大 \\
streaming DMA & 一般 cacheable 記憶體；PL 讀之前 \code{dma_sync_single_for_device}（寫回 cache），PL 寫完 \code{dma_sync_single_for_cpu}（作廢 cache） & 每幀都要做 cache 維護，app 也要配合 \\
硬體一致性 & 改用 HPC port，PL 的存取經過 CCI（會查 CPU cache 的互連）；AXI master 送 cacheable、shareable 的屬性；節點加 \code{dma-coherent} & 要改 RTL 與 block design、重建 bitstream \\
\bottomrule
\end{tabularx}
\end{center}

kernel 不知道 HP port 是什麼，它看 device tree：節點沒有 \code{dma-coherent}，就把裝置當成 non-coherent，\code{dma_alloc_coherent} 給出 non-cacheable 映射。

\textbf{\code{dma_alloc_coherent} 實際做了什麼。}沒有 SMMU 時：從 CMA 拿實體連續的頁（一般的配置器一次最多給 4 MB，而且久了記憶體會零碎，不保證拿得到）；把這些頁在 cache 裡的內容寫回並作廢；另外映射一個 Normal non-cacheable 的 kernel 位址回傳；DMA 位址等於實體位址。有 SMMU 時：PL 看到的是 IOVA，實體頁不必連續，DMA API 從一般記憶體配頁，再由 SMMU 排成一段連續的 IOVA，所以不佔 CMA。driver 的程式兩種情況完全一樣。

\textbf{為什麼不用 \code{kmalloc}。}它的記憶體在 kernel 的線性映射（linear map：kernel 頁表裡把整個 DDR 一對一對應進來的那一段，cacheable）裡，要給 PL 就得自己做 streaming DMA；一次最多 4 MB；也不給 \code{dma_addr_t}，而有 SMMU 時實體位址根本不是 PL 該用的位址。

\textbf{兩種「不經過 cache」的記憶體。}arm64 頁表的每一頁都帶屬性。這裡牽涉兩種，CPU 存取都不經過 cache，差別在寫入：

\begin{center}
\small
\begin{tabularx}{\textwidth}{l X X}
\toprule
 & Device-nGnRnE & Normal non-cacheable \\
\midrule
誰會得到 & \code{/dev/mem} 映射 no-map 區、\code{pgprot_noncached()} & \code{dma_mmap_coherent()}、\code{pgprot_dmacoherent()} \\
寫入 & 每個 store 都是一筆獨立的匯流排交易，不合併（nG）；回應必須由終點發出（nE） & 連續的小 store 可以在 write buffer 裡合併成一筆大的寫入（write gathering） \\
讀取 & 每次都到 DDR & 每次都到 DDR \\
對齊 & 不能非對齊存取 & 可以 \\
\bottomrule
\end{tabularx}
\end{center}

所以寫入差 3 到 5 倍、讀取一樣，和量到的一致。讀取 81 MB/s 表示每 8 bytes 約 100 ns，大約是一趟 DDR 讀取的延遲，兩種屬性都要付。只有 \code{write-in} 反映到 fps，因為寫輸入在 producer 執行緒的關鍵路徑上（寫完才能啟動 PL）；讀 heatmap 在另一個執行緒，和 PL 重疊。

\textbf{實驗的限制。}\code{uio_68} 只換位置、fps 不變，排除了「buffer 在 DDR 哪裡」；但 \code{uio_68}→\code{drv_dev} 一次換了兩件事，要分開還需要一組「driver 版、buffer 在保留區」，這次沒有做，所以只能說「輪詢位置與 buffer 來源合起來看不出影響」。

\subsection{面試會問}

\begin{qa}{為什麼用 dma\_alloc\_coherent？HP port 不是 coherent，cache 怎麼處理？driver 為什麼比較快？}
\code{dma_alloc_coherent} 一次給我 CPU 位址和裝置用的 DMA 位址。不用 \code{kmalloc}，因為它是 cacheable、最大 4 MB、也不給 DMA 位址；不用 \code{/dev/mem}，因為它能讀寫整個實體記憶體。HP port 不查 CPU 的 cache，節點沒有 \code{dma-coherent}，kernel 就把 buffer 映射成 non-cacheable。沒有 SMMU 時它從 CMA 拿連續的頁，有 SMMU 時從一般記憶體配、由 SMMU 排成連續的 IOVA。driver 版本比 UIO 版本快約 1\%（同一組開機檔上 1.1\%），我做四組 A/B 實驗，把記憶體屬性單獨拆出來，只差這一項的兩組差 1.2\%：原因是 CPU 端屬性從 Device 變成 Normal non-cacheable，寫入可以合併，CPU 寫 buffer 快 3 到 5 倍，寫輸入從 \splitDevIn{} 降到 \splitNcIn{} ms。
\end{qa}

\subsection{關鍵函式速查}

\begin{center}
\small
\begin{tabularx}{\textwidth}{l l X}
\toprule
名稱 & 位置 & 做什麼 \\
\midrule
probe 第 3 步 & \code{twn_fuse.c}:\lineof{drv_probe_dma} & 40-bit mask、配 buffer、檢查對齊 \\
\code{arena_mb} & \code{twn_fuse.c}:\lineof{drv_param_arena} & buffer 大小（MB），預設 8 \\
buffer 配置 & \code{mbv2_fuse_ps.c}:\lineof{app_arena_map} & app 在 8 MB 裡的區段 \\
\code{run()} & \code{test/ab_measure.sh}:\lineof{ab_run} & A/B 實驗：切 driver、跑 pipeline 三次 \\
第 3 部分 & \code{test/twn_bench.c}:\lineof{bench_copy} & CPU 寫、讀 buffer 的頻寬 \\
\code{dmam_alloc_coherent} & kernel & coherent buffer，driver 離開時自動歸還 \\
\bottomrule
\end{tabularx}
\end{center}

\needspace{20\baselineskip}
\section{mmap：只把 buffer 給 user space}
\label{sec:mmap}

\begin{oneline}
driver 的 \code{mmap} 只把 DMA buffer 那幾頁放進 app 的頁表，app 不必複製資料（zero-copy）就能讀寫 buffer，而暫存器從來不經由 driver 映射出去。
\end{oneline}

\begin{before}
v1.0.0 的 app 自己做兩次映射：暫存器（\code{/dev/uio4} 的 64 KB，第 1 章「原本怎麼做」的 \code{fuse_hal_init()}）和整塊實體記憶體（\code{/dev/mem} 從 \code{0x6000_0000} 起 256 MB，第 3 章「原本怎麼做」的 \code{fuse_arena_init()}）。板上 app 執行時，它的位址空間最後兩行就是這兩個映射；app 的任何一個錯誤指標，都可能寫到暫存器或那 256 MB 的任何地方：

\boardfile{A4_uio_state.txt}{5-12}
\end{before}

\begin{tech}
\begin{itemize}
\item \code{file_operations.mmap}：app 呼叫 \code{mmap()} 時，kernel 先建立一段位址範圍，再呼叫 driver 的這個函式，由 driver 決定這段範圍對應到哪些實體頁、用什麼屬性。
\item \code{struct vm_area_struct}（VMA）：描述這段範圍，包括起點 \code{vm_start}、終點 \code{vm_end}，以及 app 要求的 offset \code{vm_pgoff}（以頁為單位）。
\item \code{dma_mmap_coherent()}：把 \code{dma_alloc_coherent} 配的 buffer 映射到 user space，屬性和 kernel 端一樣是 non-cacheable。
\end{itemize}
\end{tech}

\begin{gain}
\begin{itemize}
\item app 的位址空間裡只剩一段 8 MB 的 \code{/dev/twn_fuse}，沒有暫存器，也沒有 \code{/dev/mem}。
\item 要求超過 buffer 大小，或 offset 不是 0，都會被拒絕（\code{-EINVAL}）。
\item 仍然是 zero-copy：輸入影像與 heatmap 不必經過 \code{read()}／\code{write()} 複製；CPU 端是 Normal non-cacheable，寫入可以合併（第 3 章）。
\end{itemize}
\end{gain}

\subsection{實際例子}

\textbf{driver 端：只映射 buffer。}先檢查 offset 與長度，再交給 DMA API。中間的 \code{map_device} 分支只給第 3 章的 A/B 實驗用：它改用 \code{remap_pfn_range()} 加 \code{pgprot_noncached()}，得到和 \code{/dev/mem} 一樣的 Device 屬性；有 SMMU 時 buffer 沒有固定的實體位址可以這樣映射，所以拒絕：

\snip{drv_mmap}

\code{file_to_td(f)} 從 \code{struct file} 找回 driver 的 \code{struct twn_dev}（第 5 章）。

\textbf{app 端：問大小、映射、回報「基底位址 0」。}

\snip{app_arena_init}

\code{*phys = 0}：回報基底位址 0，app 用「基底加 offset」算出的位址就變成 offset，正好是 ioctl 要的東西（為什麼這樣 app 就不用改，見第 9 章）。

\textbf{板上：app 的位址空間只剩 buffer。}pipeline app 執行中，看它的 \code{/proc/<pid>/maps} 與開著的檔案：

\boardfile{H10_driver_state.log}{56-58}

\code{0xffff95afd000 - 0xffff952fd000 = 0x800000}，正好 8 MB。和前面 v1.0.0 的輸出相比，\code{/dev/mem} 與 \code{/dev/uio4} 都不見了。越界的要求被拒絕：

\boardfile{H4_board_check.log}{10-12}

\subsection{原理}

\textbf{\code{mmap} 做了什麼。}app 呼叫 \code{mmap(NULL, len, ..., fd, offset)} 時，kernel 在 app 的位址空間找一段空的範圍，建立 VMA，記下長度與 offset（\code{vm_pgoff = offset / 4096}），再呼叫 driver 的 \code{.mmap} 填頁表。driver 回傳錯誤碼，整個 \code{mmap()} 就失敗。

\textbf{\code{dma_mmap_coherent} 的內部。}沒有 SMMU 時走 \code{kernel/dma/direct.c} 的 \code{dma_direct_mmap}：用 \code{dma_pgprot} 把屬性改成 non-cacheable（同一塊記憶體若被 kernel 和 user 用不同的 cache 屬性映射，CPU 可能看到不一致的資料），檢查範圍不超出 buffer，再用 \code{remap_pfn_range} 把這段連續的實體頁一次填進頁表。發佈映像打開了 SMMU，buffer 的實體頁不連續（第 3 章），改走 \code{drivers/iommu/dma-iommu.c} 的 \code{iommu_dma_mmap}，逐頁映射。對 driver 來說都是同一個呼叫。

\textbf{app 為什麼碰不到暫存器。}app 的位址空間裡有什麼，完全由 kernel 建的頁表決定。CSR 只在 kernel 裡做過 \code{ioremap}，driver 的 \code{.mmap} 只映射 buffer，所以 app 的頁表裡沒有任何一項指向 \code{0xA000_0000}。剩下的旁路是 root 自己開 \code{/dev/mem}：發佈映像用 \code{STRICT_DEVMEM} 與 \code{IO_STRICT_DEVMEM} 擋住 System RAM 與 driver 佔用中的暫存器，但有邊界（第 10 章）。

\textbf{映射的生命週期。}VMA 持有 \code{struct file} 的參考。就算 app 關掉 fd，只要映射還在，檔案就不算關閉，module 也就不能卸載（第 7 章的板上測試）。這保證 buffer 不會在 app 還在用的時候被釋放。

\subsection{面試會問}

\begin{qa}{app 為什麼只看得到 buffer，看不到暫存器？為什麼用 mmap 不用 read／write？}
app 的頁表完全由 kernel 決定。我的 driver 的 \code{.mmap} 只把 \code{dma_alloc_coherent} 配的 buffer 用 \code{dma_mmap_coherent} 映射給 app，並檢查長度不超過 buffer、offset 必須是 0；暫存器只在 kernel 裡 \code{ioremap}，從來沒有映射出去。板上 \code{/proc/<pid>/maps} 可以看到：UIO 版本有 256 MB 的 \code{/dev/mem} 和 64 KB 的 \code{/dev/uio4}，driver 版本只剩 8 MB 的 \code{/dev/twn_fuse}。資料走 mmap 是因為每幀幾百 KB，用 \code{read}／\code{write} 要經 \code{copy_from_user}／\code{copy_to_user} 多複製兩次；mmap 讓 app 和 PL 看同一塊記憶體。
\end{qa}

\subsection{關鍵函式速查}

\begin{center}
\small
\begin{tabularx}{\textwidth}{l l X}
\toprule
名稱 & 位置 & 做什麼 \\
\midrule
\code{twn_mmap()} & \code{twn_fuse.c}:\lineof{drv_mmap} & 檢查 offset 與長度，映射 buffer \\
\code{map_device} 的檢查 & \code{twn_fuse.c}:\lineof{drv_mmap_iommu} & 有 SMMU 時拒絕 Device 映射 \\
\code{fuse_arena_init()} & \code{pl_hal_fuse_drv.c}:\lineof{app_arena_init} & 問大小、\code{mmap}、回報基底 0 \\
\code{dma_mmap_coherent} & kernel & 映射 coherent buffer，屬性與 kernel 端一致 \\
\bottomrule
\end{tabularx}
\end{center}

\needspace{20\baselineskip}
\section{file operations 與 ioctl：控制介面、copy\_from\_user 與 offset}
\label{sec:ioctl}

\begin{oneline}
\code{/dev/twn_fuse} 只實作 \code{open}、\code{release}、\code{mmap}、\code{unlocked_ioctl} 四個 file operation：資料走 mmap，指令走 ioctl，參數一律用 \code{copy_from_user} 取，DDR 位址只收 buffer 內的 offset。
\end{oneline}

\begin{before}
app 自己寫暫存器，寫進去的是實體位址（\code{g_arena_pa + OFF_IN} 這種算法）。\code{wr64()} 直接把它拆成兩個 word 寫進 PL（第 2 章「原本怎麼做」的 \code{pl_hal_fuse.c} 第 80 行）。PL 拿到什麼位址就讀寫什麼位址，不經過任何檢查；app 算錯一個位址，或被惡意修改，PL 就會改寫 kernel 或其他程式的記憶體。每個 app 也都要自己記住「先等 idle、清 done、寫位址、再寫 START」的順序。
\end{before}

\begin{tech}
\begin{itemize}
\item misc device：主裝置號固定是 10、次裝置號自動分配的字元裝置，一個 \code{misc_register()} 就出現 \code{/dev/twn_fuse}。
\item \code{struct file_operations}：把 \code{open}、\code{mmap}、\code{ioctl} 等系統呼叫接到 driver 的函式。
\item \code{_IO}／\code{_IOR}／\code{_IOW}／\code{_IOWR}：把方向、參數大小、類別、編號編進一個 32-bit 指令碼的巨集。
\item uapi header \code{twn_fuse_ioctl.h}：driver 與 app 共用，只用 \code{__u32}、\code{__u64} 這種兩邊都有的固定寬度型別。
\item \code{copy_from_user}／\code{copy_to_user}／\code{memdup_user}：在 user 與 kernel 記憶體之間安全地搬資料；\code{in_arena()}：driver 自己的範圍檢查。
\item \code{void __user *}：標記「這是 user space 的指標」的型別，和第 2 章的 \code{__iomem} 一樣是給靜態檢查工具 sparse 看的，直接解這種指標會被它擋下。
\end{itemize}
\end{tech}

\begin{gain}
\begin{itemize}
\item 每個指令就是一段完整、順序固定的暫存器操作，app 不可能漏掉某一步。
\item 寫進 PL 的位址一定落在 driver 的 buffer 裡；head 的傳輸長度已知，連「從起點到終點」整段都檢查。
\item 錯誤用標準錯誤碼回報：越界回 \code{EINVAL}、壞指標回 \code{EFAULT}、不認得的指令回 \code{ENOTTY}。
\item 介面本身幾乎不花時間：一次 ioctl 來回 \ioctlUs{} µs；跑一幀時，系統呼叫、拿鎖、參數複製與讀回六個 cycle counter 合計 \bbSysUs{} µs，不到幀時間的 0.1\%。
\end{itemize}
\end{gain}

\subsection{實際例子}

\textbf{四個 file operation。}沒有 \code{.read}／\code{.write}：app 呼叫 \code{read()} 會直接得到 \code{-EINVAL}。

\snip{drv_fops}

\code{open} 用 \code{test_and_set_bit} 讓一次只有一個 process 能打開（第 7 章）；\code{file_to_td()} 用 \code{container_of}，從 \code{misc_register} 註冊的 \code{struct miscdevice} 反推回整個 \code{struct twn_dev}，所有 file operation 都靠它找到暫存器、buffer 與鎖：

\snip{drv_file_to_td}

\textbf{介面定義：九個指令。}\code{0xBC} 是這個 driver 的類別（magic），Linux 6.1 的 \code{ioctl-number.rst} 裡沒有人登記；第三個參數是型別，巨集用 \code{sizeof} 把大小編進指令碼：

\snip{uapi_cmds}
\snip{uapi_bb_run}

\textbf{最簡單的完整範例：設定 base。}複製參數、檢查範圍、寫暫存器：

\snip{drv_set_bases}

\textbf{範圍檢查。}起點 16-byte 對齊（PL 一次搬一個 16-byte 的 beat）、起點在 buffer 內、長度不超出 buffer 結尾。寫成 \code{len <= size - off} 而不是 \code{off + len <= size}，是為了避免 \code{off + len} 在 64 bits 溢位、繞回一個很小的數而通過檢查：

\snip{drv_in_arena}

\textbf{head 的描述子：長度已知，整段都檢查。}app 已經把形狀、旗標打包好，driver 只檢查三個位址範圍、寫進去、commit。operator 0 的輸入多保留 64 KB，因為 head 會在兩個輸入 slot 之間切換：

\snip{drv_hd_commit}

描述子表（17 個 block 各 16 個 word）用 \code{memdup_user} 一次複製進來（\code{twn_set_table()}，第 \lineof{drv_set_table} 行），它等於 \code{kmalloc} 加 \code{copy_from_user}，用完要 \code{kfree}。

\textbf{分派：一個 switch、兩把鎖。}\code{STATUS} 只讀六個暫存器、不拿鎖；backbone 與 head 的指令各拿自己的 mutex（第 7 章）；不認得的指令回 \code{-ENOTTY}：

\snip{drv_ioctl}

\textbf{板上：該拒絕的都拒絕了。}\code{Bad address} 就是 \code{EFAULT}：測試程式故意把參數指標設成 16（沒有映射的位址），kernel 回報錯誤而沒有當掉：

\boardfile{H4_board_check.log}{1,13-20}

\textbf{板上：ioctl 要花多少時間。}\code{twn_bench} 第 1 部分把不拿鎖的 \code{STATUS} 連打 100,000 次；第 2 部分同一幀跑 1000 次，比較三個時鐘：app 在 ioctl 前後讀的時間、driver 在 kernel 裡量的時間（拿到參數、檢查完之後開始計時，清完 done 停錶；介面標頭的註解寫成「START -> done」是簡寫）、PL 自己數的 cycle 除以 240 MHz：

\boardfile[：用中斷時]{E3_bench.log}{6-10}

app 比 driver 多的 \bbSysUs{} µs，是系統呼叫進出、拿鎖、\code{copy_from_user}／\code{copy_to_user}，以及讀回六個 cycle counter：driver 在 \code{twn_bb_run()} 裡先停錶、才讀 counter（第 6 章的程式）。同樣讀六個暫存器的 \code{STATUS} 整趟只要 \ioctlUs{} µs，所以這 6 次 MMIO 讀取最多約佔 \bbSysUs{} µs 的一半，這次沒有拆開量。driver 比 PL 多的 \irqWaitUs{} µs 主要是「等」的成本，另含 START 前後的幾次暫存器存取與 block 之間的空檔，第 6 章比較輪詢與中斷時用的就是它。

\subsection{原理}

\textbf{為什麼沒有 read／write。}\code{read}／\code{write} 是位元組串流的介面，適合一段段流進流出的資料（序列埠、感測器 FIFO）。這顆加速器兩件事都不是：資料每幀幾百 KB，用 \code{read}／\code{write} 要多複製兩次，所以走 mmap；指令有結構（設表、設位址、啟動並等待），用 \code{ioctl} 加固定格式的 struct 最直接。kernel 提供的其他介面各有用途：

\begin{center}
\small
\begin{tabularx}{\textwidth}{l X l}
\toprule
介面 & 適合 & 這個 driver \\
\midrule
\code{read}／\code{write} & 串流資料；可以用 \code{cat}／\code{echo} 操作 & 沒用 \\
\code{ioctl} & 對裝置下動作，參數是結構化的 & 九個指令 \\
\code{mmap} & 大量資料，零複製 & 8 MB buffer \\
sysfs & 一個檔案一個值的屬性，給人和腳本；算正式的介面約定（ABI） & 狀態與計數器（第 8 章） \\
debugfs & 只給開發者的除錯介面，不保證穩定 & 暫存器傾印、短幀（第 8 章） \\
\bottomrule
\end{tabularx}
\end{center}

\textbf{指令碼怎麼編。}\code{_IOC} 把 32-bit 數字切成四段：bits 31--30 是方向，29--16 是參數大小，15--8 是類別，7--0 是編號。\code{TWN_IOC_BB_RUN} = \code{0xC030BC03}：讀加寫、48 bytes、類別 \code{0xBC}、編號 3。struct 的大小一改，指令碼就不同，舊 app 送來的指令 driver 不認得而回 \code{ENOTTY}，不會讀錯資料；但只改欄位順序、大小不變時偵測不到。

\textbf{為什麼一定要用 \code{copy_from_user}。}ioctl 的 \code{arg} 是 app 位址空間裡的指標，kernel 不能直接解：位址可能沒有映射，直接解會 oops，\code{copy_from_user} 遇到這種 fault 由 kernel 的 exception table（記錄「出錯時跳到哪裡修復」的表）接手並回傳失敗，driver 據此回 \code{-EFAULT}；位址也可能故意指向 kernel，\code{copy_from_user} 先用 \code{access_ok} 確認整段都在 user 範圍。ARMv8.1 起有 PAN（Privileged Access Never），kernel 直接存取 user 記憶體會發生例外；這塊板子的 Cortex-A53 是 ARMv8.0，沒有 PAN，也沒有開軟體模擬，所以更要靠 \code{copy_from_user}。struct 裡要帶指標時（描述子表的 \code{words}），型別用 \code{__u64}，32 與 64 位元的 app 看到的版面才一樣。

\textbf{為什麼只收 offset。}PL 的 DMA 拿到什麼位址就存取什麼。如果 driver 接受 app 給的實體位址，等於讓任何能開 \code{/dev/twn_fuse} 的程式控制一個能讀寫全部記憶體的 DMA 引擎。只收 offset，檢查後加上 driver 自己的 DMA 位址，PL 收到的位址就一定在 buffer 裡。這層檢查管的是「driver 寫進暫存器的起點」：head 的三個位址連長度都檢查；backbone 只檢查起點，輸出寫多長取決於描述子表，driver 不逐項解讀；描述子表本身由 PL 解讀，欄位理論上仍可讓 PL 越界。PL 的每一筆存取都檢查，要靠 SMMU（第 11 章）。

\textbf{介面沒有實作的部分。}沒有 \code{.compat_ioctl}：struct 都是固定寬度型別，32 位元的 app 要用時，填 kernel 現成的 \code{compat_ptr_ioctl} 即可，這次沒有做（app 是 64 位元）。沒有 \code{.poll}：需要等待的指令都會在 ioctl 裡睡，app 要同時做別的事就用執行緒，或把「啟動」和「等待」拆成兩個指令（head 的 \code{HD_START}／\code{HD_WAIT} 就是這樣）。

\subsection{面試會問}

\begin{qa}{ioctl 和其他介面怎麼選？為什麼要 copy\_from\_user、只收 offset？}
我依資料的性質選：大量資料用 mmap 零複製；對裝置的動作和結構化參數用 ioctl，指令碼用 \code{_IOWR} 把方向和大小編進去；給人看的狀態用 sysfs，除錯用 debugfs。所以我的 driver 只有 open、release、mmap、ioctl，沒有 read／write。ioctl 的參數是 user 指標，可能沒映射、也可能故意指向 kernel，所以一律 \code{copy_from_user}，失敗回 \code{-EFAULT}；A53 沒有 PAN，硬體不會幫我擋。位址只收 buffer 內的 offset，檢查對齊與範圍，連 offset 加長度的溢位都擋掉，再加上 DMA 位址寫進硬體，app 給不出任意實體位址。板上量到一次 ioctl \ioctlUs{} µs，跑一幀的系統呼叫、複製加讀回 counter 合計 \bbSysUs{} µs，介面不是瓶頸。
\end{qa}

\subsection{關鍵函式速查}

\begin{center}
\small
\begin{tabularx}{\textwidth}{l l X}
\toprule
名稱 & 位置 & 做什麼 \\
\midrule
\code{twn_fops} & \code{twn_fuse.c}:\lineof{drv_fops} & 四個 file operation \\
九個指令碼 & \code{twn_fuse_ioctl.h}:\lineof{uapi_cmds} & driver 與 app 共用的介面 \\
\code{in_arena()} & \code{twn_fuse.c}:\lineof{drv_in_arena} & 對齊、起點、長度，含溢位防護 \\
\code{twn_set_bases()} & \code{twn_fuse.c}:\lineof{drv_set_bases} & 最簡單的完整 ioctl \\
\code{twn_set_table()} & \code{twn_fuse.c}:\lineof{drv_set_table} & \code{memdup_user} 複製描述子表 \\
\code{twn_hd_commit()} & \code{twn_fuse.c}:\lineof{drv_hd_commit} & head 描述子：三段範圍完整檢查 \\
\code{twn_ioctl()} & \code{twn_fuse.c}:\lineof{drv_ioctl} & 依指令分派、上鎖 \\
第 1、2 部分 & \code{test/twn_bench.c}:\lineof{bench_status} & ioctl 與一幀的時間 \\
\bottomrule
\end{tabularx}
\end{center}

\needspace{20\baselineskip}
\section{等待完成：從輪詢到中斷}
\label{sec:irq}

\begin{oneline}
driver 啟動一幀後要等 PL 做完：v1.1 在 RTL 加了一根 level 中斷線，driver 睡在 completion 上、由中斷處理函式（ISR）叫醒；沒有中斷的舊硬體則退回每 13--50 µs 讀一次 done 位元的輪詢，兩者都有 3 秒逾時。
\end{oneline}

\begin{before}[原本怎麼做（v1.0.0 的 app 輪詢，driver v1.0 的 kernel 輪詢）]
v1.0.0 的 app 在 user space 做完整個流程：清殘留的 done、等 busy 結束、寫位址、寫 START，再用 \code{nanosleep} 輪詢（第 263 行），逾時直接 \code{exit(1)}（第 243、271 行）：

\snip{v100_uio_run_frame}

driver v1.0 把同樣的順序搬進 kernel，等待用 \code{readl_poll_timeout()}，但 bitstream 沒有中斷線（block design 把 PS 的中斷輸入接成常數 0），所以仍是輪詢。代價可以量：v1.1 的 driver 以 \code{use_irq=0} 載入時，輪詢的程式和 driver v1.0 相同（\code{readl_poll_timeout}，每次最多睡 50 µs），在開機組 E 量到 driver 比 PL 多等 \pollWaitUs{} µs，而且這個 process 每幀花掉 \pollCpuUs{} µs 的 CPU 時間（11\%）。
\end{before}

\begin{tech}
\begin{itemize}
\item \textbf{RTL}：新暫存器 IRQ（offset \code{0x0064}）放兩個致能位元；中斷輸出是暫存器輸出、level high。
\item \textbf{IP 封裝與 block design}：\code{ipx::infer_bus_interface} 把 \code{interrupt} 這個 port 宣告成 Vivado 的 interrupt 介面（\code{SENSITIVITY = LEVEL_HIGH}），block design 接到 \code{ps/pl_ps_irq0}；PetaLinux 讀 XSA（Vivado 匯出的硬體描述檔）後自動在節點寫 \code{interrupts = <0 89 4>}。
\item \textbf{kernel}：\code{platform_get_irq_optional()}、\code{devm_request_irq()}、ISR 回傳 \code{irqreturn_t}；\code{struct completion} 的 \code{init_completion()}、\code{reinit_completion()}、\code{complete()}、\code{wait_for_completion_timeout()}；沒有中斷時的 \code{readl_poll_timeout()}；module parameter \code{use_irq}。
\end{itemize}
\end{tech}

\begin{gain}
\begin{itemize}
\item 啟動與等待的整套順序只有一份，在 kernel 裡；逾時回 \code{-ETIMEDOUT} 並在 \code{dmesg} 印出暫存器值，呼叫端可以自己處理，重新載入描述子表就能繼續用。
\item \textbf{CPU 時間}：每幀從 \pollCpuUs{} µs 降到 \irqCpuUs{} µs；\textbf{等待}：driver 比 PL 多的時間從 \pollWaitUs{} µs 降到 \irqWaitUs{} µs（開機組 E，同一次開機、同一個程式，只差 \code{use_irq}）。
\item \textbf{吞吐量}：同一次開機交替跑三輪，中斷 \hIrqFps{} fps、輪詢 \hPollFps{} fps（+0.8\%），三次中斷都比三次輪詢快。
\item 1000 幀剛好 1000 次中斷，沒有多出來的空中斷；device tree 沒有中斷的舊 bitstream，driver 自動退回輪詢。
\end{itemize}
\end{gain}

\subsection{硬體：RTL 加一根線}

每個 clock 都把「致能 AND done」重新算一次，存進暫存器 \code{irq}。\code{ap_done} 是 backbone 做完、\code{hd_done_l} 是 head 做完的黏著位元（設起來後一直保持，直到被清掉），所以 \code{irq} 在 done 設起來的下一個 cycle 拉高，在 done 被清掉的下一個 cycle 放下；\code{ifdef CSR_IRQ_BREAK_MASK} 是故意做壞的版本，只給模擬用：

\snip[Verilog]{csr_irq_reg}

致能放在自己的 IRQ 暫存器，不沿用 CTRL 的 bit 2，因為 CTRL 每一次寫入都會把 bit 2 一起寫掉。driver 啟動寫的是 \code{0x1}、清 done 寫的是 \code{0x2}，bit 2 都是 0，若致能放在 CTRL[2]，每次啟動都會把中斷關掉：

\snip[Verilog]{csr_irq_wr}
\snip[Verilog]{csr_ctrl_wr}

\textbf{gate。}單元 testbench 跑三次：原設計必須 PASS，兩個故意做壞的版本（忽略致能、致能放在 CTRL[2]）必須 FAIL，證明 testbench 抓得到這兩種錯：

\boardfile[：PC 上的 Vivado xsim]{S1_csr_irq_gate.log}{2,4-7,9-15}

head 路徑的整合 gate 用改過的 \code{twn_csr.v} 重跑，4 個 operator、434,176 bytes 逐位元相同：

\boardfile{S2_hd_int_gate.log}{86,626,631,634}

\textbf{從 IP 到 device tree。}封裝 IP 時把 port 宣告成 interrupt 介面，沒有這一步，Vivado 只把它當一條普通的線，PetaLinux 也不會寫出中斷；block design 再把它接到 PS：

\snip[tcl]{pkg_irq}
\snip[tcl]{bd_irq}

PetaLinux 產生的節點（反組譯發佈的 \code{system.dtb}）：\code{interrupts = <0x00 0x59 0x04>} 的 0 表示 SPI（Shared Peripheral Interrupt，所有 CPU 共用的外部中斷），\code{0x59} = 89 是 SPI 編號，4 是 \code{IRQ_TYPE_LEVEL_HIGH}；GIC（ARM 的中斷控制器）內部的編號是 SPI 編號加 32，所以 \code{/proc/interrupts} 顯示 121：

\boardfile[：PC 上反組譯，節點原文]{S3_binding_check.log}{24-34}

改 RTL 就得重建 bitstream，而這個設計在 240 MHz 沒有時序餘裕；第 12 章講怎麼讓它收斂。

\subsection{driver：ISR、completion 與輪詢的退路}

\textbf{probe 時申請中斷。}device tree 有中斷、且 \code{use_irq} 沒被關掉，就用 \code{devm_request_irq()} 註冊 ISR，再寫 IRQ 暫存器打開兩個致能；沒有中斷就留著輪詢：

\snip{drv_probe_irq}

\code{platform_get_irq_optional()} 找不到中斷時不在 kernel log 印錯誤，因為對舊 bitstream 這是正常的。它回傳 \code{-EPROBE_DEFER}（中斷控制器還沒準備好）時，probe 直接回傳，kernel 稍後再呼叫一次 probe，不能當成「沒有中斷」而退回輪詢。之後 \code{misc_register()} 若失敗，probe 先關掉致能再回傳。

\textbf{ISR。}看哪一個 done 設起來，就清掉它、叫醒等它的人；兩個都看完之後讀回一次 IRQ 暫存器，確保清除已經到達 PL 才返回：

\snip{drv_isr}

\textbf{等待。}\code{wait_done()} 有中斷時睡在 completion 上，沒有時走 \code{wait_bit()} 輪詢；兩者都是 3 秒逾時：

\snip{drv_wait_done}
\snip{drv_wait_bit}

\textbf{啟動一幀。}等 busy 變 0（PL 在忙時寫的 START 會被丟掉）；清 done（START 不清 done，殘留的 done 會讓等待立刻結束）；\code{reinit_completion()}；寫兩個位址，再寫 START（\code{writel} 的 barrier 保證 app 寫的輸入比 START 先被 PL 看到，第 2 章）；等 done，清 done，停錶，再讀六個 cycle counter，連同耗時交回 app：

\snip{drv_bb_run}

head 的兩個指令也一樣：\code{twn_hd_start()} 在 START 之前 \code{reinit_completion(&td->hd_done)}，\code{twn_hd_wait()} 呼叫 \code{wait_done()}。啟動與等待拆成兩個指令，app 才能在兩者之間讓 backbone 跑下一幀（第 7 章）。

\subsection{板上}

driver 開機時就用中斷，\code{/proc/interrupts} 有一條 GIC 121、Level：

\boardfile{H3_boot_state.log}{11-15}

\textbf{中斷與輪詢比一比。}同一次開機、同一個 \code{twn_bench}，先用中斷，再用 \code{use_irq=0} 重新載入改成輪詢：

\boardfile[：前一段是中斷，後一段是輪詢]{E3_bench.log}{6-10,14-23}

輪詢時 process 每 13--50 µs 醒來一次，用中斷時一幀只醒一次。剩下的 \irqWaitUs{} µs 裡有 START 前的幾次暫存器存取、中斷延遲、ISR、喚醒與排程、done 之後清 done 的那次寫入，以及 PL 在 block 之間的空檔（\code{cycles_total} 不算空檔）；讀六個 cycle counter 發生在 driver 停錶之後，算在 app 那 \bbSysUs{} µs 裡（第 5 章）。推測空檔佔其中一部分，但沒有計數器可以拆開，沒有量。

\textbf{吞吐量：同一次開機交替量。}每一輪先載入中斷版、再載入輪詢版，各跑一次 120 幀的 pipeline，共三輪：

\boardfile[：交替的中斷與輪詢]{H7_irq_vs_poll_fps.log}{1-5,7,9,11,13,15,17,19,21,23,25}

每一次中斷都比每一次輪詢快，兩組都隨時間略降（233.84→232.71、231.33→231.19）。每一輪都是先中斷後輪詢，這個固定順序不能抵銷漂移；但輪詢三次之間只差 0.14 fps，遠小於兩組約 1.87 fps 的差，結論不受影響。換算成每幀週期，中斷平均 4.290 ms、輪詢 4.324 ms，短了約 35 µs（開機組 H，pipeline）。開機組 E 的 \code{twn_bench}（只跑 backbone、1000 幀）量到 backbone 的等待少了 25.3 µs；兩者是不同的開機與程式，推測其餘來自 head 的等待也改成中斷（沒有分開量）。

\textbf{每一幀剛好一次中斷。}跑 1000 幀前後，ISR 自己的計數和 kernel 的計數（\code{/proc/interrupts}，ISR 回傳 \code{IRQ_NONE} 的呼叫也算在內）都剛好加 1000：

\boardfile{E6_irqcount.log}{4-11}

\textbf{逾時在 3 秒回來。}測試程式先載入一張 PL 會拒絕的描述子表，PL 因此忽略 START，done 永遠不會來；\code{dmesg} 同時印出 driver 的訊息，之後重新載入正確的表，fixture 照樣逐位元相同：

\boardfile{H4_board_check.log}{20-27}

\subsection{原理}

\textbf{輪詢的退路怎麼睡。}\code{readl_poll_timeout(addr, val, cond, 50, 3000000)} 反覆「讀暫存器、檢查條件、\code{usleep_range(13, 50)}」，直到成立或逾時；逾時後再讀一次，避免「睡過頭的那一刻其實已經完成」被誤判。不用 \code{msleep(1)}：這個 kernel 的 \code{HZ=250}，\code{msleep} 至少睡一個 4 ms 的 jiffy，等於白多等一幀；\code{usleep_range} 用 hrtimer（高解析度計時器）。也不 busy-wait：kernel 是 \code{PREEMPT_NONE}，系統呼叫裡空轉 4 ms 會佔住一顆核心。\code{usleep_range} 的上限是硬到期時間，CPU 閒著時多半睡滿 50 µs，一幀大約醒 80 次（估計，沒有量），每次要排程進出、設定下一個 hrtimer、讀一次暫存器，這些算在 process 的 CPU 時間裡，累積起來每幀約 0.45 ms；計時器中斷本身多半在 CPU 閒置時發生，算在 idle（下一段）。

\textbf{CPU 時間怎麼量。}\code{twn_bench} 用 \code{getrusage(RUSAGE_SELF)}，只算這個 process 的 user 與 kernel 時間。這顆 kernel 沒開 \code{IRQ_TIME_ACCOUNTING}，CPU 閒置時進來的中斷（ISR 與 hrtimer）算在 idle，所以 \irqCpuUs{} µs 不含 ISR 本身的成本。

\textbf{為什麼用 level，ISR 為什麼一定要清 done。}level 中斷是「線拉高的期間一直成立」。GIC 在 ISR 結束（EOI）後重新看這條線，done 沒清，線還是高，ISR 會立刻再被呼叫，形成中斷風暴（interrupt storm）。level 的好處是不會漏：兩個 done 幾乎同時設起來，線會一直高到兩個都被處理。

\textbf{為什麼清完還要讀一次。}\code{writel()} 是 posted write：CPU 送出就繼續往下走，不等寫入真的到達 PL。ISR 若馬上返回，GIC 可能還看到舊的高電位，多進一次 ISR。讀同一個裝置的暫存器會等前面的寫入完成，讀回來時，清除的寫入已經到了 PL。這是「完成」的問題，barrier 只管「順序」，解決不了（第 2 章）。

\textbf{\code{reinit_completion()} 為什麼放在 START 之前。}completion 是一個計數加一個等待佇列：\code{complete()} 把計數加一並叫醒等待者，\code{wait_for_completion_timeout()} 在計數大於 0 時減一並返回。若歸零放在 START 之後，PL 很快做完時，ISR 的 \code{complete()} 可能先發生、接著被歸零，等待的人就睡到 3 秒逾時。放在 START 之前，PL 還沒開始，不可能有中斷搶先。

\textbf{ISR 為什麼可以在硬體中斷裡做完。}它只做幾次暫存器讀寫和 \code{complete()}，都不會睡，可以在 interrupt context（CPU 被中斷打斷、還沒回到任何 process 時執行的程式，不能睡眠）執行，不需要 threaded IRQ（把處理函式放到 kernel 執行緒、可以睡眠的處理方式）。

\textbf{初始化、申請、打開致能的順序。}probe 先 \code{init_completion()}，再 \code{devm_request_irq()}，最後打開致能；remove 先關掉致能，再由 \code{devm} 釋放中斷。對這條沒有共用的線，順序反過來其實也不會出現「送到 CPU 卻沒人處理」的中斷：GIC 端的線在 \code{request_irq()} 裡才啟動，最後一個 handler 釋放時就關掉。這個順序真正保護的是兩件事：ISR 一被呼叫就會用到 completion，所以它必須先初始化；若中斷線和別的裝置共用，沒人清的 level 中斷會被 kernel 判成「nobody cared」而把整條線關掉。

\subsection{面試會問}

\begin{qa}{你怎麼把輪詢改成中斷？從硬體講到軟體。}
硬體上，我在控制暫存器模組加了一個暫存器輸出的 level-high 中斷，等於「致能 AND done」；致能放在獨立的暫存器，因為 CTRL 每次寫入都會覆蓋原本的 bit 2，driver 一啟動就會把中斷關掉，這個錯誤我在 testbench 裡故意做出來、確認抓得到。IP 封裝時宣告成 interrupt 介面，block design 接到 \code{pl_ps_irq0}，PetaLinux 就在 device tree 寫出 SPI 89、level high。driver 在 probe 用 \code{devm_request_irq()} 註冊 ISR；ISR 清 done、\code{complete()}，返回前讀回一次確保清除已經到達；等待的一邊用 \code{wait_for_completion_timeout()}，\code{reinit_completion()} 放在 START 之前避免漏掉中斷。在同一次開機上只切換 \code{use_irq} 比較：每幀 CPU 時間從 \pollCpuUs{} µs 降到 \irqCpuUs{} µs；另一次開機交替量三輪，pipeline 從 \hPollFps{} 變 \hIrqFps{} fps。device tree 沒有中斷時自動退回 \code{readl_poll_timeout} 輪詢。
\end{qa}

\subsection{關鍵函式速查}

\begin{center}
\small
\begin{tabularx}{\textwidth}{l l X}
\toprule
名稱 & 位置 & 做什麼 \\
\midrule
\code{irq} 暫存器 & \code{rtl/fuse/twn_csr.v}:\lineof{csr_irq_assign} & 致能 AND done，level high \\
interrupt 介面 & \code{vivado/package_fuse_ip.tcl}:\lineof{pkg_irq} & \code{SENSITIVITY = LEVEL_HIGH} \\
接到 PS & \code{vivado/build_bd_fuse.tcl}:\lineof{bd_irq} & \code{pl_ps_irq0} \\
\code{use_irq} & \code{twn_fuse.c}:\lineof{drv_param_irq} & 可以關掉中斷，退回輪詢 \\
probe 的中斷段 & \code{twn_fuse.c}:\lineof{drv_probe_irq} & 申請中斷、打開致能 \\
\code{twn_isr()} & \code{twn_fuse.c}:\lineof{drv_isr} & 清 done、\code{complete()}，返回前讀回 \\
\code{wait_done()} & \code{twn_fuse.c}:\lineof{drv_wait_done} & completion 或輪詢，3 秒逾時 \\
\code{wait_bit()} & \code{twn_fuse.c}:\lineof{drv_wait_bit} & \code{readl_poll_timeout}，50 µs、3 s \\
\code{twn_bb_run()} & \code{twn_fuse.c}:\lineof{drv_bb_run} & 等閒置、清 done、寫位址、START、等 done \\
\bottomrule
\end{tabularx}
\end{center}

\clearpage
\section{並行控制：兩把 mutex 與獨占開啟}
\label{sec:lock}

\begin{oneline}
同一時間只有一次 \code{open()} 會成功，而在開著的這一份裡，backbone 和 head 各用一把 mutex，同一個引擎的暫存器操作不會被兩個 thread 交錯執行。
\end{oneline}

\begin{before}
v1.0.0 完全沒有鎖，也沒有獨占：

\snip{uio_map_regs}

第 130 行的 \code{open(uio, O_RDWR | O_SYNC)} 誰都能開、開幾次都可以。兩個 app 一起跑，就會互相改寫對方的描述子表和位址；pipeline 該怎麼安排兩個引擎，也完全靠 app 自己遵守順序。
\end{before}

\begin{tech}
\begin{itemize}
\item \code{struct mutex}、\code{mutex_lock_interruptible()}：可以睡眠的鎖；等鎖時被訊號（例如 Ctrl-C）打斷就回傳 \code{-EINTR}。
\item \code{test_and_set_bit()} / \code{clear_bit()}：不可分割（atomic）的位元操作，用來實作「只能開一次」。
\item \code{file_operations.open} / \code{.release}：開檔、最後一次關檔時被呼叫。
\item \code{container_of()}：已知「某個欄位的位址」，算出「包含它的 struct 的位址」。
\item \code{.owner = THIS_MODULE}：檔案開著時，module 不能卸載。
\item \code{suppress_bind_attrs}：不提供 sysfs 的 \code{bind} / \code{unbind}。
\end{itemize}
\end{tech}

\begin{gain}
\begin{itemize}
\item 第二次 \code{open()} 會得到 \code{-EBUSY}；裝置開著時、以及關掉之後只剩映射時，\code{rmmod} 都被拒絕（兩者都在板上測過，見下）。
\item 同一個引擎的一整段暫存器操作（例如 \code{BB_RUN} 的「等閒置、清 done、寫位址、START、等 done」）不會被另一個 thread 插進來。
\item 兩個引擎各一把鎖：兩個 thread 分別驅動兩個引擎時互不阻擋。這是設計上的保證；這個 app 由同一個 thread 依序驅動兩者，量不出差別（見下方「踩過的坑」）。
\end{itemize}
\end{gain}

\subsection{實際例子：程式碼與板上輸出}

\textbf{獨占開啟。}\code{in_use} 的第 0 位元就是「有人開著」。\code{test_and_set_bit} 一次完成「讀舊值、設成 1」，兩個 process 同時開也只有一個會讀到 0：

\snip{drv_open}

\textbf{從 \code{struct file} 找回 driver 的資料。}\code{misc_open()}（Linux 6.1 \code{drivers/char/misc.c} 第 138 行）把 \code{struct miscdevice} 的位址存進 \code{file->private_data}。這個 \code{miscdevice} 是 \code{struct twn_dev} 裡的一個欄位，\code{container_of} 依欄位在 struct 裡的位置往回推，就得到整個 \code{struct twn_dev}：

\snip{drv_file_to_td}
\snip{drv_struct}

\textbf{分派與上鎖。}\code{twn_ioctl} 先決定這個指令屬於哪個引擎、要拿哪把鎖，再呼叫對應的函式。\code{STATUS} 只讀暫存器，不拿鎖，所以 backbone 跑到一半時也能查狀態：

\snip{drv_ioctl}

\textbf{app 怎麼讓兩個引擎重疊。}pipeline 的 producer thread 每一輪依序做三件事：啟動 head（處理上一幀）、跑 backbone（這一幀）、等 head（第 517、529、532 行）：

\snip{app_frame_pl}

\begin{center}
\begin{tikzpicture}[font=\footnotesize,x=10.5mm,y=7mm]
\draw[-{Stealth[length=2mm]}] (0,0) -- (12,0) node[right]{時間};
\node[left] at (0,2) {head（PL）};
\node[left] at (0,1) {backbone（PL）};
\fill[qabg,draw=qac] (0.3,1.7) rectangle (4.3,2.3); \node at (2.3,2) {head：第 $k-1$ 幀};
\fill[techbg,draw=techc] (0.5,0.7) rectangle (4.9,1.3); \node at (2.7,1) {backbone：第 $k$ 幀};
\fill[qabg,draw=qac] (5.3,1.7) rectangle (9.3,2.3); \node at (7.3,2) {head：第 $k$ 幀};
\fill[techbg,draw=techc] (5.5,0.7) rectangle (9.9,1.3); \node at (7.7,1) {backbone：第 $k+1$ 幀};
\draw[dashed] (0.3,0) -- (0.3,2.6) node[above,font=\scriptsize]{HD\_START};
\draw[dashed] (0.5,0) -- (0.5,-0.4) node[below,font=\scriptsize]{BB\_RUN 開始};
\draw[dashed] (4.9,0) -- (4.9,2.6) node[above,font=\scriptsize]{HD\_WAIT};
\end{tikzpicture}
\end{center}

兩個引擎在硬體上同時跑；CPU 這邊則是同一個 thread 依序發出三個 ioctl。

\textbf{板上：重複開啟與卸載都被擋下。}

\boardfile{A2_board_check_uio_dtb.log}{10-10}
\boardfile{A2_board_check_uio_dtb.log}{31-38}

第二個框的第 3 行，是程式先 \code{mmap}、再 \code{close} 掉 fd、只留下映射時試 \code{rmmod}（\code{twn_errtest hold}）。

\begin{pit}
\textbf{第一個：sysfs 的 unbind 會繞過所有保護。}一開始 \code{bind.sh} 用 \code{/sys/bus/platform/drivers/twn_fuse/unbind} 把裝置交還給 UIO。後來發現這條路很危險：unbind 會直接呼叫 \code{remove()}，\code{devm_} 的 buffer 隨即被釋放，就算還有 app 把那塊 buffer 映射著也一樣，app 之後寫進去的資料會落在已經還給系統的記憶體上。這是讀程式時發現的風險，沒有實際發生過。

修法是在 driver 加上 \code{.suppress_bind_attrs = true}（第 1 節 \code{twn_driver}），拿掉 \code{bind} / \code{unbind} 這兩個檔案。裝置要離開 driver 只剩 \code{rmmod} 一條路，而 \code{rmmod} 在檔案還開著、或 buffer 還映射著時會被 kernel 拒絕。\code{bind.sh} 也改成用 \code{rmmod} 加 \code{drivers_probe} 切換。

\textbf{第二個：我先前說錯的一件事。}規劃時我說「整個裝置只用一把鎖的話，重疊會消失，fps 會掉」。這對這個 app 不成立：\code{HD_START} 一下就返回，鎖跟著釋放，硬體已經在跑；接著 \code{BB_RUN} 拿鎖、等 backbone；最後 \code{HD_WAIT}。三個呼叫來自同一個 thread、依序執行，所以只用一把鎖，硬體照樣重疊。兩把鎖真正的作用，是當兩個引擎由兩個 thread 分別驅動時（例如另開一個 thread 專門等 head），等其中一個引擎不會擋住另一個。這是設計上的正確性，不是這個 app 量得到的效能差別。
\end{pit}

\subsection{原理}

\textbf{為什麼用 mutex，不用 spinlock。}拿著鎖的時候，driver 會在 \code{readl_poll_timeout} 裡睡眠，一等就是好幾 ms。spinlock 拿著時不能睡（等鎖的 CPU 會一直空轉），所以只能用 mutex：等不到鎖的 thread 會被排程器換下去，把 CPU 讓給別人。

\textbf{為什麼用 \code{_interruptible}，又為什麼回 \code{-ERESTARTSYS}。}一般的 \code{mutex_lock} 等鎖時連 \code{kill} 都叫不醒。\code{mutex_lock_interruptible} 被訊號打斷時回傳 \code{-EINTR}（Linux 6.1 \code{kernel/locking/mutex.c} 第 667 行）。driver 收到之後改回傳 \code{-ERESTARTSYS}，這是 kernel 內部的特殊值：系統呼叫返回 user space 前，kernel 依訊號處理函式有沒有設 \code{SA_RESTART}，決定要自動重做這個 ioctl，還是讓 app 收到 \code{EINTR}。

\textbf{為什麼一次只開一次。}描述子表、base 位址、BB2HD 模式都是整顆加速器共用的狀態。兩個 process 同時用，各自載入自己的表，就會互相蓋掉。讓整個裝置同一時間只有一次 \code{open()} 成功，比替每個開檔維護一份狀態簡單得多，也符合「一顆加速器同一時間只跑一個網路」的實際用法。精確地說，它擋的是第二次 \code{open()}：同一個 process \code{fork()} 出來的子程序、\code{dup()} 出來的 fd，共用同一個開啟的檔案，不會再經過 \code{twn_open()}。

\textbf{rmmod 為什麼會被拒絕。}\code{misc_open} 用 \code{fops_get} 取得 driver 的 \code{file_operations}，因為 \code{.owner = THIS_MODULE}，這一步會把 module 的參考計數加一，檔案完全關閉時才減回去。\code{mmap} 建立的 VMA 也持有這個檔案。所以只要有人開著裝置或映射著 buffer，module 的參考計數就不是 0，\code{rmmod} 回報 \code{Module twn_fuse is in use}（上方板上輸出）。

\subsection{面試會問}

\begin{qa}{你的 driver 怎麼處理並行？為什麼用 mutex 不用 spinlock？}
兩層。第一層是獨占開啟：\code{open} 用 \code{test_and_set_bit}，同一時間只有一次 open 會成功，因為描述子表和 base 位址是整顆加速器共用的狀態。第二層是 thread：backbone 和 head 各一把 mutex，同一個引擎的一整段暫存器操作不會被另一個 thread 插進來，兩個引擎之間不互相阻擋；只讀的 \code{STATUS} 不拿鎖。用 mutex 是因為拿著鎖時我會在 \code{readl_poll_timeout} 裡睡好幾 ms，spinlock 不能睡。我也拿掉了 sysfs 的 unbind，因為它會在 buffer 還被映射時釋放 buffer；這樣卸載只剩 \code{rmmod}，板上測過裝置開著、或只剩映射時，\code{rmmod} 都被拒絕。
\end{qa}

\subsection{關鍵函式速查}

\begin{center}
\small
\begin{tabularx}{\textwidth}{l l X}
\toprule
名稱 & 位置 & 做什麼 \\
\midrule
\code{struct twn_dev} & \code{twn_fuse.c}:\lineof{drv_struct} & 兩把鎖、\code{in_use}、內嵌的 \code{miscdevice} \\
\code{file_to_td()} & \code{twn_fuse.c}:\lineof{drv_file_to_td} & \code{container_of} 找回 \code{struct twn_dev} \\
\code{twn_open()} & \code{twn_fuse.c}:\lineof{drv_open} & 獨占開啟，第二次回 \code{-EBUSY} \\
\code{twn_ioctl()} & \code{twn_fuse.c}:\lineof{drv_ioctl} & 選鎖、上鎖、分派、解鎖 \\
\code{mbv2_fuse_frame_pl()} & \code{mbv2_fuse_ps.c}:\lineof{app_frame_pl} & head start → backbone run → head wait \\
\bottomrule
\end{tabularx}
\end{center}

\needspace{20\baselineskip}
\section{看得見硬體：sysfs 與 debugfs}
\label{sec:observe}

\begin{oneline}
給人和腳本看的狀態（六個 cycle counter、STATUS、NBLK）放 sysfs，一個檔案一個值、\code{cat} 就能看；只給開發者的除錯工具（暫存器傾印、短幀、中斷計數、SMMU 測試）放 debugfs，不讓 app 依賴。
\end{oneline}

\begin{before}[原本怎麼做（v1.0.0 的 UIO 版本，與 driver v1.0）]
v1.0.0 的 counter 只能由 app 在跑完一幀後經 \code{/dev/uio4} 讀；沒有 app 在跑時，想看硬體狀態得另外寫程式或用 \code{devmem}。UIO 版本的除錯功能（只跑前 n 個 block 的「短幀」、像素相位的觀察位元、陣列的第一個輸出）都是 app 直接寫暫存器。driver v1.0 不讓 app 碰暫存器，這三個函式就失效了：

\snip{v10_hal_stubs}

\code{fuse_set_nblk()} 至少印訊息、回傳 $-1$；另外兩個安靜地回傳 0，呼叫的人分不出「硬體讀到 0」和「功能不存在」，比直接失敗更容易誤導。
\end{before}

\begin{tech}
\begin{itemize}
\item sysfs：\code{DEVICE_ATTR_RO(name)} 宣告唯讀屬性，讀檔時呼叫 \code{name_show()}；\code{sysfs_emit()} 保證輸出不超過一頁；\code{ATTRIBUTE_GROUPS()} 加 \code{driver.dev_groups}，讓 driver core 在綁定後自動建檔、解除前移除。
\item debugfs：掛在 \code{/sys/kernel/debug}，只有 root 能讀寫，格式不保證穩定。\code{debugfs_create_regset32()} 一次印出一張暫存器表；\code{DEFINE_DEBUGFS_ATTRIBUTE()} 加 \code{debugfs_create_file_unsafe()} 做可寫的數字檔；\code{debugfs_create_u32()} 把一個變數開成唯讀檔；\code{debugfs_remove_recursive()} 整個目錄一起移除。
\end{itemize}
\end{tech}

\begin{gain}
\begin{itemize}
\item driver 綁定期間，任何時候都能 \code{cat} 硬體狀態，不用開裝置；檔案的建立與移除交給 driver core，和 probe／remove 沒有先後的競爭。
\item 短幀回來了，而且比 UIO 安全：寫 \code{nblk} 檢查 1 到 17（寫 99 被拒絕），並拿 backbone 的鎖，不會在一幀進行到一半時改掉 NBLK。
\item 多了兩個工具：\code{irq_count}（ISR 處理了幾次，第 6 章用它驗證沒有空中斷）與 \code{dma_fault}（故意越界讀取，第 11 章測 SMMU）。
\end{itemize}
\end{gain}

\subsection{實際例子}

\textbf{sysfs：一個巨集產生一個屬性。}\code{TWN_REG_ATTR(name, off, fmt)} 展開成一個 \code{name_show()}（讀一個暫存器、照格式印出）加上 \code{DEVICE_ATTR_RO(name)}；用它宣告八個屬性，列成表交給第 1 章 \code{twn_driver} 的 \code{.dev_groups}：

\snip{drv_reg_attr}
\snip{drv_reg_attr_use}
\snip{drv_attrs}

\textbf{板上：跑完一幀後讀出來。}

\boardfile{H10_driver_state.log}{25-34}

PL 時脈 240 MHz。\code{cycles_total} 是 17 個 block 執行期間的 cycle 總和（block 之間的空檔不算），\cyclesTotal{} cycles 約 4.13 ms，是 START 到 done 的下限；\code{cycles_array} 佔 98.5\%，三值運算陣列幾乎一直在工作；\code{cycles_stall} 只有 167，記憶體幾乎沒有讓運算停下來。\code{status} 的欄位：

\begin{center}
\small
\begin{tabularx}{\textwidth}{l l X}
\toprule
位元 & 名稱 & 這次的值與意思 \\
\midrule
0、1 & \code{d_valid}、\code{tab_checked} & 1、1：描述子表通過檢查、檢查已經跑過 \\
6:2 & \code{cur_block} & 16：最後執行的 block（0 到 16） \\
10:7、15:11 & \code{err_code}、\code{err_block} & 0、0：表格沒有錯 \\
16 & \code{ovr_sticky} & 1：見下段 \\
\bottomrule
\end{tabularx}
\end{center}

STATUS bit 16 是 ring overrun 的黏著旗標：PL 裡某個列緩衝區在滿的時候被寫入，就會設成 1（只記錄、不停機），只有 PL reset 會清掉。它是 bitstream 本身在正常執行時就會設起來的，和 driver 無關：v1.0.0 時期回退到原本的 dtb、driver 完全沒載入，只用 UIO 版本跑 4 張 fixture，跑之前是 0、跑完變成 \code{0x00010043}，輸出照樣逐位元正確：

\boardfile{C1_rollback_to_uio_dtb.txt}{10-21}

原因還沒查明，記錄在 \code{results/board_results.json}，留待之後追查。

\textbf{debugfs：建立目錄與檔案。}probe 的最後一步呼叫 \code{twn_debugfs_init()}；\code{regs} 背後是一張「名字、offset」的表，可以寫的 \code{nblk} 在寫入時檢查範圍、拿 backbone 的鎖：

\snip{drv_debugfs_init}
\snip{drv_nblk_set}

\textbf{app 這一側。}HAL 用 \code{fopen}／\code{fgets} 讀 \code{regs}、找出某個名字的那一行；設定短幀就是寫 \code{nblk} 再讀回確認：

\snip{hal_dbg_reg}
\snip{hal_set_nblk}

\textbf{板上。}\code{IRQ = 0x00000003} 表示兩個中斷致能都打開了：

\boardfile{E4_debugfs.log}{4-25}

接著用短幀跑一次，再試一個不合法的值：

\boardfile{E4_debugfs.log}{26-36}

\code{PIX_NZ} 與 \code{CAP0}--\code{CAP3} 都是 0：這兩組觀察邏輯只在 RTL 用 \code{FUSE_OBS} 編譯時才會建出來，發佈的 bitstream 沒有，暫存器固定讀到 0。debugfs 只是讓它們讀得到。

\subsection{原理}

\textbf{這些檔案的 read／write 是誰寫的。}app 讀 \code{regs}、寫 \code{nblk}，底層確實是 \code{read()}／\code{write()} 系統呼叫，但那組 file operations 是 kernel 產生的：\code{regs} 由 \code{debugfs_create_regset32()} 處理，只用到 driver 給的名字與 offset 表；\code{nblk} 由 \code{DEFINE_DEBUGFS_ATTRIBUTE()} 展開，寫入時才呼叫 driver 的 \code{twn_nblk_set()}；sysfs 的屬性檔讀取時呼叫 driver 的 \code{show()}。「文字轉數字、處理 offset 與長度、\code{copy_to_user}」這些瑣事都由 kernel 做，driver 只寫取值與設值。這也是 \code{/dev/twn_fuse} 不需要自己的 \code{.read}／\code{.write} 的原因之一（第 5 章）。

\textbf{為什麼用 \code{.dev_groups}，不在 \code{probe()} 裡自己建檔。}driver core 在 \code{probe()} 成功之後才建立 \code{dev_groups} 的檔案，在呼叫 \code{remove()} 之前先移除，所以 \code{show()} 被呼叫時，\code{dev_get_drvdata()} 一定拿得到 probe 設好的資料。

\textbf{為什麼除錯工具放 debugfs。}各種介面怎麼選在第 5 章。sysfs 的檔案一旦給人用，改格式就是改了 ABI（給使用者的介面約定）；除錯工具不希望 app 依賴，所以放 debugfs：kernel 文件說它「不打算當成給 user space 的穩定 ABI」（\code{Documentation/filesystems/debugfs.rst}），同時也提醒最好還是當成要長期維護的介面來設計。

\textbf{為什麼用 \code{_unsafe} 版本。}\code{debugfs_create_file()} 每次讀寫都多包一層保護，防止檔案被移除後還被呼叫；\code{DEFINE_DEBUGFS_ATTRIBUTE()} 產生的 file operations 自己已經處理了這件事，所以搭配 \code{debugfs_create_file_unsafe()} 避免重複保護（\code{fs/debugfs/inode.c} 的註解）。

\textbf{為什麼寫 \code{nblk} 要拿鎖。}NBLK 決定 PL 這一幀要跑幾個 block，一幀進行中被改，PL 的 sequencer 可能在中途看到不同的值。\code{TWN_IOC_BB_RUN} 從頭到尾都拿著 \code{bb_lock}，所以只要沒有發生過逾時，拿到鎖就表示沒有幀在跑；逾時的 ioctl 會釋放鎖，但 RTL 沒有中止功能，PL 若其實還在跑，這個保證就不成立。

\subsection{面試會問}

\begin{qa}{除了 ioctl，你怎麼讓人看到硬體狀態、怎麼提供除錯功能？}
狀態用 sysfs：六個 cycle counter、STATUS、NBLK 各是一個唯讀檔，一個巨集產生 \code{show} 加 \code{DEVICE_ATTR_RO}，再用 \code{ATTRIBUTE_GROUPS} 交給 platform driver 的 \code{.dev_groups}，driver core 在 probe 成功後建檔、remove 前移除。板上跑完一幀，\code{cat cycles_total} 是約 99 萬個 cycle、4.13 ms，運算陣列忙碌 98.5\%。除錯工具放 debugfs，不讓 app 依賴：\code{debugfs_create_regset32()} 一次印出暫存器；可寫的 \code{nblk} 用 \code{DEFINE_DEBUGFS_ATTRIBUTE()}，寫入時檢查範圍、拿 backbone 的 mutex，所以比 UIO 時代直接寫暫存器安全。這些檔案的 read／write 都由 kernel 提供，我只寫取值與設值的函式。
\end{qa}

\subsection{關鍵函式速查}

\begin{center}
\small
\begin{tabularx}{\textwidth}{l l X}
\toprule
名稱 & 位置 & 做什麼 \\
\midrule
\code{TWN_REG_ATTR} & \code{twn_fuse.c}:\lineof{drv_reg_attr} & 一個巨集產生 \code{show()} 與 \code{DEVICE_ATTR_RO} \\
\code{twn_attrs[]} & \code{twn_fuse.c}:\lineof{drv_attrs} & 八個 sysfs 屬性與 \code{ATTRIBUTE_GROUPS(twn)} \\
\code{twn_debugfs_init()} & \code{twn_fuse.c}:\lineof{drv_debugfs_init} & 建 debugfs 目錄與四個檔案 \\
\code{twn_nblk_set()} & \code{twn_fuse.c}:\lineof{drv_nblk_set} & 範圍檢查、拿鎖、寫 NBLK \\
\code{dbg_reg()} & \code{pl_hal_fuse_drv.c}:\lineof{hal_dbg_reg} & app 從 \code{regs} 讀一個暫存器 \\
\code{fuse_set_nblk()} & \code{pl_hal_fuse_drv.c}:\lineof{hal_set_nblk} & app 設定短幀 \\
\bottomrule
\end{tabularx}
\end{center}

\needspace{20\baselineskip}
\section{App 端：同一份 HAL 介面，以及影像怎麼進 buffer}
\label{sec:app}

\begin{oneline}
app 的演算法與位址計算沒有改：\code{pl_hal_fuse.h} 有兩個實作，編譯時用 \code{HAL=drv}（預設）或 \code{HAL=uio} 選一個；driver 版本回報「buffer 基底位址是 0」，app 算出的位址就自動變成 offset。
\end{oneline}

\begin{before}
v1.0.0 的 \code{Makefile.cross} 只有一個實作可選（\code{pl_hal_fuse.c}）：

\snip[make]{v100_make_common}

\code{pl_hal_fuse.c} 裡混著兩種程式：要碰暫存器的（UIO、\code{/dev/mem}），和完全不碰硬體的（head 描述子的編碼與檢查、印 counter）。app（\code{mbv2_fuse_ps.c}）只呼叫這份介面、不直接碰硬體，這個分層正是能換實作的原因。
\end{before}

\begin{tech}
\begin{itemize}
\item HAL（Hardware Abstraction Layer）：上層只依賴一份 header，底下可以換不同的實作；\code{Makefile.cross} 用 \code{HAL} 變數決定編哪一個 \code{.c}。
\item 重構：不碰硬體的函式搬到 \code{pl_hal_fuse_enc.c}，兩個實作共用，內容一字不改。
\item 「基底 0」：driver 版本的 \code{fuse_arena_init()} 回報實體位址 0。
\item \code{_Static_assert}：編譯時期的檢查，buffer 配置超過 6 MB 就編譯失敗。
\end{itemize}
\end{tech}

\begin{gain}
\begin{itemize}
\item 兩個版本跑的是同一份 app，量到的差異只來自 HAL 底下那一層，比較才有意義（第 3、13 章）。
\item UIO 版本保留下來，給 v1.0.0 的開機檔與 A/B 實驗用；v1.1 的映像只用 driver 版本。
\item app 本身只改了一處：要求的 buffer 從整塊 256 MB 改成實際需要的 6 MB。
\end{itemize}
\end{gain}

\subsection{實際例子}

\textbf{在 Makefile 選實作。}\code{HAL} 是 \code{drv} 就編 \code{pl_hal_fuse_drv.c}，否則編 \code{pl_hal_fuse.c}；輸出放在不同資料夾：

\snip[make]{app_make_hal}

\textbf{每個 HAL 函式都變成一個 ioctl。}共用的錯誤處理，與跑一幀的函式：

\snip{app_xioctl}
\snip{app_run_frame}

和 UIO 版本 61 行的 \code{fuse_run_frame()}（第 6 章）比，這裡只剩填參數、呼叫、存下 counter。注意其中的 \code{exit(1)}：\code{fuse_run_frame()} 的介面沒有錯誤回傳值，為了不改介面，driver 回的錯誤碼在這裡仍然變成結束程式；要讓 app 自己處理逾時，得連 \code{pl_hal_fuse.h} 一起改。

\textbf{「基底 0」怎麼讓 app 不用改。}app 算位址的方式一直是「基底加區段 offset」：

\snip{app_run_call}

UIO 版本的 \code{g_arena_pa} 是 \code{0x6000_0000}，傳下去的是實體位址，HAL 直接寫進暫存器；driver 版本的 \code{g_arena_pa} 是 0（第 4 章），傳下去的就是 offset，HAL 原封不動交給 ioctl，driver 再加上真正的 DMA 位址。同一行 app 程式，在兩個實作底下都對。

\textbf{板上：兩個版本，同樣的答案。}v1.0.0 的開機檔上，同一份 app 分別連結兩個 HAL，四張合成 fixture 都逐位元相同：

\boardfile{D4_synthetic.log}{1,7-8,12,19-20,24}

\subsection{影像怎麼進 buffer：以攝影機 app 為例}

\code{sw/app/main_cam.c} 從 USB 攝影機讀影像，經過五步才變成 PL 能吃的格式：
\begin{enumerate}
\item \textbf{從攝影機拿一幀}：攝影機由 kernel 的 uvcvideo driver 管，app 用 V4L2（Video4Linux2）的 ioctl 操作：\code{VIDIOC_S_FMT} 設 640×480 與格式（YUYV 或 MJPEG），\code{VIDIOC_REQBUFS} 請 kernel 配 capture buffer 並 \code{mmap} 進來，每幀 \code{VIDIOC_DQBUF} 取一格、用完 \code{VIDIOC_QBUF} 還回去。用法和 \code{/dev/twn_fuse} 一樣是「ioctl 下指令、mmap 給資料」；uvcvideo 的 \code{.read} 直接回 \code{-EINVAL}。
\item \textbf{解碼成 RGB}：YUYV 用固定小數點換算；MJPEG 用 libturbojpeg。
\item \textbf{裁切與縮放}：從畫面中間取正方形、略往下移，nearest 取樣成 256×256。
\item \textbf{量化成 int8}：訓練時的輸入是「像素 ÷ 255、減 ImageNet 平均、除標準差、除模型的輸入 scale、四捨五入」。這串浮點運算事先算成每個 channel 256 格的表（\code{tools/gen_input_qlut.py} 從模型的 \code{model_meta.c} 產生 \code{input_qlut.h}），執行時查表：
\snip{cam_prep_input}
\item \textbf{換排列、寫進 buffer}：上一步得到的是 CHW（三張平面），PL 要的是 HWC（每個像素的 R、G、B 連在一起），每列補齊到 128 bytes 的倍數，直接寫進 mmap 來的 buffer（第 3 章的 \code{OFF_IN}）：
\snip{app_write_input}
\end{enumerate}

接著 \code{ioctl(TWN_IOC_BB_RUN)} 只傳這塊輸入的 offset；driver 寫 START 時，\code{writel} 的 barrier 保證剛寫的影像比 START 先被 PL 看到（第 2 章）。推論結束後，app 讀回 head 的 heatmap，找出 14 個關節點，畫在 \code{/dev/fb0}（framebuffer，又是一個 mmap 出來的記憶體）上。

PL 吃的是量化後的 256×256 int8，攝影機給的是 640×480 的 YUYV 或 JPEG，每個像素都要 CPU 處理一次，所以讓 PL 直接讀攝影機的 buffer（dma-buf 共用）在這裡沒有意義。查表這一步也沒有 fixture 驗證，因為 fixture 存的是量化好的 tensor；app 留了 \code{CAM_QSTAT=1} 印出量化後的範圍供人工檢查。攝影機 app 的 \code{--selftest}（把 fixture 走完同樣的推論並逐位元比對）只在 v1.0.0 的 UIO 版本上做過（\code{results/board_results.json} 的 \code{camera_app_selftest}）；driver 版本的攝影機 app 有編出來，但 v1.1 的板上檢查沒有跑它。

\subsection{原理}

\textbf{為什麼在連結時選，不在執行時選。}執行時選（例如函式指標表）要多一層間接呼叫，兩個實作也都會編進同一個執行檔。這裡兩個版本不會在同一個程式裡同時用到，用 Makefile 在連結時選最簡單：每個執行檔只含一個實作，函式名稱完全相同，app 看不出差別。

\textbf{為什麼不乾脆讓 app 直接呼叫 ioctl。}那得改動 app 裡每一處算位址、每一次啟動與等待，UIO 版本也保留不了。守住 HAL 這一層，差異都關在兩個 \code{.c} 裡，app 不用改，也就不會因為改 app 而引入新的錯誤。不碰硬體的編碼函式搬到共用檔、內容不改，兩個版本就不會悄悄產生不同的結果。

\subsection{面試會問}

\begin{qa}{換成 driver 之後，上層的 app 改了多少？攝影機的影像怎麼送進去？}
演算法和位址計算都沒改。app 本來就只透過 \code{pl_hal_fuse.h} 碰硬體，我替它寫了第二個實作，用 Makefile 的 \code{HAL=drv} 在連結時選。關鍵是 driver 版本回報 buffer 基底位址為 0：app 用「基底加 offset」算位址，基底是 0 時算出來的就是 offset，正好是 ioctl 要的。攝影機那邊：用 V4L2 的 ioctl 和 mmap 拿到一幀，解碼成 RGB、裁切成 256×256、每個 channel 查表量化成 int8，再轉成 PL 要的 HWC 排列直接寫進 mmap 來的 DMA buffer，最後一個 ioctl 只傳 offset。driver 寫 START 時 \code{writel} 的 barrier 保證影像比 START 先被 PL 看到。
\end{qa}

\subsection{關鍵函式速查}

\begin{center}
\small
\begin{tabularx}{\textwidth}{l l X}
\toprule
名稱 & 位置 & 做什麼 \\
\midrule
\code{HAL} 變數 & \code{Makefile.cross}:\lineof{app_make_hal} & 選 \code{pl_hal_fuse_drv.c} 或 \code{pl_hal_fuse.c} \\
\code{xioctl()} & \code{pl_hal_fuse_drv.c}:\lineof{app_xioctl} & 呼叫 ioctl，失敗時印出是哪一個 \\
\code{fuse_run_frame()} & \code{pl_hal_fuse_drv.c}:\lineof{app_run_frame} & 一個 \code{TWN_IOC_BB_RUN} \\
\code{prep_input()} & \code{main_cam.c}:\lineof{cam_prep_input} & 裁切、縮放、查表量化 \\
\code{mbv2_fuse_write_input()} & \code{mbv2_fuse_ps.c}:\lineof{app_write_input} & CHW 轉 HWC，寫進 buffer \\
\code{pl_hal_fuse_enc.c} & \code{sw/app/} & 兩個實作共用的編碼與檢查 \\
\bottomrule
\end{tabularx}
\end{center}

\needspace{20\baselineskip}
\section{隔離之一：限制 /dev/mem、拿掉保留區}
\label{sec:devmem}

\begin{oneline}
driver 自己配 DMA buffer 之後，v1.0.0 為 \code{/dev/mem} 保留的 256 MB 就拿掉了；再打開 \code{STRICT_DEVMEM} 與 \code{IO_STRICT_DEVMEM}，\code{/dev/mem} 不能再映射 kernel 管理的記憶體（System RAM）與 driver 佔用中的暫存器。
\end{oneline}

\begin{before}
device tree 保留 256 MB 給 app 經 \code{/dev/mem} 映射（第 3 章），kernel 沒有限制 \code{/dev/mem}。v1.0.0 開機時的紀錄：保留區佔掉 256 MB（\code{MemTotal} 只有 1,773,764 kB）；\code{STRICT_DEVMEM} 沒開；任何 root 程式都能用 \code{devmem} 讀寫 DDR 和加速器的暫存器：

\boardfile{D5_boot_state.log}{26-43}
\end{before}

\begin{tech}
\begin{itemize}
\item \textbf{拿掉 \code{reserved-memory}}：driver 的 buffer 由 DMA API 配，不再需要保留區。
\item \code{CONFIG_STRICT_DEVMEM}：\code{/dev/mem} 不能存取 System RAM（\code{/proc/iomem} 裡標成 \code{System RAM} 的範圍）。
\item \code{CONFIG_IO_STRICT_DEVMEM}：\code{/dev/mem} 也不能存取被 driver 佔用（busy）的 I/O 區間。
\item kernel 設定片段（\code{.cfg}）加 \texttt{linux-xlnx\_\%.bbappend}：在 PetaLinux 預設的 kernel 設定上只改這幾個選項。
\end{itemize}
\end{tech}

\begin{gain}
\begin{itemize}
\item kernel 可用的記憶體多了約 \memGainMB{} MB：\code{MemTotal} 從 1,773,764 kB 變成 2,035,396 kB。兩次開機的 kernel 不同（v1.1 多開了 SMMU、\code{STRICT_DEVMEM} 等選項），這個差裡也混著兩顆 kernel 本身大小的變化。
\item \code{/dev/mem} 映射 System RAM 與加速器暫存器都被拒絕（\code{Operation not permitted}）。
\item driver 載入期間，user space 只能經過 driver 碰暫存器與 buffer，「只給 buffer、不給暫存器」（第 4 章）不能再被 \code{/dev/mem} 繞過。這個保證有邊界：凡是「不是 System RAM、也沒有 driver 佔用」的位址，\code{/dev/mem} 都還能映射，例如 root 卸載 driver 之後的暫存器、DDR 最上面 1 MB（下面實測）。
\item 代價：v1.0.0 的 UIO 版本 app 在 v1.1 映像上不能跑，v1.1 的 app 一律用 driver。
\end{itemize}
\end{gain}

\subsection{實際例子}

\textbf{kernel 設定。}PetaLinux 的 kernel recipe 把 \code{SRC_URI} 裡的 \code{.cfg} 片段合併進 kernel 設定。片段的註解說 \code{/dev/mem} 不能再映射「any RAM」，精確地說是 System RAM：

\snip[{}]{kcfg}
\snip[{}]{kbbappend}

\textbf{板上：發佈映像開機後。}\code{reserved-memory} 節點不在了，兩個選項都開了，\code{devmem} 讀 DDR 與暫存器都被拒絕：

\boardfile{H3_boot_state.log}{16-31}

\textbf{限制的邊界。}\code{v11_check.sh devmem} 只讀不寫，依序試 System RAM 裡的一個位址、driver 載入時的暫存器、卸載後的暫存器、重新載入後的暫存器，最後是 System RAM 以外的 DDR：

\boardfile{H8_devmem.log}{1-20}

System RAM 到 \code{0x7fef_ffff} 為止。System RAM 與 driver 佔用中的暫存器都被拒絕；\code{rmmod} 之後 \code{/proc/iomem} 裡沒有 \code{twn_fuse}，同一個暫存器讀得到（\code{0x11} = 17，NBLK 的 reset 值），重新載入後又被拒絕；\code{0x7ff0_0000} 不是 System RAM、也沒有 driver 佔用，所以讀得到。這最上面的 1 MB 不是閒置的記憶體：Vivado 建 block design 時就報告它保留給 PMU（平台管理單元）內部使用，PetaLinux 產生的 \code{/memory} 節點因此沒有涵蓋它。\code{devmem_is_allowed()} 不分讀寫，這裡只測了讀：

\boardfile[：PC 上的 Vivado 建置訊息]{S8_pmu_ddr_reserved.log}{2-3}

\subsection{原理}

\textbf{\code{/dev/mem} 為什麼危險。}它把整個實體位址空間開給 root：kernel 的程式碼與資料、其他 process 的記憶體、所有裝置的暫存器。拿到 root 的程式能直接改 kernel，也能繞過 driver 去動硬體，driver 做的 offset 檢查（第 5 章）全被繞過。

\needspace{13\baselineskip}
\textbf{兩個選項各擋什麼。}\code{/dev/mem} 的 \code{mmap()} 對每一頁呼叫 \code{devmem_is_allowed()}（Linux 6.1 \code{lib/devmem_is_allowed.c}），參數 \code{pfn}（page frame number）是實體位址除以頁大小 4096：

\nopagebreak
\begin{lstlisting}[numbers=none,xleftmargin=0pt,title={\scriptsize lib/devmem\_is\_allowed.c（Linux 6.1，節錄）}]
int devmem_is_allowed(unsigned long pfn)
{
	if (iomem_is_exclusive(PFN_PHYS(pfn)))
		return 0;
	if (!page_is_ram(pfn))
		return 1;
	return 0;
}
\end{lstlisting}

\code{STRICT_DEVMEM} 讓這個函式生效：\code{page_is_ram()} 為真就拒絕，依據的是 kernel 管理的記憶體範圍（System RAM），不是「這是不是 DDR」。System RAM 是 \code{/memory} 節點的範圍，扣掉標成 \code{no-map} 的保留區，所以 \code{/memory} 以外、以及 \code{no-map} 保留區裡的 DDR 都不會被擋；v1.0.0 的 256 MB 保留區就是 \code{no-map}，就算當時開了 \code{STRICT_DEVMEM}，UIO 版本也還能經 \code{/dev/mem} 映射它。\code{IO_STRICT_DEVMEM} 讓 \code{iomem_is_exclusive()} 把 busy 的區間也算成獨佔：probe 用的 \code{devm_platform_ioremap_resource()} 會先佔用 \code{reg} 範圍（第 2 章），所以加速器的 64 KB 暫存器被擋下；driver 卸載時釋放這個區間，之後那段暫存器又能映射；重新載入、再次登記時，kernel 會撤銷 \code{/dev/mem} 對這段位址既有的映射（\code{kernel/resource.c} 的 \code{revoke_iomem()}）。沒有登記佔用的周邊也擋不住，例如中斷控制器 GIC 的 driver 用 \code{of_iomap()} 映射、沒有登記，root 仍能經 \code{/dev/mem} 映射它的暫存器（由原始碼推得，沒有在板上測）。要完全不給 \code{/dev/mem}，得關掉 \code{CONFIG_DEVMEM}；這個映像沒有關。

\textbf{為什麼以前需要保留區，現在不需要。}\code{/dev/mem} 只是把一個實體位址映射給 app，不會配置記憶體；要讓 kernel 不去用那塊記憶體，只能在 device tree 事先保留。driver 版本改用 DMA API，記憶體是 kernel 自己配給 driver 的，kernel 知道它被誰用，不需要保留。

\textbf{少掉的 512 kB。}256 MB 是 262,144 kB，\code{MemTotal} 只多了 261,632 kB。kernel 收回這塊記憶體時要為它建立管理結構，例如線性映射（第 3 章）的分頁表；以 4 KB 的頁、每項 8 bytes 計算，256 MB 的分頁表剛好 512 kB。這是推算，沒有逐項查證；而且比較的兩邊是不同設定的 kernel，kernel 映像本身的大小差也混在裡面，沒有拆開。

\subsection{面試會問}

\begin{qa}{為什麼要限制 /dev/mem？限制之後 app 怎麼存取硬體？}
\code{/dev/mem} 把整個實體位址空間開給 root，等於任何 root 程式都能改 kernel、繞過 driver 動硬體。v1.0.0 需要它，因為 app 要自己映射 DMA buffer，device tree 還得保留 256 MB。driver 版本用 DMA API 配 buffer、只把 buffer mmap 給 app，所以 app 不再需要 \code{/dev/mem}，保留區也拿掉了。我打開 \code{STRICT_DEVMEM} 擋 System RAM、\code{IO_STRICT_DEVMEM} 擋 driver 佔用的暫存器，板上 \code{devmem} 都回 \code{Operation not permitted}，kernel 多了大約 \memGainMB{} MB。邊界我也講得出來：不是 System RAM、也沒有 driver 佔用的位址都還能映射，板上實測卸載 driver 後的暫存器、DDR 最上面那 1 MB 的 PMU 保留區都讀得到，所以這是限制，不是把 \code{/dev/mem} 關掉。
\end{qa}

\subsection{關鍵函式速查}

\begin{center}
\small
\begin{tabularx}{\textwidth}{l X p{40mm}}
\toprule
名稱 & 位置（\code{meta-user/} 指 \code{linux/project-spec/meta-user/}） & 做什麼 \\
\midrule
\code{twn-v11.cfg} & \code{meta-user/recipes-kernel/linux/linux-xlnx/twn-v11.cfg}:\lineof{kcfg} & kernel 選項片段 \\
\texttt{linux-xlnx\_\%.bbappend} & \code{meta-user/recipes-kernel/linux/}\texttt{linux-xlnx\_\%.bbappend}:\lineof{kbbappend} & 把片段加進 kernel recipe \\
\code{devmem_is_allowed()} & kernel 的 \code{lib/devmem_is_allowed.c} & \code{/dev/mem} 每一頁的檢查 \\
\code{iomem_is_exclusive()} & kernel 的 \code{kernel/resource.c} & 開了 IO 選項後，busy 的區間也算獨佔 \\
\code{v11_check.sh devmem} & \code{test/v11_check.sh}:\lineof{chk_devmem} & 限制的邊界（只讀） \\
\bottomrule
\end{tabularx}
\end{center}

\needspace{20\baselineskip}
\section{隔離之二：SMMU 限制 PL 能碰的記憶體}
\label{sec:smmu}

\begin{oneline}
v1.1 打開 SMMU（ARM 的 IOMMU）並登記加速器四個 HP 埠的 stream ID，讓帶著這些 ID 的 Non-secure 交易只能存取 driver 映射的那 8 MB；第一次打開沒有效果，因為 PL 送出的是 Secure 交易。
\end{oneline}

\begin{before}[原本怎麼做（沒有 SMMU：v1.0.0 的開機檔，v1.1 的 driver）]
driver 的 \code{in_arena()} 只能保證 driver 自己寫進暫存器的基底位址落在 buffer 內（第 5 章）。PL 執行時用「基底加上描述子裡的偏移」去存取 DDR，描述子由 app 提供，PL 也可能因為 RTL 的錯誤而越界。沒有 SMMU 時，越界的存取不會被任何東西擋下。v1.0.0 的開機檔上（開機組 D，沒有 SMMU），\code{dma_fault}（第 8 章）讓 PL 從 buffer 結尾之後讀輸入，結果「frame done」，沒有任何錯誤：

\boardfile{D3_debugfs.log}{68-71}
\end{before}

\begin{tech}
\begin{itemize}
\item \textbf{SMMU（MMU-500）}：替 CPU 以外的匯流排主控端（PS 的 DMA 周邊，以及 PL 經 HP 埠進來的交易）做位址轉換與權限檢查，概念和 CPU 的 MMU 相同。
\item \textbf{stream ID}：每筆交易進 SMMU 時帶的來源編號。ZynqMP 的組成是「TBU 編號（bit 14:10）、埠的 master ID（bit 9:6）、AXI ID（bit 5:0）」。
\item \textbf{device tree 的 \code{iommus}}：告訴 kernel 這個裝置用哪些 stream ID；kernel 就把裝置放進一個 IOMMU group（一組必須共用同一套轉換的裝置），並給它一個 IOMMU domain（一套轉換表）。
\item \textbf{DMA API 與 IOMMU}：driver 的資料路徑不用改，\code{dma_alloc_coherent()} 在有 IOMMU 時自動配一段 IOVA、建好映射，回傳的 \code{dma_addr_t} 就是 IOVA（第 3 章）。只有量測用的 \code{map_device} 路徑加了一個檢查：有 IOMMU 時 buffer 沒有固定的實體位址可以映射，driver 拒絕（第 4 章）。
\item \textbf{AxPROT 與 TrustZone}：AXI 交易的 bit 1 標示 Secure 或 Non-secure；Linux 只設定 SMMU 的 Non-secure 那一側。
\item \textbf{診斷參數 \code{arm_smmu.disable_bypass=1}}：沒登記的 stream ID 一律擋下並回報，用來找出真正的 stream ID。發佈的 kernel 沒有加這個參數，沒登記的 stream ID 直接通過（bypass），見「原理」最後一段。
\end{itemize}
\end{tech}

\begin{gain}
\begin{itemize}
\item PL 越界的存取被擋下：同一個越界讀取測試，SMMU 打開時 kernel 回報 context fault，關掉時只有「frame done」。前提是交易是 Non-secure、帶著登記過的 stream ID（「原理」最後一段）。
\item 正確性與效能不受影響：同一組開機檔只換 \code{system.dtb}，SMMU 開與關的 4 張 fixture 都逐位元相同，pipeline 三次平均 \smmuOnFps{} 對 \smmuOffFps{} fps，差距在兩次執行之間的正常變動之內。
\item driver 的資料路徑不用改：DMA API 自己處理 IOVA 與頁表。
\item v1.1 的發佈映像預設就打開 SMMU；PS 的 DMA 周邊也因此一起被放到 SMMU 後面（Ethernet 與 SD 卡在所有測試裡都用到，USB 與 DisplayPort 沒有測）。
\end{itemize}
\end{gain}

\subsection{第一次：打開了，卻沒有效果}

\code{smmu_dtb.sh} 在 \code{system.dtb} 的副本上做兩件事：打開 SMMU 節點，並在加速器節點寫入 stream ID。第一次用的版本寫入的是我推算的三個 ID（\code{0xe80}、\code{0x1280}、\code{0x1680}）；下面印出的是量到 ID 之後的現行版本，註解與 \code{iommus} 的四個值都是量到的：

\snip[bash]{smmu_dtb}

開機後 SMMU 正常啟動，加速器也被放進自己的 IOMMU group，buffer 位址變成 IOVA \mbox{\code{0xff_ff80_0000}}：

\boardfile{first_bitstream/F1_smmu_guessed_ids.log}{11-20}

但四張 fixture 全錯，而且越界讀取也沒有觸發任何 SMMU 錯誤。再用 \code{arm_smmu.disable_bypass=1} 開機（沒登記的 stream ID 應該被擋下並回報），kernel log 裡仍然一條 SMMU 錯誤都沒有。這次開機沒有保存紀錄，\code{first_bitstream/README.md} 只記下「沒有任何 fault」。兩個現象合起來的意思是：PL 的交易根本沒有經過 Linux 設定的那一套轉換。

\textbf{原因：AXI master 送的是 Secure 交易。}加速器的 AXI master 把 AxPROT 接成常數 \code{3'b000}，bit 1 是 0，代表 Secure。SMMU 對 Secure 與 Non-secure 交易各有一套設定；Linux 跑在 Non-secure，只能設定 Non-secure 那一套。Secure 那一套不是 Linux 設定的；從板上的結果看，它直接放行、不轉換。所以 PL 拿 IOVA 當成實體位址去存取，讀寫到錯的地方。

\snip[Verilog]{axi_rd_prot}

修正是把 bit 1 設成 1（\code{3'b010}：Non-secure、非特權、資料存取）。\code{twn_axi_rd.v} 與 \code{twn_axi_wr.v} 是加速器所有 AXI master 共用的模組，改這兩行就涵蓋四個 HP 埠。

\subsection{第二次：找出真正的 stream ID}

改成 Non-secure 交易之後，HP1、HP2、HP3 的 stream ID 仍然是我推算的，推算可能是錯的。這時 \code{arm_smmu.disable_bypass=1} 就派上用場：SMMU 遇到沒有登記的 stream ID 時會擋下交易，並在 global fault 的 \code{GFSYNR1} 欄位回報它的 stream ID（開機組 F）：

\boardfile[：只登記了推算的 0xe80、0x1280、0x1680]{F0_smmu_find_stream_ids.log}{2-3,5-8,11-12,17-22}

HP0 的交易已經經過 Non-secure 的轉換：第 22 行有一個 context fault 帶著 \code{cbfrsynra=0xe80}（這次開機做過哪些存取沒有記下，這一筆是哪一個存取造成的沒有查）。另外三個埠的 ID 沒有登記，SMMU 回報 global fault，沒登記的 ID 是 \code{0x12c0}、\code{0x1300}、\code{0x1740}。\code{uniq -c} 數到的 6、1、3 次加起來正好 10，等於 kernel 印這類訊息的速率上限（5 秒內最多 10 次），所以這些次數不代表各埠的流量，某個 ID 也可能被限流吃掉；\code{F1_smmu_measured_ids.log} 登記了四個 ID 之後一次 global fault 都沒有，補上了這個漏洞。拆開來看：TBU 4 的 master ID 1011、TBU 4 的 1100、TBU 5 的 1101，正好接在 HP0 的 1010 後面。所以 HP1 是 \code{0x12c0}、HP2 是 \code{0x1300}、HP3 是 \code{0x1740}；我原本以為 HP1 與 HP2 共用 \code{0x1280}、HP3 是 \code{0x1680}，兩個都錯了。

把四個 ID 都登記上（現行的 \code{smmu_dtb.sh}），同樣開著 \code{disable_bypass}：fixture 全部逐位元相同，越界讀取被擋下（context fault），global fault 一個都沒有，表示 PL 的每一筆交易都找得到登記：

\boardfile[：四個 ID 都登記]{F1_smmu_measured_ids.log}{16-21,30,37-38,51-54}

\subsection{發佈映像：預設打開 SMMU}

量到 stream ID 之後，\code{system-user.dtsi} 直接打開 SMMU 並登記這四個 ID（第 1 章的 \code{fuse_0} 節點），PetaLinux 產生的節點因此帶著 \code{iommus}。發佈映像開機後，加速器在 iommu group 15，DMA buffer 的位址是 IOVA；越界讀取被 SMMU 擋下、回報 context fault，\code{cbfrsynra=0xe80} 就是 HP0 的 stream ID：

\boardfile{H6_smmu.log}{11-15}
\boardfile{H6_smmu.log}{30-31,37-38}

同一個測試在 SMMU 關掉時（開機組 G）只說 frame done，buffer 位址是 CMA 裡的實體位址：

\boardfile[：SMMU 關]{G4_debugfs.log}{37-40}

context fault 裡的 \code{iova=0x00000000} 是因為 buffer 剛好在 $2^{40}$ 下面，buffer 結尾之後的第一個位址是 $2^{40}$，寫進 40 bits 的位址暫存器就變成 0；不論是哪一個，都不在映射範圍內。擋下之後，下一個 app 照常執行，4 張 fixture 仍然正確（\code{H6_smmu.log} 第 39 行起）。

開與關的效能比較要用同一組開機檔、只換 \code{system.dtb}，這是開機組 E 與 G（發佈的開機檔加候選版 rootfs；它的 driver 和發佈版只差兩個 probe 錯誤路徑，不影響這裡量的東西）：

\begin{center}
\small
\begin{tabularx}{\textwidth}{X r r}
\toprule
同一個映像，只差 \code{system.dtb} & E：SMMU 開 & G：SMMU 關 \\
\midrule
pipeline，三次的平均（\code{E2}/\code{G2} 第 37--39 行） & \smmuOnFps{} fps & \smmuOffFps{} fps \\
每幀 CPU 時間，中斷（\code{E3}/\code{G3} 第 10 行） & \irqCpuUs{} µs & 10.8 µs \\
等待延遲，中斷（第 9 行） & \irqWaitUs{} µs & 14.3 µs \\
越界讀取（\code{E4}/\code{G4} 第 37--41 行） & context fault & 沒有察覺 \\
\bottomrule
\end{tabularx}
\end{center}

\textbf{為什麼 SMMU 沒有拖慢 PL。}每次轉換都要查頁表，但 SMMU 的 TBU 有自己的 TLB（轉換結果的快取）；PL 又是順序地串流讀寫同一塊 buffer，大部分的存取會落在剛查過的頁上。量到的吞吐量差距在正常變動之內；TLB 命中率本身沒有量，用了多大的頁也沒有查。buffer 的實體頁在 SMMU 開與關時從哪裡來，見第 3 章。

\subsection{原理}

\textbf{SMMU 怎麼轉換。}SMMU 的結構和 CPU 的 MMU 很像：每個 IOMMU domain 有一套頁表，裝置發出的位址（IOVA）經頁表轉成實體位址，沒有映射的位址會產生 context fault。kernel 的 DMA API 在有 IOMMU 時，從裝置的 DMA mask 往下配 IOVA（這裡的 mask 是 40 bits，所以 IOVA 是 \mbox{\code{0xff_ff80_0000}}，緊貼在 $2^{40}$ 下面），再把 buffer 的實體頁映射上去。

\textbf{stream ID 怎麼算。}ZynqMP 的 SMMU 有 6 個 TBU（轉換單元），各自負責幾個匯流排埠。PL 經 HP0 進來的交易走 TBU 3，HP1、HP2 走 TBU 4，HP3 走 TBU 5；HP0 的 master ID 是 \code{1010}，HP1 到 HP3 依序是 \code{1011}、\code{1100}、\code{1101}（開機組 F 量到的）。AXI ID 由 RTL 決定，這個設計的 master 都送 0。所以 HP0 是 \code{(3 << 10) | (0b1010 << 6) = 0xe80}。

\textbf{為什麼診斷時也看不到錯誤。}\code{disable_bypass} 只影響 Non-secure 那一套設定裡「沒有登記的 stream ID」。Secure 交易根本不走這一套，所以不會被擋，也不會被回報。這種「完全沒有反應」比錯誤訊息更難查：要從「測試的結果和預期的差異，連錯誤都沒有」反推出交易沒有進到被設定的路徑。

\textbf{SMMU 能保護什麼、不能保護什麼。}對帶著登記過的 stream ID 的 Non-secure 交易，它保證 PL 只能存取 driver 為它映射的頁，描述子給錯的偏移、或 PL 的位址計算錯誤，都碰不到 kernel 和其他 process 的記憶體。但有兩種交易它擋不住：
\begin{itemize}
\item \textbf{沒有登記的 stream ID}：發佈的 kernel 沒有開 \code{ARM_SMMU_DISABLE_BYPASS_BY_DEFAULT}，命令列也沒有 \code{arm_smmu.disable_bypass=1}，所以 SMMU 對沒登記的 ID 直接放行。stream ID 的低 6 位元是 AXI ID，只要 PL 送出非 0 的 AXI ID（RTL 的錯誤，或換了另一顆 bitstream），交易就以實體位址直通整個 DDR。
\item \textbf{Secure 交易}（AxPROT[1] = 0）：完全不經過 Linux 設定的轉換。這正是本章第一次失敗的情形，也是一種 RTL 錯誤。
\end{itemize}
要擋住第一種，可以在命令列加 \code{arm_smmu.disable_bypass=1}；發佈映像沒有這樣做，因為它會擋下所有沒有登記 stream ID 的主控端，這個組合只在開機組 F 用過，沒有完整測過 PS 的各個周邊。第二種只能靠 RTL 本身正確。另外，SMMU 也不檢查 buffer 內部的存取是否合理（例如輸出寫到權重區），那仍然要靠 driver 的 offset 檢查。

\subsection{面試會問}

\begin{qa}{你怎麼用 SMMU 隔離 FPGA 的 DMA？遇到什麼問題？}
在 device tree 打開 SMMU，給加速器的節點寫 \code{iommus}，列出它 HP 埠的 stream ID；stream ID 是 TBU 編號、埠的 master ID 和 AXI ID 拼起來的。driver 的資料路徑不用改，DMA API 在有 IOMMU 時自動配 IOVA、建頁表，帶著登記過的 ID 的交易只能存取那 8 MB。第一次打開時 fixture 全錯，而且連 \code{disable_bypass} 都看不到任何錯誤，表示交易根本沒進 Linux 設定的那套轉換。原因是我的 AXI master 把 AxPROT 接成 0，送的是 Secure 交易，而 Linux 只設定 SMMU 的 Non-secure 那一側。把 AxPROT[1] 改成 1 之後，越界讀取會變成 context fault。後來發現另外三個埠的 stream ID 我推算錯了兩個，就用 \code{disable_bypass} 讓 SMMU 回報沒有登記的 ID，量出真正的值。開 SMMU 前後的吞吐量是 \smmuOnFps{} 對 \smmuOffFps{} fps，看不出差別，所以發佈映像預設打開。它的限制我也會講：發佈的 kernel 對沒登記的 stream ID 是 bypass，PL 若送出非 0 的 AXI ID 或 Secure 交易，SMMU 擋不住；要更嚴格得加 \code{disable_bypass}，代價是每個 PS 周邊都得登記。
\end{qa}

\subsection{關鍵函式速查}

\begin{center}
\small
\begin{tabularx}{\textwidth}{l l X}
\toprule
名稱 & 位置 & 做什麼 \\
\midrule
\code{smmu_dtb.sh} & \code{linux/driver/twn_fuse/smmu_dtb.sh}:\lineof{smmu_dtb} & 打開 SMMU、寫入 \code{iommus} \\
\code{arprot}、\code{awprot} & \code{rtl/src/twn_axi_rd.v}:\lineof{axi_arprot}、\code{twn_axi_wr.v} & Non-secure 交易 \\
\code{twn_fault_set()} & \code{twn_fuse.c}:\lineof{drv_fault_set} & 越界讀取的測試 \\
\code{map_device} 的檢查 & \code{twn_fuse.c}:\lineof{drv_mmap_iommu} & 有 IOMMU 時拒絕 Device 映射 \\
\bottomrule
\end{tabularx}
\end{center}

\needspace{20\baselineskip}
\section{改了 RTL 就要重建：timing closure}
\label{sec:timing}

\begin{oneline}
v1.1 改了兩次 RTL（加中斷、把 AXI 交易改成 Non-secure），兩次都用 v1.0.0 發佈時的建置配方重建並在 240 MHz 收斂，第二次 route 後是 $-$0.032 ns，靠 route 之後的 \code{phys_opt_design}（strategy 內建一次，建置腳本再補兩次）補到 0.000 ns，和 v1.0.0 發佈時的情形一樣。
\end{oneline}

\begin{before}[原本怎麼做（v1.0.0 的 bitstream）]
v1.0.0 的 bitstream（\code{52b400785924efe0}）在 240 MHz 收斂得很勉強：第一次 route 後 WNS（Worst Negative Slack，最差路徑的時間餘裕；負的表示這條路徑來不及）是 $-$0.030 ns，建置腳本自動跑 post-route \code{phys_opt_design}，補到 0.000 ns，完全沒有餘裕。出處：\code{results/board_results.json} 的 \code{timing_bitstream_52b400785924efe0}，報告原文在 \code{results/raw/thesis_model_20261012/timing_impl_fuse_ternpose.rpt}。

這個設計的 LUT（look-up table，FPGA 裡實作組合邏輯的基本單元）用了 93\%、block RAM 用了 97\%，改一點點 RTL，place \& route（placement 決定每個邏輯單元放在晶片的哪個位置，route 決定它們之間怎麼繞線）的結果就可能不同。
\end{before}

\begin{tech}
\begin{itemize}
\item \textbf{固定的建置配方}：\code{vivado/build.cmd} 裡的變數，與 timing 有關的是 implementation strategy（Vivado 預先組好的一整組 place \& route 選項）\code{Performance_NetDelay_high}、HP3 拆出獨立的 interconnect（\code{FUSE_HP3_SPLIT}）、一個列緩衝區（\code{RB_C}，RTL 裡叫 RingC）做成三列深（\code{FUSE_CC_ROWS3}）、一組 block RAM 的位置限制（pblock，\code{FUSE_T1_PBLOCK}）；其餘是建出 head 引擎（\code{FUSE_HD_PROBE}）、套用開發板的 PS 設定（\code{BOARD_PS}）、跑 implementation 並寫出 bitstream（\code{FUSE_IMPL}、\code{FUSE_BIT}）。兩次都沒改。
\item \textbf{post-route 恢復}：\code{build_bd_fuse.tcl} 在 route 之後若 WNS < 0，依序最多跑三次 \code{phys_opt_design}（\code{AggressiveExplore}、\code{AlternateFlowWithRetiming}、再一次 \code{AggressiveExplore}；directive 是同一個指令的不同最佳化模式），還不行就重新 route 再跑一次。
\item \textbf{資料路徑檢查}：合成後與 route 後各查一次，25 個資料路徑實例都還在 netlist（合成產生的電路元件與連線清單）裡。
\item \textbf{備案：重現加 ECO}：用同一個 commit 在另一個資料夾重建，確認 Vivado 到 placement 為止給出一模一樣的中間結果，準備在已經收斂的設計上只改那幾根常數線（ECO，engineering change order：直接修改已經 route 好的設計，不重跑整個流程）。這次沒有用到。
\end{itemize}
\end{tech}

\begin{gain}
\begin{itemize}
\item 最終的 bitstream（\code{11b4d88f583b8e23}）收斂：WNS 0.000 ns、WHS +0.008 ns，和 v1.0.0 一樣靠 post-route 恢復。
\item 資源幾乎不變：LUT 44,157（v1.0.0 是 44,057，多 0.2\%），block RAM tile（晶片上的記憶體區塊）145 個、DSP（硬體乘法器單元）217 個都相同。
\item 可以從 repo 的原始碼重現：\code{build.cmd} 不需要任何手動步驟。
\end{itemize}
\end{gain}

\subsection{實際例子}

\textbf{建置配方。}\code{build.cmd} 把配方寫成環境變數。兩次重建只設了 \code{FUSE_TAG}（\code{v11irq}、\code{v11}），讓輸出檔名不蓋掉 v1.0.0 版：

\snip[bash]{build_recipe}

\textbf{post-route 恢復。}第一次重建 route 完已經是正的，這段第一行就跳出；第二次重建跑了前兩次：

\snip[tcl]{postroute}

\textbf{最終那次建置的紀錄}（只留 timing 與檢查的行；第 4 行確認 \code{FUSE_CC_ROWS3} 生效，第 13--14 行是 pblock 的結果，見後文）：

\boardfile[：PC 上的 Vivado 2023.2]{S4_build_timing.log}{1,3-14}

由上往下讀：placement 後 $-$0.472 ns；placement 後的 \code{phys_opt} 拉到 $-$0.116 ns；route 後 $-$0.032 ns；strategy 自己在 route 後的 \code{phys_opt} 到 $-$0.015 ns；接著建置腳本的恢復：\code{AggressiveExplore} 到 $-$0.006 ns，\code{AlternateFlowWithRetiming} 到 0.000 ns。

\textbf{三次建置比一比}（都是 Vivado 2023.2、同一份配方、同一顆 xczu2cg-sfvc784-1-e）：

\begin{center}
\small
\begin{tabularx}{\textwidth}{X r r r}
\toprule
 & v1.0.0（發佈） & v1.1 第一版 & v1.1 最終 \\
\midrule
RTL 的差別 & & 加中斷 & 再改 AxPROT[1] \\
placement 後的 WNS & 公開紀錄沒有 & $-$0.295 ns & $-$0.472 ns \\
route 後的 WNS & $-$0.030 ns & +0.037 ns & $-$0.032 ns \\
最後的 WNS／WHS & 0.000／0.002 ns & +0.037／0.008 ns & 0.000／0.008 ns \\
post-route 恢復 & 用到 & 沒用到 & 用到 \\
CLB LUT & 44,057 & 44,081 & 44,157 \\
CLB Register & 36,339 & 36,253 & 36,203 \\
Block RAM tile／DSP & 145／217 & 145／217 & 145／217 \\
\bottomrule
\end{tabularx}
\end{center}

出處：v1.0.0 見「原本怎麼做」；v1.1 第一版是 \code{first_bitstream/S4_*}，最終是 \code{S4_*}（報告第 142 行是 WNS／WHS，第 10 行的 Design State 是報告產生時設計所在的步驟，從它看得出有沒有跑 post-route \code{phys_opt}）。表中的 CLB（configurable logic block）是放 LUT 與暫存器的基本區塊。AxPROT 是 AXI 交易附帶的保護屬性，第 11 章解釋為什麼要改它的 bit 1。

\textbf{pblock 的變化。}配方裡的 pblock 希望某一組 block RAM（\code{RB_T1}）放在第 0 到第 2 欄。第一版有 6 顆落在範圍外（\code{first_bitstream/S4_build_timing.log} 第 11 行），最終版有 5 顆（紀錄第 13--14 行）；v1.0.0 版沒有公開的紀錄。它在這個配方裡只是建議性的位置限制，落在外面不算錯；若之後重建 timing 收不了，這是第一個要看的地方。

\textbf{備案：重現加 ECO。}第二次重建 placement 後明顯比較差，當時不確定能不能收斂，所以同時在另一個資料夾，用 v1.1 第一版的 commit 重建。它的合成報告和第一版相同（日期、主機名稱、指令那幾行除外），placement 後 \code{phys_opt} 的第一個估計值也一模一樣（\code{WNS=-0.622}、\code{TNS=-443.497}；TNS 是所有違規路徑的負餘裕加總，\code{first_bitstream/S5_rebuild_determinism.log}）：到 placement 為止，同樣的輸入給出同樣的結果。這次重建在 placement 中途就停了，所以「可以重現第一版那個 +0.037 ns 的設計」只是推測，沒有跑完驗證。計畫是重現之後，用 ECO 把 PS 端六根 AxPROT[1] 的 pin 從接地改接 VCC：這幾根線都是常數，沒有任何 timing 路徑，改了也不會動到 slack。最終的重建先收斂了，這個備案就停掉，沒有執行。

\subsection{原理}

\textbf{只改常數，為什麼 timing 差這麼多。}place \& route 是一個從某個起點出發的最佳化。netlist 只要有一點不同，起點就不同，結果像換了一個亂數種子。AxPROT[1] 從 0 改成 1（\code{3'b000} → \code{3'b010}）看起來只是常數，但 interconnect 裡若有暫存器存這些位元（這是推論）：資料是 0、reset 也是 0 時，暫存器的輸出永遠是 0，合成時會被當成常數刪掉；資料改成 1 之後，reset 期間輸出 0、之後輸出 1，就不再是常數，必須保留。最終版比第一版多了 76 個 LUT，就是 netlist 確實變了的證據（哪些邏輯變了，沒有逐一比對）。

\textbf{同一份配方，結果為什麼不一樣。}v1.0.0 版原本的建置（\code{board_results.json} 的 \code{opt1_archived_20261008}）第一次 route 是 $-$0.030 ns、LUT 44,057；2026-10-11 從乾淨的 clone 重建（\code{rebuilt_from_clone_20261011}），第一次 route 也是 $-$0.030 ns、資源數相同（兩次的 bitstream 雜湊不同，相同的是 WNS 與資源數），表示同樣的輸入會得到同樣的結果；這次的 S5 到 placement 為止也是如此。但輸入一變，結果就在 $-$0.03 到 +0.04 ns 之間跳動。這個設計在 240 MHz 本來就沒有餘裕，所以 post-route 恢復是配方的必要部分，不是偶爾才需要的補救。

\textbf{為什麼 ECO 不會影響 timing。}timing 分析看的是暫存器到暫存器之間的路徑。接到常數的 pin 沒有起點暫存器，不在任何路徑上；把它從接地改接 VCC，只是換一條常數線，其他所有 cell 的位置與繞線都不變。

\textbf{retiming 要再驗一次。}最後一步用了 \code{AlternateFlowWithRetiming}，它會做 retiming：把暫存器搬過組合邏輯，讓前後兩段路徑的延遲比較平均。工具保證功能不變，但這是會改變電路結構的最佳化，所以上板後仍然重跑了全部的逐位元比對（第 13 章，開機組 E 與 H）。

\subsection{面試會問}

\begin{qa}{改了 RTL 之後 timing 收不了，你會怎麼做？}
先確定建置配方沒變，差異只來自 RTL。這個設計 LUT 用 93\%、block RAM 用 97\%，在 240 MHz 本來就沒什麼餘裕，所以我的建置腳本在 route 後若 WNS 小於 0，會自動依序跑幾種 \code{phys_opt_design}，還不行就重新 route。這次我改了兩次 RTL：加中斷那次 route 完就是 +0.037 ns；把 AxPROT[1] 從 0 改成 1 那次，只改常數，placement 卻差了 0.18 ns，route 完 $-$0.032 ns，靠 route 後的 \code{phys_opt_design}（strategy 內建一次、建置腳本再兩次）補到 0。當時我也準備了備案：先確認 Vivado 對同樣的輸入、到 placement 為止給出同樣的結果，打算重現上一版收斂的設計，再用 ECO 只改那幾根常數 pin，timing 一點都不會動。最後沒有用到，重現也沒有跑完。因為最後一步用了 retiming，上板後我重跑了全部的逐位元比對。
\end{qa}

\subsection{關鍵函式速查}

\begin{center}
\small
\begin{tabularx}{\textwidth}{l l X}
\toprule
名稱 & 位置 & 做什麼 \\
\midrule
建置配方 & \code{vivado/build.cmd}:\lineof{build_recipe} & strategy、時脈、pblock 等變數 \\
post-route 恢復 & \code{vivado/build_bd_fuse.tcl}:\lineof{postroute} & WNS < 0 時跑 \code{phys_opt}，再不行重新 route \\
\code{FUSE_TAG} & \code{vivado/build.cmd}:\lineof{build_tag} & 輸出檔名，不蓋掉 v1.0.0 版 \\
\bottomrule
\end{tabularx}
\end{center}

\needspace{20\baselineskip}
\section{板上驗證與靜態檢查：證明了什麼、沒證明什麼}
\label{sec:verify}

\begin{oneline}
\code{board_check.sh} 一條指令跑完四類檢查（輸出是否逐位元相同、該拒絕的有沒有拒絕、卸載重載是否正常、吞吐量），發佈映像上全部通過；v1.1 新增的機制各有 \code{v11_check.sh} 的一個模式；程式另外過了 \code{W=1} 與 \code{checkpatch}，官方的 \code{dt_binding_check} 與 sparse 沒有跑。
\end{oneline}

\begin{before}
v1.0.0 的板上檢查是 \code{board/deploy.sh} 的 \code{verify}（4 張 fixture）與 \code{speed}（吞吐量）：

\snip[bash]{deploy_verify}

只檢查「算得對不對」。沒有錯誤路徑可以檢查，因為 UIO 版本根本不拒絕任何東西。
\end{before}

\begin{tech}
\begin{itemize}
\item bit-exact gate：4 張 fixture 的輸出和 golden 逐位元比對，head 在 PL 和在 CPU 兩種擺法都跑。公開 repo 用的是合成的 fixture（\code{models/mbv2_pose/fixtures_synthetic/}，第 14 章），MPII 的 fixture 只在作者的板子上跑。
\item \code{twn_errtest}：\errChecks{} 項檢查，其中 \errRefusals{} 項故意做錯一件事、確認 driver 回傳預期的錯誤碼，另外 2 項確認正常的操作會成功；參數 \code{timeout} 加跑逾時，\code{hold} 只留映射。
\item 卸載測試：裝置開著時、以及只剩映射時，\code{rmmod} 都必須失敗；之後連續 5 次卸載、重載，每次都跑 4 張 fixture。
\item 吞吐量：pipeline 跑 120 幀、3 次，每次都檢查關節點和 fixture 0 相同。
\item 靜態檢查：\code{make W=1}（kernel 額外的編譯器警告）與 \code{scripts/checkpatch.pl --strict}（kernel 的程式風格檢查）。
\end{itemize}
\end{tech}

\begin{gain}
\begin{center}
\small
\begin{tabularx}{\textwidth}{X l l}
\toprule
發佈的 driver 上的檢查（沒有註明的都是開機組 H） & 結果 & 紀錄 \\
\midrule
4 張合成 fixture，head 在 PL／在 CPU & 4/4、4/4 & \code{H4} \\
4 張 MPII fixture，同上 & 4/4、4/4 & \code{H4m} \\
錯誤路徑 \errChecks{} 項（\errRefusals{} 項拒絕、2 項成功） & \errChecks/\errChecks{} 符合預期 & \code{H4} 第 7--23 行 \\
被拒絕的表之後，下一個 app 再跑 & 4/4、4/4 & \code{H4} 第 26--27 行 \\
裝置開著時／只剩映射時 \code{rmmod} & 都被拒絕 & \code{H4} 第 29--30 行 \\
5 次卸載重載，每次 4 張 fixture & 5 次都 4/4 & \code{H4} 第 31--35 行 \\
SMMU 關（開機組 I）：10 次各要 64 MB，總量超過 CMA & 10/10 & \code{I1} \\
SMMU 開：10 次各要 64 MB 的重載 & 10/10 & \code{H10} \\
pipeline 吞吐量（3 次） & \relFpsA{}／\relFpsB{}／\relFpsC{} fps & \code{H4} 第 37--39 行 \\
\code{W=1}、\code{checkpatch --strict} & 0 個警告 & \code{S6} \\
\bottomrule
\end{tabularx}
\end{center}
\end{gain}

\subsection{實際例子}

\textbf{fixture 檢查。}每一種擺法跑一次 app，數有幾行 \code{PASS | final max|diff|=0}，要剛好 4 行：

\snip[bash]{check_fixtures}

\textbf{錯誤路徑測試的寫法。}每一項都是「做一件事，看回傳值和錯誤碼對不對」：

\snip{test_check}
\snip{test_call}

\begin{pit}
\code{twn_errtest} 第一次跑，兩項 \code{mmap} 檢查顯示 \code{FAIL ... -> Success}，看起來像是 driver 沒拒絕。其實錯在測試程式：當時 \code{mmap()} 寫在 \code{CALL} 外面，\code{CALL} 先把 \code{errno} 設成 0 才檢查，於是 \code{mmap} 留下的 \code{EINVAL} 被蓋掉，印出來的「Success」其實是 \code{strerror(0)}。\code{mmap} 本身是失敗的，driver 確實拒絕了。修法是把 \code{mmap()} 放進 \code{CALL} 裡，讓 \code{errno} 在呼叫完的下一刻就被讀取。教訓是：測試說「沒拒絕」時，先檢查測試本身有沒有讀對東西，不要直接去改受測的程式。
\end{pit}

\textbf{發佈映像上的一次 board check。}中間的錯誤路徑與卸載已經在第 4 到 7 章逐段列出，這裡只印開頭與吞吐量：

\boardfile{H4_board_check.log}{1-6}
\boardfile{H4_board_check.log}{36-40}

\textbf{v1.1 新增機制的檢查。}\code{v11_check.sh} 的每個模式只測一件事，每條指令先印成一行 \code{\# ...} 再印輸出，紀錄因此可以逐行引用：

\begin{center}
\small
\begin{tabularx}{\textwidth}{l X l}
\toprule
模式 & 測什麼 & 紀錄與章節 \\
\midrule
\code{ab} & 四組 A/B 的吞吐量（要 UIO，只能在 v1.0.0 的開機檔上跑） & \code{D1}，第 3 章 \\
\code{bench} & ioctl 成本、backbone 時間拆解、CPU 存取 buffer 的頻寬 & \code{D2}、\code{E3}，第 3、5、6 章 \\
\code{split} & pipeline 各階段的時間，兩種映射 & \code{D6}，第 3 章 \\
\code{irqcount} & 1000 幀對應 1000 次中斷，沒有空中斷 & \code{E6}，第 6 章 \\
\code{irqfps} & 同一次開機，中斷與輪詢交替量吞吐量 & \code{H7}，第 6 章 \\
\code{debugfs} & debugfs 檔案、短幀、不合法的值、越界讀取 & \code{D3}、\code{E4}，第 8、11 章 \\
\code{synthetic} & 合成 fixture，兩種 head 擺法 & \code{D4}、\code{E5}、\code{H5}，第 9 章 \\
\code{boot} & 開機自動載入、中斷、\code{/dev/mem}、記憶體、密碼期限 & \code{H3}，第 10、14 章 \\
\code{devmem} & \code{/dev/mem} 還碰得到什麼（只讀） & \code{H8}，第 10 章 \\
\code{smmu} & IOMMU group、fixture、吞吐量、越界讀取 & \code{H6}，第 11 章 \\
\code{state} & 載入中的 driver：module 資訊、綁定、buffer、sysfs、映射、10 次重載 & \code{H10}、\code{I1}，第 1、3、8 章 \\
\bottomrule
\end{tabularx}
\end{center}

\textbf{driver 版本和 UIO 版本比快慢。}driver v1.0 剛完成時（開機組 A：v1.0.0 的開機檔，driver 靠輪詢），同一次開機、同一份 app 分別連結兩個 HAL，board check 的吞吐量是 UIO \fpsUioMean{} fps、driver \fpsDrvMean{} fps：

\boardfile{A2_board_check_uio_dtb.log}{40-45}

當時 driver 一起改了三件事（等硬體的地方、buffer 位置、CPU 端的記憶體屬性），只能說「沒有變慢」。v1.1 的 A/B 實驗把記憶體屬性單獨拆出來，差距來自它；另外兩件事是一起換的，合起來看不出影響（第 3 章）。中斷對吞吐量的影響另外在同一次開機上量（第 6 章）。

\textbf{靜態檢查。}審稿時重跑發現 \code{twn_fuse_ioctl.h} 有 1 個 WARNING（檔案裡寫了自己的檔名）與 1 個 CHECK（一行超過 100 欄），修掉之後，對著發佈映像那一顆 kernel 的建置樹重跑，兩個檔案都是 0：

\boardfile[：第 3--5 行是發佈 kernel 的選項，第 18--19 行證明 W=1 真的傳到了編譯器]{S6_static_checks.log}{1-6,10-13,16-20,22,24,28,30}

\subsection{沒有做到的兩項，以及怎麼補}

\textbf{\code{make dt_binding_check}。}kernel 官方的 binding 檢查會用 JSON Schema 逐條驗證每個屬性的格式（例如 \code{reg} 的 cell 數），需要 Python 套件 \code{dtschema}（\code{yamllint} 沒裝時只會略過 YAML 的格式檢查）。這次的工作環境不允許從網路安裝新套件，所以沒有跑，只有 \code{binding_check.py} 的部分檢查（第 1 章）。補的方法：

\begin{cmdline}
pip install dtschema yamllint     # yamllint is optional
cp linux/driver/twn_fuse/wjhuang,twn-fuse.yaml <kernel>/Documentation/devicetree/bindings/misc/
make -C <kernel> dt_binding_check DT_SCHEMA_FILES=misc/wjhuang,twn-fuse.yaml
\end{cmdline}

要送進 mainline 還得登記 vendor prefix（第 1 章），這次沒有做。

\textbf{sparse（\code{make C=2}）。}kernel 專用的靜態分析工具，最常抓到的是 \code{__user} 指標被直接取值、\code{__iomem} 指標沒經過 \code{readl()} 就使用（第 5 章的 \code{__user} 與第 2 章的 \code{__iomem} 就是給它看的標記）、大小端型態混用。同樣需要另外安裝，所以沒有跑。補的方法：\code{apt install sparse} 之後在 driver 資料夾執行 \code{make C=2}（\code{C=1} 只檢查這次重新編譯的檔案，driver 已經編過時要先 \code{make clean}）。

\subsection{原理}

\textbf{每一項證明了什麼、沒證明什麼。}

\begin{center}
\small
\begin{tabularx}{\textwidth}{l X X}
\toprule
檢查 & 證明了 & 沒有證明 \\
\midrule
fixture bit-exact & driver 版本對這 4 張輸入算出和 golden 完全相同的結果，兩種 head 擺法都是 & 其他輸入。證明的是這顆 v1.1 bitstream 加上 v1.1 driver 的組合；bitstream 和 v1.0.0 不同（加中斷、改 AxPROT，第 12 章） \\
\errChecks{} 項錯誤路徑 & 每一條檢查規則至少在一種錯誤上有效，而且回報的是預期的錯誤碼 & 所有可能的錯誤組合 \\
\code{rmmod} 被拒、5 次重載 & module 的參考計數有效（開著與只剩映射兩種）；重新 probe 之後硬體仍正常 & 長時間反覆載入（5 次，加上 \code{state} 的 10 次） \\
10 次 64 MB（SMMU 關） & 每次卸載都把 buffer 還給 CMA（第 3 章） & SMMU 開著時從一般記憶體配的頁有沒有歸還（不經過 CMA，640 MB 也小於可用記憶體，量不出來） \\
吞吐量 & driver 版本沒有變慢；中斷比輪詢快約 0.8\%（第 6 章） & 長時間執行的穩定性 \\
\code{W=1}、\code{checkpatch} & 沒有編譯器警告與風格問題 & sparse 的位址空間檢查、binding 的格式檢查 \\
\bottomrule
\end{tabularx}
\end{center}

\textbf{量測時不碰板子。}吞吐量都是在同一條 ssh 指令裡連續跑完的，量測期間沒有其他連線去查板子的狀態，避免查詢本身干擾量測。

\textbf{合成 fixture 與 MPII fixture。}合成的 fixture 為什麼也能當驗證，見第 14 章；同一個映像上 MPII fixture 也全部通過（\code{H4m}），兩組結果一致。

\subsection{面試會問}

\begin{qa}{你怎麼證明 driver 是對的，而且沒有讓系統變慢？}
正確性：4 張 fixture，head 放在 PL 和放在 CPU 兩種擺法，都和 golden 逐位元相同。錯誤路徑：\errChecks{} 項測試，其中 \errRefusals{} 項故意給越界 offset、壞指標、重複開啟、PL 會拒絕的表，都回傳預期的錯誤碼，逾時在 \timeoutSec{} 秒回來，之後下一個 app 照常跑。生命週期：裝置開著、或只剩映射時，\code{rmmod} 都被拒絕；連續 5 次卸載重載都 4/4；SMMU 關掉時連續 10 次各配 64 MB，總量超過 CMA，證明 buffer 有歸還。效能：driver v1.0 剛完成時比 UIO 版本快約 0.8\%，但同時改了三件事分不出原因，後來用四組 A/B 實驗把記憶體屬性單獨拆出來，差距來自它。靜態檢查做了 \code{W=1} 和 \code{checkpatch --strict}；官方的 \code{dt_binding_check} 與 sparse 因為環境不能安裝工具沒有跑，我會直接說出來，並說明怎麼補。
\end{qa}

\subsection{關鍵函式速查}

\begin{center}
\small
\begin{tabularx}{\textwidth}{l l X}
\toprule
名稱 & 位置 & 做什麼 \\
\midrule
\code{board_check.sh} & \code{linux/driver/twn_fuse/test/} & 四類檢查一次跑完，最後印 PASS／FAIL \\
\code{fixtures()} & \code{board_check.sh}:\lineof{check_fixtures} & 兩種 head 擺法各數 4 行 PASS \\
\code{CALL}／\code{check()} & \code{twn_errtest.c}:\lineof{test_call} & 呼叫後立刻讀 \code{errno}，比對預期錯誤碼 \\
\code{v11_check.sh} & \code{linux/driver/twn_fuse/test/} & v1.1 新增機制的板上檢查，一個模式一件事 \\
\code{bind.sh} & \code{linux/driver/twn_fuse/} & 用 \code{driver_override} 在 UIO 與 driver 之間切換，也可帶 module 參數重新載入；UIO 版本的 app 只能在 v1.0.0 的開機檔上跑（要有保留區與 \code{/dev/mem}） \\
\bottomrule
\end{tabularx}
\end{center}

\needspace{20\baselineskip}
\section{進映像與發佈：recipe、預設密碼、合成 fixture}
\label{sec:release}

\begin{oneline}
v1.1 用 PetaLinux（Yocto）的 module recipe 把 \code{twn_fuse.ko} 編進映像、開機自動載入；發佈的 root file system 帶著公開、但建置時就設成過期的預設密碼，另附四張不含人物的合成 fixture 與板上工具包，拿到板子的人不需要 PetaLinux、交叉編譯器或 MPII 資料集就能驗證加速器。
\end{oneline}

\begin{before}[原本怎麼做（v1.0.0 的映像與 driver v1.0）]
v1.0.0 的 \code{system-user.dtsi} 把節點交給 UIO（第 1 章「原本怎麼做」）。driver v1.0 不在映像裡。每次開機要把 \code{twn_fuse.ko} 複製到板上，再用 \code{bind.sh} 透過 \code{driver_override} 把裝置從 UIO 搶過來；driver 是在 PC 上對著另外複製出來的 kernel 編譯樹編的，要靠 vermagic（module 裡記錄的 kernel 版本與主要設定字串，載入時必須和執行中的 kernel 完全相同）一致才能載入。

發佈方面，v1.0.0 沒有附 root file system（rootfs，開機後掛成根目錄的整個檔案系統），因為建好的映像裡有板子的 root 密碼雜湊（\code{/etc/shadow}）；也沒有附 fixture，因為四張 fixture 都是 MPII（一個人體姿態資料集）的人物照片，不能再散布。想在板子上驗證，得自己建 rootfs、自己下載 MPII 重新產生 fixture、自己交叉編譯 app。
\end{before}

\begin{tech}
\begin{itemize}
\item \textbf{Yocto module recipe}（\code{.bb} 檔）：\code{inherit module} 讓這個 recipe 依賴映像自己的 kernel，用同一棵 kernel 樹編譯 out-of-tree module（原始碼不在 kernel 原始碼樹裡、另外編譯的 module）。\code{KERNEL_MODULE_AUTOLOAD} 產生開機載入清單，\code{IMAGE_INSTALL:append} 把套件放進 rootfs。
\item \textbf{用標籤參照改寫節點}：\code{system-user.dtsi} 裡的 \texttt{\&fuse\_0 \{...\}} 覆寫 PetaLinux 產生的節點的 \code{compatible}（第 1 章）。這是建置時的覆寫，不是執行期才套上的 device tree overlay（\code{.dtbo}）。
\item \textbf{公開的預設密碼，一開始就過期}：\code{build_linux.sh} 沒有給 \code{BOARD_PW} 時用 \code{twn-v1.1}，雜湊用 \code{openssl passwd -6}（SHA-512 crypt）算好，經 Yocto 的 \code{EXTRA_USERS_PARAMS} 寫進映像；同一個變數裡的 \code{passwd-expire root} 讓第一次登入一定要先改密碼。
\item \textbf{合成 fixture}：用公式畫出四張圖，標準答案由與加速器逐位元相同的整數參考模型算出。
\item \textbf{板上工具包}（\code{board/make_kit.sh}）：driver 版的 app、測試程式、權重、合成 fixture 打成一個 tar 檔，附 \code{SHA256SUMS}。
\end{itemize}
\end{tech}

\begin{gain}
\begin{itemize}
\item 開機到可以用不需要任何手動步驟：開機第 4.8 秒 driver 就載入並綁上（第 1 章的 \code{H3_boot_state.log}）；module 和 kernel 在同一次建置裡用同一棵樹編出來，vermagic 一定相符。
\item repo 裡只有一份 driver 原始碼：\code{build_linux.sh} 建置時把它複製進 recipe。
\item 開機即用：把開機檔放上 SD 卡、開機、登入並改密碼、解開工具包，\code{sh board_check.sh} 就能得到 \code{BOARD CHECK: PASS}（發佈的映像與工具包在開機組 H 驗過：\code{H0_kit_sha.log}、\code{H4_board_check.log}）。
\item 發佈檔裡沒有個人資料：映像只有公開預設密碼的雜湊；fixture 不含任何人物；執行檔裡沒有建置機器的路徑。
\item 預設密碼不能被留著用：序列埠或 ssh 的第一次登入都要先改密碼；改之前，用 ssh 直接執行指令、用 scp 傳檔都被拒絕。限制：rootfs 是 initramfs，重開機會回到過期的預設密碼。
\end{itemize}
\end{gain}

\subsection{實際例子：把 driver 編進映像}

\textbf{recipe。}整個檔案只有十幾行。\code{LIC_FILES_CHKSUM} 的 md5 是 \code{twn_fuse.c} 第 1 行（SPDX 授權宣告）的 checksum，Yocto 用它確認授權宣告沒有被改；\code{RPROVIDES} 讓套件也能用 Yocto 慣用的 \code{kernel-module-twn-fuse} 這個名字被依賴：

\snip[{}]{recipe_bb}

\textbf{給 module class 用的 Makefile。}\code{KERNEL_SRC} 由 \code{inherit module} 傳進來，指向映像的 kernel 樹：

\snip[make]{recipe_make}

\textbf{放進映像、只保留一份原始碼。}

\snip[{}]{bsp_install}
\snip[bash]{build_copy_src}

開機後的狀態（module 已載入、是映像裡的檔案、在開機載入清單裡、裝置綁在 \code{twn_fuse} 上）見第 1 章的 \code{H3_boot_state.log} 第 4--13 行。

\subsection{實際例子：預設密碼}

密碼本身不進 repo，只把雜湊與過期設定寫進 PetaLinux 專案自己的設定檔：

\needspace{14\baselineskip}
\snip[bash]{build_pw}

\needspace{14\baselineskip}
\textbf{板上：ssh 遇到過期的密碼（開機組 H，在任何人改密碼之前）。}從 PC 用 OpenSSH 試三種情形：直接執行指令、用 scp 傳檔、密碼打錯；再用 \code{expect}（照腳本對互動程式打字的工具）模擬一次互動登入：

\boardfile{H1_ssh_expired.log}{1-13,17-25}

密碼正確時，sshd 先說密碼已過期；沒有終端機（直接執行指令與 scp）就沒辦法改，連線直接結束，指令沒有執行（第 3--10 行）。密碼打錯時是一般的拒絕（第 12 行）。互動登入（第 17--23 行，同一次開機的第二次執行；第 14--16 行的說明與指令沒有印出）先要求改密碼；在 \code{New password:} 按 Ctrl-C 只結束了這一次輸入，接著問確認密碼，連線停在這裡，沒有拿到 shell。

\textbf{用 ssh 完成第一次登入（開機組 H，重開機之後）。}同一個腳本給了新密碼時，最後用互動的 ssh 把密碼改掉，接著用新密碼直接執行指令：

\boardfile[：前面四種情形與上面相同，這裡只印最後兩段]{H9_ssh_first_login.log}{1,24-39}

改完之後這個 ssh 連線 30 秒內沒有再輸出任何東西（第 33--34 行），但第 37--38 行用新密碼直接執行的指令成功了，root 的「上次修改日」也變成當天，表示密碼確實在 ssh 上改掉了。實際使用時，改完若畫面停住，就關掉這個連線（先按 Enter，再輸入 \code{~.}），用新密碼重新登入。

\textbf{板上：序列埠上的第一次登入（開機組 H）。}第 2 行列出實際按的鍵（板子不回顯密碼）。第一次登入在第一個 \code{New password:} 按 Ctrl-C；第二次先給空的新密碼，再給兩個不同的密碼：

\boardfile{H2_first_login.log}{1-9,11-21}

Ctrl-C 讓 \code{login} 整個結束，回到登入提示（第 9 行是重新出現的登入橫幅），沒有拿到 shell。空的新密碼被拒絕（第 18 行），兩次不同也被拒絕（第 21 行），兩種情形都沒有 shell。第三次才正確改掉：

\boardfile{H2_first_login.log}{26-33,37-41}

改完之後 root 的「上次修改日」變成 20752（1970 年起的天數，也就是 2026-10-26）；\code{petalinux} 仍是 0，它的第一次登入同樣要先設密碼。

\begin{pit}
\textbf{候選版 rootfs 的強制改密碼有漏洞。}候選版 rootfs（開機組 E）不讓密碼過期，而是在 \code{/etc/profile.d/} 放一個腳本：互動式 root 登入時先執行 \code{passwd}，改成功前不給 shell（\code{E7_first_login.log}）。審稿時發現它只攔得住「互動的登入 shell」：\code{ssh root@板子 指令}、\code{scp} 都不經過 \code{profile.d}，拿公開的預設密碼就能以 root 執行任何指令。開機組 E 的 E1--E6 本身就是在第一次互動登入（E7）之前、用預設密碼經 ssh 跑的。改成讓密碼過期之後，檢查發生在 PAM 的帳號檢查，所有登入方式都要經過它，腳本就拿掉了。
\end{pit}

\textbf{PetaLinux 自己的 \code{petalinux} 帳號。}映像保留了 PetaLinux 預設建立的這個帳號：密碼是空的、而且已經過期，可以用 \code{sudo} 執行任何指令，登入的 PAM 設定允許空密碼（\code{common-auth} 的 \code{nullok_secure}；這個映像的 \code{pam_unix} 沒有 Debian 那個只限 \code{/etc/securetty} 終端機的修補，等同 \code{nullok}）。由設定推得（沒有在板上測）：能接上序列埠的人可以用空密碼登入 \code{petalinux}、被要求設一個新密碼，再用 \code{sudo} 取得 root，不需要 root 的密碼。sshd 預設 \code{PermitEmptyPasswords no}，所以這只影響序列埠。能接序列埠通常也能直接換 SD 卡，所以這次只在發佈說明與這裡寫明，沒有鎖掉這個帳號；要鎖，可以在 \code{EXTRA_USERS_PARAMS} 加 \code{usermod -L petalinux;}。

\subsection{實際例子：合成 fixture 與工具包}

四張圖完全由公式產生（漸層、棋盤加圓、條紋、平滑雜訊）：

\snip[Python]{syn_images}

腳本若找得到 MPII 的 fixture，會先用同一個參考模型重算第 0 張 MPII fixture，要求結果和既有的標準答案完全相同，確認這個參考模型就是板子拿來比對的那一個，再產生合成的標準答案。四張合成 fixture 在開機組 D、E、H，兩種 head 位置都逐位元相同（\code{D4}、\code{E5}、\code{H5_synthetic.log}）。

\begin{pit}
\textbf{工具包裡的建置路徑。}發佈前掃描發佈檔時，發現工具包裡兩個 app 的執行檔帶著建置機器上的路徑：它們沒有 strip，debug 資訊裡記著編譯器、SDK 與原始碼的絕對路徑。發佈的其他檔案（BOOT.BIN、bitstream、XSA、kernel、兩個 dtb、boot.scr、rootfs 的 2,513 個檔案、權重與模型檔）都沒有這個問題：\code{tools/scan_release_paths.py} 對發佈檔逐一掃描，找到 0 處，同一個工具對 strip 之前的 app 找到 8 處（\code{S7_release_path_scan.log}）。修法是 \code{board/make_kit.sh} 用 SDK 的 \code{strip} 去掉符號與 debug 資訊，strip 之後、打包之前再掃一次工具包的每個檔案，含有建置路徑就中止。去掉 debug 資訊不改變程式碼，但檔案的雜湊變了，所以工具包重新在板上驗過（開機組 E 的 \code{E9_release_kit.log}，發佈的版本在開機組 H）。
\end{pit}

\code{make_kit.sh} 的 strip 與打包前的掃描：

\snip[bash]{kit_strip}

\subsection{原理}

\textbf{開機時 module 是怎麼被載入的。}有兩條路，這個映像兩條都有：
\begin{itemize}
\item \textbf{依 modalias}：kernel 建立平台裝置時就有一個 uevent，帶著 \code{MODALIAS}（內容是 \code{of:N...T...C} 加上節點的 \code{compatible}）。那時開機還很早，udev 還沒啟動，沒有人收；udev 啟動後把既有裝置的事件重送一次（coldplug），才用內建的 libkmod 載入它（效果等同執行 \code{modprobe}），也就是在 \code{modules.alias} 找到這個字串屬於 \code{twn_fuse}（來自第 1 章的 \code{MODULE_DEVICE_TABLE(of, ...)}）。
\item \textbf{依清單}：systemd（這個映像的 init 系統）的 \code{systemd-modules-load} 讀 \code{/etc/modules-load.d/*.conf}，逐一載入。
\end{itemize}
不論哪一條先載入，第二條都只會發現它已經載入。載入後，driver core 用 \code{compatible} 比對，呼叫 \code{probe()}。\code{driver_override} 是在 \code{compatible} 對不上時強迫綁定的手段，節點已經是 driver 認得的字串，就不再需要；\code{bind.sh} 仍然保留：它用 \code{driver_override} 在 UIO 與 driver 之間切換（第 13 章的卸載重載測試就用它），也可以帶 module 參數重新載入。

\textbf{\code{inherit module} 做了什麼。}讓 recipe 依賴 \code{virtual/kernel}，編譯時把 \code{KERNEL_SRC} 指向這次建置的 kernel 樹，安裝時把 \code{.ko} 放進 \code{/lib/modules/<版本>/extra/}，並在映像建立時跑 \code{depmod}：它掃描所有 module，產生 \code{modules.alias} 這張「modalias 字串 → module 名稱」的對照表。

\textbf{過期的密碼怎麼擋住所有登入方式。}\code{/etc/shadow} 每一列的第 3 欄是上次修改密碼的日期（1970 年起的天數），0 的意思是「必須在下次登入時修改」。序列埠的 \code{login} 與 sshd 都用 PAM 檢查帳號，\code{pam_unix} 看到 0 就回報「需要新密碼」：\code{login} 接著要求改密碼，不改就結束；sshd 在有終端機時同樣要求改，沒有終端機（直接執行指令、scp、sftp）就印出 \code{Password change required but no TTY available} 並結束連線。檢查發生在登入的那一刻，和之後執行的是哪一種 shell 無關，這是它比 \code{profile.d} 腳本可靠的原因。為什麼不放真正的密碼：雜湊可以離線暴力破解，而且同一份映像會被很多人使用，等於共用一組密碼；公開預設密碼加上第一次登入強制修改，安全性來自「第一次登入的人一定要改掉」，所以第一次開機要在自己控制的網路上。

\textbf{initramfs 的副作用。}這個映像的 root file system 是開機時載入記憶體的 initramfs，改的密碼只存在記憶體裡，重開機就回到過期的預設密碼。要保存，得把 root file system 放到 SD 卡的分割區上；這次沒有做。

\textbf{為什麼合成圖也能當驗證。}驗證的目標是「PL 的計算和參考模型逐位元相同」，不是「姿態估得準不準」。任何輸入都會走過全部 65 個 operator，只要每個 operator 的輸出逐位元相同，用什麼圖都可以。四張圖分別有平滑區、尖銳邊緣、週期圖樣和雜訊，讓數值範圍盡量分散。

\subsection{面試會問}

\begin{qa}{driver 怎麼做到開機就自動載入？}
我寫了一個 Yocto 的 module recipe：\code{inherit module} 讓它用映像自己的 kernel 樹編譯，所以 vermagic 一定對；\code{KERNEL_MODULE_AUTOLOAD} 產生 \code{/etc/modules-load.d} 的設定，開機由 systemd 載入；\code{IMAGE_INSTALL} 把它放進 rootfs。device tree 在 \code{system-user.dtsi} 把 PetaLinux 產生的節點的 \code{compatible} 改成 driver 比對表裡的字串，所以載入後 driver core 直接比對成功、呼叫 probe，不再需要 \code{driver_override}。另外 udev 也會用 modalias 載入它，兩條路都在。板上開機第 4.8 秒 driver 就綁好，用中斷模式。
\end{qa}

\begin{qa}{你的發佈映像怎麼處理 root 密碼？}
映像裡不放任何真正的密碼。建置腳本沒有給密碼時用公開的預設密碼，雜湊用 \code{openssl passwd -6} 算好、經 \code{EXTRA_USERS_PARAMS} 寫進去，同一行用 \code{passwd-expire} 讓它一開始就過期。序列埠和 ssh 的第一次登入都會先要求改密碼；改掉之前，ssh 帶指令與 scp 都回「Password change required but no TTY available」；序列埠上按 Ctrl-C 回到登入提示，空密碼和兩次不一致都被拒絕。候選版我用 \code{profile.d} 腳本，審稿時發現它擋不住非互動的 ssh，才改成讓密碼過期。限制是 rootfs 是 initramfs，重開機會回到過期的預設密碼。
\end{qa}

\subsection{關鍵函式速查}

\begin{center}
\small
\begin{tabularx}{\textwidth}{l X p{36mm}}
\toprule
名稱 & 位置（\code{meta-user/} 指 \code{linux/project-spec/meta-user/}） & 做什麼 \\
\midrule
\code{twn-fuse.bb} & \code{meta-user/recipes-modules/twn-fuse/twn-fuse.bb}:\lineof{recipe_bb} & module recipe、開機載入 \\
\code{IMAGE_INSTALL} & \code{meta-user/conf/petalinuxbsp.conf}:\lineof{bsp_install_line} & 放進 rootfs \\
複製原始碼 & \code{linux/build_linux.sh}:\lineof{build_copy_for} & recipe 用 repo 裡唯一的一份 \\
預設密碼與過期 & \code{linux/build_linux.sh}:\lineof{build_pw_line} & 沒給 \code{BOARD_PW} 時寫入雜湊與 \code{passwd-expire} \\
\code{images()} & \code{tools/make_synthetic_fixtures.py}:\lineof{syn_images} & 四張合成圖 \\
\code{make_kit.sh} & \code{board/make_kit.sh}:\lineof{kit_strip} & strip、掃描建置路徑、打包工具包 \\
\bottomrule
\end{tabularx}
\end{center}

\appendix
\clearpage
\section{名詞表}

依英文字母排序（不分大小寫）。「章」是第一次詳細說明的地方，「導」是導讀，「C」是附錄 C。

\footnotesize
\begin{xltabular}{\textwidth}{l X c}
\toprule
名詞 & 一句話說明 & 章 \\
\midrule
\endhead
A/B 實驗 & 兩組設定只差一件事，用來確定是哪一件事造成差異 & 3 \\
ABI & 給使用者程式的介面約定；改了就會讓舊程式壞掉 & 8 \\
\code{access_ok} & 檢查一段位址是否完全落在 user space 範圍內 & 5 \\
arena、buffer & PL 讀寫的那塊 DDR；v1.0.0 是 256 MB 保留區（arena），driver 版本是 8 MB DMA buffer & 導 \\
AXI、beat & ARM 的晶片內部匯流排／它一次傳輸的資料量，這個設計是 16 bytes & 導 \\
AxPROT & AXI 交易附帶的 3 個保護位元；bit 1 是 0 表示 Secure、1 表示 Non-secure & 11 \\
barrier & 限制 CPU 記憶體存取先後順序的指令 & 2 \\
binding & device tree 節點的規格檔，寫明可以有哪些屬性 & 1 \\
bit-exact & 逐位元相同：輸出和標準答案的每個 byte 都一樣 & 導 \\
block RAM tile & FPGA 晶片上固定大小的記憶體區塊，計算 block RAM 用量的單位 & 12 \\
\code{boot.scr}、bootargs & U-Boot 開機時執行的指令稿／U-Boot 傳給 kernel 的命令列參數 & C \\
bypass & SMMU 對沒有登記的 stream ID 不做轉換、直接放行；發佈的 kernel 是這樣設定 & 11 \\
CCI & Cache Coherent Interconnect，會去查 CPU cache 的互連，HPC 埠經過它 & 3 \\
checkpatch & kernel 的程式風格檢查腳本 \code{scripts/checkpatch.pl} & 13 \\
CLB & configurable logic block，FPGA 裡放 LUT 與暫存器的基本區塊 & 12 \\
CMA & Contiguous Memory Allocator，開機劃出、平時借給可搬移用途的連續記憶體區 & 3 \\
coherent（一致） & CPU 和裝置看到的記憶體內容保證相同 & 3 \\
\code{compatible} & device tree 節點上的硬體名字，格式「製造商,型號」 & 1 \\
completion & kernel 的同步機制：一方等待，另一方（常是 ISR）呼叫 \code{complete()} 叫醒它 & 6 \\
\code{container_of} & 由 struct 中某個欄位的位址，算出整個 struct 的位址 & 7 \\
context fault & 交易在 SMMU 轉換時找不到對應的頁，被擋下並回報 & 11 \\
\code{copy_from_user} & 安全地從 user 記憶體複製到 kernel，失敗時回傳未複製的 byte 數 & 5 \\
CSR & Control/Status Registers，加速器的控制暫存器 & 導 \\
debugfs & 只給開發者用的 kernel 檔案系統，掛在 \code{/sys/kernel/debug} & 8 \\
depmod & 掃描所有 module、產生 \code{modules.alias} 等對照表的工具 & 14 \\
descriptor table & 描述子表：17 個 block 的設定，經 driver 寫進 CSR，PL 依表執行 & 導 \\
device tree & 描述硬體有哪些周邊、位址在哪的資料結構，開機時交給 kernel & 1 \\
Device 記憶體 & 給周邊用的記憶體屬性：不走 cache、不預讀；nGnRnE 表示不合併、不重排、寫入要等終點回應 & 2 \\
\code{devm_*} & device-managed 資源，driver 離開或 probe 失敗時自動釋放 & 1 \\
directive & 同一個 Vivado 指令的不同最佳化模式 & 12 \\
\code{dma_addr_t} & 裝置看到的位址；沒有 IOMMU 時等於實體位址，有 SMMU 時是 IOVA & 3 \\
DMA、DMA mask & 裝置不經 CPU 直接讀寫記憶體／裝置能定址的位元數，這裡是 40 & 3 \\
\code{dmb}／\code{dsb} & 兩種 barrier：\code{dmb} 只排順序，\code{dsb} 要等存取真的完成 & 2 \\
\code{driver_override} & sysfs 檔案，強制指定裝置只能交給哪個 driver & 1 \\
DSP & FPGA 晶片上的硬體乘加單元 & 12 \\
dt-schema & 用 JSON Schema 寫 device tree binding 的格式與檢查工具 & 1 \\
ECO & engineering change order，在已經 route 好的設計上只改一小部分 & 12 \\
\code{-EPROBE_DEFER} & probe 的回傳值，表示依賴的東西還沒準備好，kernel 稍後會再呼叫一次 probe & 6 \\
errno & 系統呼叫失敗時的錯誤碼，例如 \code{EINVAL}、\code{EFAULT} & 5 \\
exception table & kernel 記錄「這些指令出錯時跳到哪裡修復」的表 & 5 \\
fixture & 測試輸入加上它的標準答案（golden） & 導 \\
GIC、SPI & ARM 的中斷控制器／所有 CPU 共用的外部中斷，GIC 編號是 SPI 編號加 32 & 6 \\
HAL & Hardware Abstraction Layer，上層只依賴介面，底下可換實作 & 9 \\
HP／HPC 埠 & PL 存取 DDR 的 AXI 埠；HP 不經 CPU cache，HPC 可以經 CCI 保持一致 & 3 \\
hrtimer & kernel 的高精度計時器，\code{usleep_range} 靠它叫醒 & 6 \\
initramfs & 開機時整個載入記憶體的 root file system，重開機就回到原樣 & 14 \\
interrupt context & CPU 被中斷打斷、還沒回到任何 process 時執行的程式，不能睡眠 & 6 \\
interrupt storm & level 中斷的來源沒清掉，ISR 被一直重複呼叫 & 6 \\
ioctl & 對裝置檔下指令的系統呼叫 & 5 \\
IOMMU group／domain & group 是必須共用同一套轉換的一組裝置；domain 是一套轉換表 & 11 \\
IOMMU、SMMU & 替裝置做位址轉換與存取檢查的硬體；ARM 上叫 SMMU，ZynqMP 用 MMU-500 & 11 \\
\code{ioremap} & 在 kernel 頁表建立對周邊實體位址的映射 & 2 \\
\code{IO_STRICT_DEVMEM} & kernel 選項：\code{/dev/mem} 也不能存取被 driver 佔用中的 I/O 區間 & 10 \\
IOVA & I/O 虛擬位址：經過 SMMU 時，裝置看到的位址 & 11 \\
ISR & Interrupt Service Routine，中斷發生時 kernel 呼叫的 driver 函式 & 6 \\
jiffy & kernel 的時間刻度，\code{HZ=250} 時是 4 ms & 6 \\
level／edge & 中斷的兩種觸發方式：線保持高電位的期間都算／只在變化的瞬間算一次 & 6 \\
linear map（線性映射） & kernel 頁表裡把整個 DDR 一對一對應進來的那一段 & 3 \\
LUT & look-up table，FPGA 裡實作組合邏輯的基本單元 & 12 \\
mainline & Linus Torvalds 維護的官方 kernel 原始碼樹 & 1 \\
misc device & 主裝置號 10、次裝置號自動分配的字元裝置 & 5 \\
\code{mmap} & 把裝置或檔案的內容映射進 process 的位址空間 & 4 \\
MMIO & Memory-Mapped I/O，用讀寫記憶體位址的方式操作暫存器 & 2 \\
modalias & 裝置的識別字串，udev 用它找出該載入哪個 module & 14 \\
module parameter & 載入 module 時可以給的參數，例如 \code{use_irq=0} & 1 \\
MPII & 一個人體姿態資料集；v1.0.0 的 fixture 是其中的照片，含人物，不能再散布 & 14 \\
mutex、spinlock & 拿不到就睡眠等待的鎖／拿不到就空轉等待的鎖，拿著時不能睡 & 7 \\
NBLK & 「這一幀跑幾個 block」的暫存器，平常是 17，reset 值也是 17 & 導 \\
netlist & 合成產生的電路元件與連線清單 & 12 \\
Normal non-cacheable & 一般記憶體屬性但不走 cache；寫入可以合併，可以非對齊存取 & 3 \\
oops & kernel 程式出錯時印出錯誤現場並終止目前工作的機制 & 2 \\
outer shareable & barrier 的範圍之一：CPU 與其他會存取記憶體的裝置 & 2 \\
out-of-tree module & 原始碼不在 kernel 原始碼樹裡、另外編譯的 module & 14 \\
page cache & kernel 快取檔案內容的記憶體 & 3 \\
PAM & Linux 的登入驗證框架；\code{login} 與 sshd 都經過它檢查密碼與帳號狀態 & 14 \\
PAN & Privileged Access Never，ARMv8.1 起的功能，禁止 kernel 直接存取 user 記憶體；A53 沒有 & 5 \\
\code{passwd-expire} & Yocto \code{EXTRA_USERS_PARAMS} 的指令，把密碼設成已過期，下次登入必須先改 & 14 \\
pblock & Vivado 的位置限制：把某些邏輯放在晶片的某一塊區域 & 12 \\
PFN & page frame number，實體位址除以頁大小（4096） & 10 \\
phys\_opt\_design & Vivado 在 placement 或 route 之後，針對 timing 做的局部最佳化 & 12 \\
pipeline & 量吞吐量用的 app，讓 backbone 與 head 在不同幀上重疊 & 導 \\
placement／route & 決定每個邏輯單元放在晶片的哪個位置／決定它們之間怎麼繞線 & 12 \\
platform device／driver & 不能自動列舉的 SoC 周邊，和它們的 driver & 1 \\
PMU & Platform Management Unit，ZynqMP 的平台管理處理器；DDR 最上面 1 MB 保留給它 & 10 \\
posted write & 寫入送出後 CPU 就繼續執行，不等寫入真正到達 & 6 \\
\code{probe()}／\code{remove()} & driver 接管、放開裝置時被呼叫的函式 & 1 \\
PS／PL & Zynq 的 CPU 部分／FPGA 邏輯部分 & 導 \\
\code{readl_poll_timeout} & 輪詢暫存器直到條件成立或逾時的巨集 & 6 \\
recipe & Yocto 的建置描述檔（\code{.bb}），說明一個套件怎麼取得、編譯、安裝 & 14 \\
retiming & 把暫存器搬過組合邏輯，讓前後兩段路徑的延遲比較平均 & 12 \\
ring overrun & PL 的列緩衝區在滿的時候被寫入；STATUS bit 16 記錄它 & 8 \\
rootfs & root file system，開機後掛成根目錄的整個檔案系統 & 14 \\
Secure／Non-secure & ARM TrustZone 的兩個世界；Linux 跑在 Non-secure & 11 \\
slack、WNS、WHS、TNS & 路徑的時間餘裕；最差的 setup 餘裕／最差的 hold 餘裕／所有違規路徑的負餘裕加總 & 12 \\
SoC & System on Chip，把 CPU 與周邊做在同一顆晶片上 & 1 \\
sparse & kernel 專用的靜態分析工具，用 \code{make C=2}（或 \code{C=1}，只檢查重新編譯的檔案）執行 & 13 \\
SPDX & 用標準化的授權名稱寫在原始檔第一行的授權宣告 & 14 \\
strategy & Vivado 預先組好的一整組 place \& route 選項 & 12 \\
stream ID & 每筆交易進 SMMU 時帶的來源編號，SMMU 用它決定用哪一組轉換 & 11 \\
streaming DMA & 用 cacheable 記憶體，每次傳輸前後手動做 cache 維護 & 3 \\
\code{STRICT_DEVMEM} & kernel 選項：\code{/dev/mem} 不能存取 System RAM & 10 \\
sysfs & 把 kernel 物件攤成檔案樹的檔案系統，掛在 \code{/sys} & 8 \\
System RAM & kernel 管理的記憶體範圍（\code{/proc/iomem} 裡的同名項目）：device tree 的 \code{/memory} 節點，扣掉標成 \code{no-map} 的保留區 & 10 \\
systemd & 這個映像的 init 系統；\code{systemd-modules-load} 是其中載入 module 的服務 & 14 \\
TBU & SMMU 的轉換單元，ZynqMP 有 6 個，各自負責一組匯流排埠 & 11 \\
threaded IRQ & 把中斷處理函式放到一個 kernel 執行緒裡、可以睡眠的處理方式 & 6 \\
TLB & translation lookaside buffer，位址轉換結果的快取 & 11 \\
TrustZone & ARM 把系統分成 Secure 與 Non-secure 兩個世界的硬體機制 & 11 \\
uapi header & kernel 和 user space 共用的介面定義檔 & 5 \\
U-Boot & 開機載入程式：從 SD 卡讀 kernel、device tree 與 rootfs，再啟動 kernel & C \\
uevent／udev & kernel 通知 user space「有裝置出現」的訊息／收到後負責載入 module 的程式；udev 啟動後重送既有裝置的事件叫 coldplug & 14 \\
UIO & Userspace I/O，把裝置的暫存器開給 user space 的通用 driver & 1 \\
vendor prefix & \code{compatible} 字串逗號前面那一段，代表廠商或作者 & 1 \\
vermagic & module 裡記錄的 kernel 版本與主要設定字串，載入時必須和執行中的 kernel 相同 & 14 \\
VMA & \code{vm_area_struct}，描述 process 位址空間中一段連續範圍 & 4 \\
W=1 & kernel 編譯時打開額外警告的選項 & 13 \\
write gathering & 把連續的小寫入在 CPU 內合併成一筆大寫入 & 3 \\
XSA & Vivado 匯出的硬體描述檔，PetaLinux 用它產生 device tree & 6 \\
Yocto & PetaLinux 底下用來建置整個 Linux 映像的系統 & 14 \\
\bottomrule
\end{xltabular}
\normalsize

\needspace{20\baselineskip}
\section{面試題：30 秒版本}

每題先講結論，再講一個板上的數字或證據。完整版在括號中的章節。

\begin{qashort}{1. compatible 字串怎麼讓 kernel 找到 driver、呼叫 probe？（第 1 章）}
開機時 device tree 的每個 PL 節點被建成 platform device。driver 的 \code{of_match_table} 放了 \code{"wjhuang,twn-fuse-1.1"} 與 \code{-1.0}，註冊時 platform bus 用 \code{platform_match()} 比對：有 \code{driver_override} 就只比 driver 名字，否則比 compatible。對上了就呼叫 \code{probe()}：映射暫存器、配 DMA buffer、申請中斷、建 \code{/dev/twn_fuse}，全部用 \code{devm_}。版本號讓舊 driver 不會去等新硬體才有的中斷。
\end{qashort}

\begin{qashort}{2. 為什麼要 ioremap？readl／writel 和 barrier 是什麼？（第 2 章）}
開了 MMU 之後指標都是虛擬位址，kernel 頁表裡沒有周邊，直接存取會 page fault。\code{ioremap} 建一筆 Device 屬性的頁表：不走 cache、不合併、不重排。存取用 \code{readl}／\code{writel}：\code{writel} 寫入前有 \code{dma_wmb}，app 寫進 DDR 的輸入一定比 START 先讓 PL 看到；\code{readl} 之後有 \code{dma_rmb}，讀到 done 之後讀輸出不會被提前。
\end{qashort}

\begin{qashort}{3. DMA buffer 怎麼配？為什麼 driver 版本比較快？（第 3 章）}
\code{dma_alloc_coherent} 同時給 CPU 位址和 DMA 位址；不用 \code{kmalloc}（cacheable、最大 4 MB、不給 DMA 位址），也不用 \code{/dev/mem}。HP 埠不查 CPU cache，kernel 把 buffer 映射成 non-cacheable。快的原因用四組 A/B 實驗把記憶體屬性單獨拆出來：「CPU 端屬性從 Device 改成 Normal non-cacheable」這一步有差，\abDrvDev{} 到 \abDrvNc{} fps，因為寫入可以合併，寫輸入從 \splitDevIn{} 降到 \splitNcIn{} ms（每次執行的最後一幀）。
\end{qashort}

\begin{qashort}{4. app 為什麼只看得到 buffer？為什麼用 mmap 不用 read／write？（第 4 章）}
app 的頁表由 kernel 決定。driver 的 \code{.mmap} 只用 \code{dma_mmap_coherent} 映射那塊 buffer，檢查長度與 offset；暫存器只在 kernel 裡 ioremap。執行中的 app 的 \code{maps} 只有 8 MB 的 \code{/dev/twn_fuse}。每幀幾百 KB 的資料用 \code{read}／\code{write} 要多複製兩次，mmap 讓 app 和 PL 看同一塊記憶體。
\end{qashort}

\begin{qashort}{5. ioctl 和其他介面怎麼選？（第 5 章）}
依資料性質選：大量資料用 mmap 零複製；對裝置的動作和結構化參數用 ioctl，\code{_IOWR} 把方向和大小編進指令碼；給人看的狀態用 sysfs，除錯用 debugfs。所以 file operations 只有 open、release、mmap、unlocked\_ioctl。一次 ioctl \ioctlUs{} µs，跑一幀的系統呼叫、複製加讀回 counter 合計 \bbSysUs{} µs，介面不是瓶頸。
\end{qashort}

\begin{qashort}{6. 為什麼用 copy\_from\_user？有了 SMMU 為什麼還只收 offset？（第 5、11 章）}
user 指標可能沒有映射，或故意指向 kernel；\code{copy_from_user} 先用 \code{access_ok} 檢查範圍，出錯時由 exception table 接住，我回 \code{-EFAULT}；A53 沒有 PAN，硬體不會擋。只收 offset、檢查對齊與範圍（含溢位）再加上 DMA 位址，是第一道防線；SMMU 只擋得住帶著登記過的 stream ID 的 Non-secure 交易，也不管 buffer 內部寫到哪一區。\errChecks{} 項錯誤測試都符合預期。
\end{qashort}

\begin{qashort}{7. 怎麼把輪詢改成中斷？（第 6 章）}
RTL 加一個暫存器輸出的 level 中斷，等於致能 AND done，致能放在獨立暫存器，免得被 START 寫掉；IP 宣告成 interrupt 介面，接 \code{pl_ps_irq0}，device tree 是 SPI 89。driver 用 \code{devm_request_irq}，ISR 清 done、\code{complete()}，返回前讀回一次確保清除到達；\code{reinit_completion} 放在 START 之前。1000 幀剛好 1000 次中斷；只切換 \code{use_irq} 比較，每幀 CPU 時間從 \pollCpuUs{} 降到 \irqCpuUs{} µs；另一次開機交替量三輪，吞吐量從 \hPollFps{} 變成 \hIrqFps{} fps。
\end{qashort}

\begin{qashort}{8. 怎麼處理並行？為什麼用 mutex 不用 spinlock？（第 7 章）}
兩層：\code{open} 用 \code{test_and_set_bit} 讓同一時間只有一次 open 成功，因為描述子表與 base 位址是整顆加速器共用的狀態；backbone 和 head 各一把 mutex，同一個引擎的暫存器操作不被插隊。拿著鎖時會睡在 completion 上好幾 ms，spinlock 不能睡。也拿掉了 sysfs 的 unbind，因為它會在 buffer 還被映射時釋放 buffer。
\end{qashort}

\begin{qashort}{9. 狀態和除錯功能怎麼給？（第 8 章）}
狀態放 sysfs：六個 cycle counter、STATUS、NBLK 各一個唯讀檔，交給 \code{.dev_groups}，driver core 在 probe 後建檔、remove 前移除；跑完一幀 \code{cycles_total} 約 99 萬 cycles。除錯放 debugfs，不當 ABI：\code{debugfs_create_regset32} 印暫存器，可寫的 \code{nblk} 檢查範圍、拿 backbone 的鎖，寫 99 回 \code{EINVAL}。
\end{qashort}

\begin{qashort}{10. 為什麼限制 /dev/mem？（第 10 章）}
它把整個實體位址開給 root。driver 用 DMA API 配 buffer、只 mmap buffer，app 不再需要 \code{/dev/mem}，256 MB 保留區也拿掉了。\code{STRICT_DEVMEM} 擋 System RAM，\code{IO_STRICT_DEVMEM} 擋 driver 佔用中的暫存器；板上 \code{devmem} 都回 \code{Operation not permitted}。邊界也講得出來：不是 System RAM、也沒有 driver 佔用的位址都還能映射，例如卸載 driver 後的暫存器、DDR 最上面 1 MB 的 PMU 保留區。
\end{qashort}

\begin{qashort}{11. SMMU 怎麼讓 PL 只能碰自己的 buffer？（第 11 章）}
device tree 打開 SMMU，給四個 HP 埠登記 stream ID；DMA API 自動把 buffer 映射到 IOVA，越界讀取變成 context fault。第一次失敗是因為 AXI master 送的是 Secure 交易，Linux 只設定 Non-secure 那一側；改成 AxPROT[1] = 1 才生效。另外三個埠的 ID 我推算錯了兩個，用 \code{disable_bypass} 讓 SMMU 回報才量出來。限制：發佈的 kernel 對沒登記的 stream ID 是 bypass。
\end{qashort}

\begin{qashort}{12. 改了 RTL 之後 timing 收不了怎麼辦？（第 12 章）}
先固定配方，只動 RTL。建置腳本在 route 後若 WNS 小於 0，自動跑 \code{phys_opt_design}，不行再重新 route。把 AxPROT[1] 從 0 改成 1，只改常數，route 完卻是 $-$0.032 ns，靠 route 後的 \code{phys_opt_design}（strategy 內建一次、建置腳本再兩次）補到 0。備案是確認 Vivado 對同樣的輸入給出同樣的結果，重現上一版收斂的設計，再用 ECO 只改那幾根常數 pin；最後沒有用到。
\end{qashort}

\begin{qashort}{13. 怎麼證明 driver 是對的？（第 13 章）}
fixture 在兩種 head 擺法都和 golden 逐位元相同；\errChecks{} 項錯誤路徑回傳預期的錯誤碼，逾時 \timeoutSec{} 秒回來；裝置開著或只剩映射時 \code{rmmod} 被拒，5 次重載都 4/4；SMMU 關掉時 10 次各配 64 MB、總量超過 CMA，證明 buffer 有歸還。\code{W=1} 與 \code{checkpatch --strict} 0 個警告；\code{dt_binding_check} 與 sparse 沒跑，我會說明原因和補法。
\end{qashort}

\begin{qashort}{14. driver 怎麼進映像？發佈映像怎麼處理密碼？（第 14 章）}
Yocto 的 module recipe：\code{inherit module} 用映像自己的 kernel 樹編，vermagic 一定對；\code{KERNEL_MODULE_AUTOLOAD} 讓 systemd 開機載入，udev 也會依 modalias 載入；device tree 的 \code{compatible} 改成 driver 的字串。密碼用公開的預設值，建置時就用 \code{passwd-expire} 讓它過期，所以第一次登入一定要改；改之前 ssh 帶指令與 scp 都被拒絕。fixture 換成公式產生的合成圖，發佈檔沒有個人資料。
\end{qashort}

\clearpage
\section{如何重現}

以下指令都在 repo 根目錄執行。主機端在 WSL（PetaLinux 2023.2 與它的 SDK），板子端用 Git Bash 的 \code{plink} / \code{pscp} 連線。

\textbf{1. 編譯 driver、測試程式與兩版 app（WSL）}

\begin{cmdline}
cd linux/driver/twn_fuse
make prepare        # once: copy the PetaLinux kernel build dir and run modules_prepare
make                # twn_fuse.ko
make test           # test/twn_errtest
cd ../../..
bash board/build_apps.sh            # build/arm      (HAL=uio)
HAL=drv bash board/build_apps.sh    # build/arm-drv  (HAL=drv)
modinfo linux/driver/twn_fuse/twn_fuse.ko | grep vermagic   # must equal `uname -r` on the board
\end{cmdline}

\textbf{2. 放到板子上（同一個目錄）}

兩版 app 編出來的檔名相同，複製到板子時要改名，加上 \code{_uio} 與 \code{_drv}，\code{board_check.sh} 找的是改名後的檔案：

\begin{cmdline}
linux/driver/twn_fuse/{twn_fuse.ko, bind.sh, test/twn_errtest, test/board_check.sh}
build/arm/mbv2_fuse_test      -> mbv2_fuse_test_uio
build/arm/mbv2_fuse_pipe      -> mbv2_fuse_pipe_uio
build/arm-drv/mbv2_fuse_test  -> mbv2_fuse_test_drv
build/arm-drv/mbv2_fuse_pipe  -> mbv2_fuse_pipe_drv
data/fuse_weights.bin  data/model_mbv2.bin  data/model_mbv2_packed.bin
golden/fix_0.bin ... golden/fix_3.bin      # regenerate from MPII, docs/QUICKSTART.md
\end{cmdline}

\textbf{3. 板上檢查}

\begin{cmdline}
sh board_check.sh            # prints BOARD CHECK: PASS only if every check passed
sh bind.sh drv ./twn_fuse.ko # or switch by hand
sh bind.sh uio
\end{cmdline}

\textbf{4. 只靠 compatible 綁定（可回退）}

\begin{cmdline}
sh linux/driver/twn_fuse/patch_dtb.sh system.dtb system_twn_fuse.dtb   # host, needs fdtput
# on the card: keep system.dtb as system_uio.dtb, copy system_twn_fuse.dtb over system.dtb, reboot
insmod twn_fuse.ko           # binds by compatible, no driver_override
# to go back: copy system_uio.dtb over system.dtb, reboot
\end{cmdline}

\textbf{5. 重建這份講義}

\begin{cmdline}
cd docs/zh-TW/driver-handout
latexmk                      # runs extract.py first; output build/handout.pdf
\end{cmdline}

這次的板上紀錄放在 \code{results/raw/linux_driver_20261026/}（各檔的用途見該資料夾的 \code{README.md}），數字整理在 \code{results/board_results.json} 的 \code{linux_driver_twn_fuse_20261026} 這一筆。

\end{document}
